\documentclass[11pt]{article}
\usepackage[utf8]{inputenc}
\usepackage{amsmath,amssymb,amsthm}
\usepackage{booktabs}
\usepackage{graphicx}
\usepackage[margin=1in]{geometry}
\usepackage{natbib}
\usepackage{xcolor}
\usepackage[colorlinks=true,linkcolor=blue!60!black,citecolor=blue!60!black,urlcolor=blue!60!black]{hyperref}
\usepackage[expansion=false]{microtype}

\newtheorem{theorem}{Theorem}
\newtheorem{proposition}[theorem]{Proposition}
\newtheorem{lemma}[theorem]{Lemma}
\newtheorem{corollary}[theorem]{Corollary}
\theoremstyle{remark}
\newtheorem{remark}[theorem]{Remark}
\newtheorem{hypothesis}[theorem]{Hypothesis}

\newcommand{\etacrit}{\eta_{\mathrm{crit}}}
\newcommand{\hetacrit}{\widehat{\eta}_{\mathrm{crit}}}
\newcommand{\Dhat}{\widehat{D}}
\newcommand{\Yhat}{\widehat{Y}}
\newcommand{\E}{\mathbb{E}}
\newcommand{\R}{\mathbb{R}}
\newcommand{\Var}{\operatorname{Var}}
\newcommand{\Cov}{\operatorname{Cov}}
\newcommand{\tr}{\operatorname{tr}}
\newcommand{\todo}[1]{\textcolor{red}{[TODO: #1]}}

\title{The Exact Mean-Square Stability Boundary of Heavy-Ball SGD:\\a Lyapunov Equation and a Secular Equation}
\author{Gabriel Paiva\\ Independent Researcher\\ \texttt{gabriel@echdesign.co.uk}}
\date{Draft of 3 November 2026\thanks{This draft is the second-moment part of a longer manuscript in preparation, on identifying the edge of stability by extrapolating the step variance, as \citet{southwell1932} located the buckling load of a column from sub-critical deflections. The estimator, its minimax analysis and its measurements on networks are in that manuscript and are not repeated here; the network checkpoints of Section~\ref{sec:phase2} come from it.}}

\makeatletter\g@addto@macro\UrlBreaks{\do\-}\makeatother
\begin{document}
\maketitle

\begin{abstract}
We study the linear mean-square stability of heavy-ball SGD on a quadratic loss, $w_{t+1}=((1+\beta)I-\eta H_{B_t})w_t-\beta w_{t-1}-\eta\xi_t$, with minibatch Hessians that need not commute with each other or with their mean $H$. We derive the exact stationary second-moment equation (Theorem~\ref{thm:F}), relate it exactly to the small-step Schur-complement expansion of \citet{andreyev2026momentum}, and show that beyond the threshold, $\rho(\mathcal T)\ge1$ for the second-moment operator, the second moment grows at least linearly (Corollary~\ref{cor:msinstab}), the mean-square part of the instability criterion they list as open for SGD with momentum. For random-feature minibatches with independent coordinates of any kurtosis the boundary is the root of a rank-one secular equation (Theorem~\ref{thm:G}); at $\beta=0$, $B=1$ and Gaussian features this is the mean-square condition of the LMS filter. For arbitrary per-sample Hessians it is the reciprocal of the top eigenvalue of a symmetrisable operator (Theorem~\ref{thm:G6}), the criterion that \citet[Theorem~F.1]{cai2026llm} obtained concurrently and independently. At the threshold an identity places a modal curvature between the two batch-sharpness plateaus $2(1-\beta)/\eta$ and $2(1+\beta)/\eta$ reported by \citet{andreyev2026momentum}, in the random-feature model a stationary batch sharpness reaches them in the two batch limits, and a variational test gives sufficient instability conditions through batch sharpness (Theorem~\ref{thm:modal-bridge}); in a quasi-stationary slow-fast model progressive sharpening saturates just below the threshold (Theorem~\ref{thm:slow-fast-saturation}).

The secular root matches sample paths with median error $0.19$--$3.82\,\%$ in 29 of 30 cells and the spectral radius of the exact operator within $6.11\times10^{-10}$ in 36 cells. On tethered Hessians of networks with $P=33$ and $3\,610$ parameters the threshold of Theorem~\ref{thm:G6} is where that radius crosses $1$. At the iterates of a $33$-parameter network trained with weight decay, runs train up to $5.1\times$ above an $|H|$ reading of the threshold without diverging; six predictions written from a first run held on three new seeds. A two-sided bracket of the threshold with bounded rounding errors places training at $0.81$--$1.10$ times the threshold of the Gauss--Newton surrogate on networks with $3\,610$ and $12\,826$ parameters, in an exploratory measurement. Code and raw files reproduce every number; failed predictions are reported with the rest.
\end{abstract}

\section{Introduction}

With minibatch noise, the quantity that hovers at $2/\eta$ during training is the batch sharpness, the edge of stochastic stability \citep{andreyev2024eoss}, and the linear mean-square stability threshold lies below the deterministic $2(1+\beta)/\lambda_1$ \citep{wu2018sgd,andreyev2026momentum}. Without momentum, \citet{mulayoff2024exact} give the exact mean-square threshold of SGD on a quadratic in any dimension. With heavy-ball momentum, \citet{andreyev2026momentum} obtain the threshold to leading order, in the noise-dominated regime and for small $\eta$, and state that the precise instability criterion of SGD with momentum is an open question. In high-dimensional limits, \citet{paquette2021momentum} and \citet{lee2022minibatch} describe mini-batch SGD with momentum through Volterra equations; in the proof of their Proposition~2, \citet{lee2022minibatch} compute the norm of the Volterra kernel and state two sufficient conditions for it to be below one.

This draft gives the exact mean-square boundary of heavy-ball SGD on the quadratic model at every $\eta>0$ and $\beta\in[0,1)$, and identifies it with the conditions above where they apply:
\begin{itemize}
\item The exact stationary second-moment equation with a matrix curvature noise that need not commute with $H$ (Theorem~\ref{thm:F}), its exact relation $I-T_{\mathrm{slow}}=\mathcal S\circ\mathcal L$ to the Schur complement of \citet[Appendix~D]{andreyev2026momentum} and their remainder to leading order (Lemma~\ref{lem:remainder}), the reduction of the spectral radius to the symmetric matrices (Lemma~\ref{lem:russodye}) and that for $\rho(\mathcal T)\ge1$ the second moment grows at least linearly, which also meets \citet[Definition~A.1]{andreyev2026momentum}, the item they list as open for SGD with momentum (Corollary~\ref{cor:msinstab}; additive noise of unbounded support meets that definition at every $\eta$, so the moment bound is the content); its relation to batch sharpness at the threshold (Theorem~\ref{thm:modal-bridge}) and saturation in a slow-fast model (Theorem~\ref{thm:slow-fast-saturation}).
\item For random-feature minibatches, the boundary as the root of a rank-one secular equation with bounds by $\lambda_1$ and $\tr H$ (Theorem~\ref{thm:G}), the two plateaus and the Rademacher separation (Corollary~\ref{cor:plateaus}), and the gap between mean-square and almost-sure stability along the top mode (Proposition~\ref{prop:G3}).
\item A curvature-noise covariance $\Sigma_z\ne H$ (Theorem~\ref{thm:G5}), where a mode with small curvature and non-zero curvature noise drives the ceiling to zero, and arbitrary non-commuting per-sample Hessians (Theorem~\ref{thm:G6}), where the boundary is the reciprocal of the top eigenvalue of a symmetrisable operator; at $\beta=0$ this is Theorem~5 of \citet{mulayoff2024exact}. Appendix~\ref{app:more} bounds the gap between the threshold and its diagonal part, which closes as $\beta\to1$, and gives the burn-in of the second moment in Loewner order.
\item Measurements of the secular root against sample paths and against the exact second-moment operator, of the linearised minibatch iteration on two network checkpoints, of the operator of Theorem~\ref{thm:G6} on tethered network Hessians, and of the threshold at the iterates of training runs of a small network, replicated on new seeds, with predictions written before each run; and a certified two-sided bracket of the threshold at training checkpoints of networks with thousands of parameters (Section~\ref{sec:exp}).
\end{itemize}

\section{Setting and the one-dimensional ceiling}\label{sec:setting}

Let $\mathcal L(\Theta;B)$ be the loss of a minibatch $B$ at parameters $\Theta\in\R^P$, $\mathcal L(\Theta)$ its mean, and heavy-ball SGD $\Theta_{t+1}=\Theta_t+\beta(\Theta_t-\Theta_{t-1})-\eta\nabla\mathcal L(\Theta_t;B_t)$ with i.i.d.\ minibatches. Near a checkpoint $\Theta_0$ the quadratic model replaces $\nabla\mathcal L(\Theta;B_t)$ by $\nabla\mathcal L(\Theta_0;B_t)+H_{B_t}(\Theta-\Theta_0)$ with $H_{B_t}=\nabla^2\mathcal L(\Theta_0;B_t)$ and mean $H=\nabla^2\mathcal L(\Theta_0)$; with $w_t=\Theta_t-\Theta_0$, $\xi_t=\nabla\mathcal L(\Theta_0;B_t)-\nabla\mathcal L(\Theta_0)$ and $\nabla\mathcal L(\Theta_0)=0$ this is the recursion \eqref{eq:hbP} of Section~\ref{sec:lyapunov}. We write $\Delta w_t=w_{t+1}-w_t$ for the step, $x_t=\langle w_t,u_1\rangle$ and $\Delta x_t=x_{t+1}-x_t$ for their components along the top eigenvector. Let $H=\sum_k\lambda_ku_ku_k^\top$ with $\lambda_1\ge\lambda_2\ge\dots\ge\lambda_P$ (a gap $\lambda_1>\lambda_2$ is assumed where $u_1$ must be unique, in Propositions~\ref{prop:oneD} and~\ref{prop:G3} and in Corollary~\ref{cor:plateaus}), and allow a proximal tether $\frac\mu2\|\Theta-\Theta_0\|^2$ with $\mu\ge0$ and $\lambda_P+\mu>0$, which shifts every eigenvalue by $\mu$ and keeps the eigenvectors. Write $\kappa_\beta=\frac{1+\beta}{1-\beta}$ and $\etacrit^{(\mu)}(\beta)=2(1+\beta)/(\lambda_1+\mu)$, the deterministic threshold of heavy-ball gradient descent.

\begin{proposition}[One-dimensional mean-square ceiling; \citealt{wu2018sgd,andreyev2026momentum}]\label{prop:oneD}
On $\operatorname{span}(u_1)$, let the restoring force be evaluated on a fresh minibatch $B_t$ at each step, with directional curvature variance $\nu_B^2=\Var(u_1^\top\nabla^2\mathcal L(\Theta_0;B_t)u_1)$, and let additive noise of variance $\sigma_1^2>0$ drive the mode. The stationary second moment is finite if and only if $\eta<\etacrit^{(\mu),\mathrm{EoSS}}(B,\beta)$, where
\begin{equation}\label{eq:eoss}
\etacrit^{(\mu),\mathrm{EoSS}}(B,\beta)=\frac{\etacrit^{(\mu)}(\beta)}{1+\kappa_B(\beta)},\qquad
\kappa_B(\beta)=\kappa_\beta\,\frac{\nu_B^2}{(\lambda_1+\mu)^2}.
\end{equation}
\end{proposition}
This is the exact one-dimensional mean-square threshold of heavy-ball SGD of \citet[Section~4.1 and Appendix~B]{andreyev2026momentum}, $\eta_{\max}=2a(1+\beta)/\big(a^2+\kappa_\beta\sigma_b^2\big)$ with $a=\lambda_1+\mu$ and $\sigma_b=\nu_B$, obtained from the linear recursion of $(\E[x_t^2],\E[x_tx_{t-1}],\E[x_{t-1}^2])$; for $\beta=0$ it is the linear-stability condition $\eta<2\lambda_1/(\lambda_1^2+\nu_B^2)$ of \citet{wu2018sgd}, and it is the necessary condition~(17) of \citet{mulayoff2024exact} along $u_1$. \citet{andreyev2026momentum} use it to explain the two plateaus $2(1+\beta)/\eta$ and $2(1-\beta)/\eta$ of batch sharpness. Without curvature noise, the stationary normalised step variance of the mode is $\E[\Delta x_t^2]/\eta^2=\sigma_1^2/\big((1-\beta^2)(1-\eta(\lambda_1+\mu)/(2(1+\beta)))\big)$, by the Yule--Walker equations of the second-order autoregression. The rest of this draft is about $\R^P$.


\section{The exact second-moment equation in \texorpdfstring{$\R^P$}{R\^P} and the stochastic ceiling}\label{sec:lyapunov}

Proposition~\ref{prop:oneD} is one-dimensional. The next two statements give the $P$-dimensional second-moment equation of heavy-ball SGD with a matrix curvature noise that need not commute with $H$, and the exact mean-square ceiling on a quadratic with random-feature minibatches. Let $H=\sum_k\lambda_ku_ku_k^\top\succ0$ with $\lambda_1\ge\lambda_2\ge\dots\ge\lambda_P>0$, and
\begin{equation}\label{eq:hbP}
w_{t+1}=\big((1+\beta)I-\eta H_{B_t}\big)w_t-\beta w_{t-1}-\eta\xi_t,
\end{equation}
with $(H_{B_t},\xi_t)$ i.i.d.\ and independent of the past, $\E[H_{B_t}]=H$, $\E[\xi_t]=0$, $\E[\xi_t\xi_t^\top]=\Sigma_\xi$, and $\mathcal V(X)\triangleq\E[(H_{B_t}-H)X(H_{B_t}-H)]$. Write $\Sigma_0=\E[w_tw_t^\top]$, $D_\Theta=\E[\Delta w_t\Delta w_t^\top]$ and $[X,Y]=XY-YX$.

\begin{theorem}[Exact Lyapunov equation and virial identity]\label{thm:F}
In any covariance-stationary regime of \eqref{eq:hbP} with $\eta>0$ and $\E[w_t]=0$ (automatic when $H\succ0$, see the proof of Theorem~\ref{thm:G}), without any commutation assumption:
(1) $D_\Theta=\frac{\eta}{1+\beta}(H\Sigma_0+\Sigma_0H)$; in the eigenbasis of $H$, $\frac{1+\beta}{2\eta}\,\E[\langle\Delta w_t,u_k\rangle^2]/\E[\langle w_t,u_k\rangle^2]=\lambda_k$ for every mode with $\E[\langle w_t,u_k\rangle^2]>0$, identically in $\beta$, $\mathcal V$ and $\Sigma_\xi$.
(2) $(H\Sigma_0+\Sigma_0H)-\frac{\eta}{1+\beta}H\Sigma_0H+\frac{\eta\beta}{(1-\beta)(1-\beta^2)}[H,[H,\Sigma_0]]-\frac{\eta}{1-\beta}\mathcal V(\Sigma_0)=\frac{\eta}{1-\beta}\Sigma_\xi$.
(3) $u_j^\top[H,[H,\Sigma_0]]u_k=(\lambda_j-\lambda_k)^2(\Sigma_0)_{jk}$ vanishes on the diagonal; for $\mathcal V=0$ and any $\Sigma_\xi$ the normalized step cross-covariance $C_{jk}=\eta^{-2}(D_\Theta)_{jk}$ equals
\[
\frac{u_j^\top\Sigma_\xi u_k}{(1-\beta^2)\Big[1-\frac{\eta\lambda_j\lambda_k}{(1+\beta)(\lambda_j+\lambda_k)}+\frac{\eta\beta(\lambda_j-\lambda_k)^2}{(1-\beta)(1-\beta^2)(\lambda_j+\lambda_k)}\Big]},
\]
whose diagonal is the one-dimensional stationary step variance of Section~\ref{sec:setting} mode by mode.
\end{theorem}
At $\beta=0$, part~(2) has the form of the stationary covariance equation $(2C-\eta D)\operatorname{vec}\Sigma_0=\eta p\operatorname{vec}\Sigma_g$ of \citet[Theorem~12]{mulayoff2024exact}, which they state at a regular minimum for sampling without replacement, $p=(n-B)/(B(n-1))$; sampling with replacement gives $p=1/B$, their large-$n$ value, and with that $p$ the stability operator is theirs. Theorem~\ref{thm:F} extends that exact equation to heavy-ball momentum with a curvature noise that need not commute with $H$; what $\beta>0$ adds is the term $\frac{\eta\beta}{(1-\beta)(1-\beta^2)}[H,[H,\Sigma_0]]$, which vanishes on the diagonal of the eigenbasis and is stabilizing off it, and the factors $1/(1\pm\beta)$. \citet[Appendix~D]{andreyev2026momentum} reach the stability boundary of the same system through the Schur complement $T_{\mathrm{slow}}(\eta,\beta)$ of the second-moment operator, which they expand for small $\eta$ with a remainder they state to be of order $\eta^2\|H\|^2/(1-\beta)^3$ (their Appendix~H.1.1 evaluates this scale to $16.28$ on their network). Theorem~\ref{thm:F}(2) gives that boundary without the expansion, and the two are related exactly: $I-T_{\mathrm{slow}}=\mathcal S\circ\mathcal L$, where $\mathcal L(\Sigma)$ is $\eta/(1-\beta)$ times the left side of the equation of Theorem~\ref{thm:F}(2), curvature-noise term included, and $\mathcal S^{-1}(X)=X+\frac{\beta\eta}{(1-\beta)^2(1+\beta)}(HX+XH)$ does not depend on the curvature noise (Appendix~\ref{app:thmFG}). For $H\succeq0$, $\mathcal S$ is positive definite at every $\eta\ge0$, so $T_{\mathrm{slow}}$ exists at every $\eta$, not only small ones, and the two equations have the same solutions for every $\Sigma_\xi$ and are singular at the same $\eta$; $\mathcal S=I$ only at $\beta=0$. Their remainder follows exactly (Lemma~\ref{lem:remainder}): $R(\Sigma)=\frac{2\beta\eta^2}{(1-\beta)^3}H\Sigma H+R_3(\Sigma)$ with an explicit $R_3=O(\eta^3)$, so the scale $\eta^2\|H\|^2/(1-\beta)^3$ of their bound is attained, with constant $2\beta$. It is carried by $H\Sigma H$: the $H^2\Sigma$ and $\Sigma H^2$ parts of our double commutator are absorbed by the left factor $\mathcal S$. The two readings differ off the diagonal of the eigenbasis, where $(H\Sigma H)_{jk}=\lambda_j\lambda_k\Sigma_{jk}$ and $[H,[H,\Sigma]]_{jk}=(\lambda_j-\lambda_k)^2\Sigma_{jk}$, and an entry with $\lambda_k=0$ and no curvature noise on it decides between them: there the remainder vanishes identically, while the double-commutator entry does not. The expansion is also not uniform in $\beta$: the next term is not small against the leading one once $\beta\eta\|H\|\gtrsim(1-\beta)^2$ (Lemma~\ref{lem:remainder} gives the full condition). On the run of their Appendix~H.1.1 ($\eta=0.01$, $\beta=0.9$, $\lambda_{\max}=12.76$, where they evaluate the scale to $16.28$) the ratio $\beta\eta\lambda_{\max}/(1-\beta)^2$ is $11.5$ and the multiplier of $\mathcal S$ on the top mode is $0.076$, so neither their expansion nor the leading order of the remainder describes the operator there; equation~(2) gives the boundary at any $\eta$. Part~(1) is a matrix form of the kind of the first fluctuation--dissipation relation of \citet{yaida2019}; the co-diagonal additive-noise case of the stationary covariance is in \citet{gitman2019momentum}. Proofs in Appendix~\ref{app:thmFG}.

\begin{lemma}[The second-moment operator is decided on the symmetric matrices]\label{lem:russodye}
Let $\mathcal T(Q)=\E[F_tQF_t^\top]$ act on all real $2P\times2P$ matrices, with $F_t=\begin{pmatrix}(1+\beta)I-\eta H_{B_t}&-\beta I\\I&0\end{pmatrix}$ the transition of \eqref{eq:hbP}. Then $\mathcal T$ commutes with transposition and is positive, $\|\mathcal T^k\|=\|\mathcal T^k(I)\|$ for every $k$ in the norm induced by the spectral norm, and
\[
\rho(\mathcal T)=\lim_{k\to\infty}\|\mathcal T^k(I)\|^{1/k}=\rho\big(\mathcal T|_{\mathbb S^{2P}}\big),\qquad\|\mathcal T^k(Q)\|\le\|Q\|\,\|\mathcal T^k(I)\|\ \text{for every }Q .
\]
\end{lemma}
\begin{proof}
$\mathcal T(Q)^\top=\E[F_tQ^\top F_t^\top]=\mathcal T(Q^\top)$, so $\mathcal T$ preserves the symmetric and the antisymmetric matrices. $\mathcal T$ is in Kraus form, hence positive (even completely positive) on the complexified matrix algebra, and so is every $\mathcal T^k$; the Russo--Dye theorem gives $\|\mathcal T^k\|=\|\mathcal T^k(I)\|$. By Gelfand's formula, which does not depend on the norm, $\rho(\mathcal T)=\lim_k\|\mathcal T^k\|^{1/k}=\lim_k\|\mathcal T^k(I)\|^{1/k}\le\rho(\mathcal T|_{\mathbb S^{2P}})$, because $\mathcal T^k(I)\in\mathbb S^{2P}$; the reverse inequality is trivial. The last bound is $\|\mathcal T^k(Q)\|\le\|\mathcal T^k\|\,\|Q\|$.
\end{proof}
So the stability of the whole recursion is decided by the orbit of the identity, the antisymmetric part never dominates, and Theorem~\ref{thm:G} can be stated on $\mathbb S^{2P}$ without loss. The orbit also gives a matrix-free way to compute the ceiling: iterate $Q\mapsto\mathcal T(Q)/\|\mathcal T(Q)\|$ from $Q_0=I$, which is a power iteration on the positive semidefinite cone at the cost of one application of $\mathcal T$ per step.

\begin{corollary}[Mean-square instability criterion]\label{cor:msinstab}
In the setting of \eqref{eq:hbP} with $\Sigma_\xi\succ0$, let $\mathcal T$ be as in Lemma~\ref{lem:russodye}, and assume either that $\xi_t$ is independent of $H_{B_t}$ or that $\E[w_0]=\E[w_{-1}]=0$. Then $\sup_t\E\|w_t\|^2<\infty$ if and only if $\rho(\mathcal T)<1$. If $\rho(\mathcal T)\ge1$, then $\E\|w_t\|^2+\E\|w_{t-1}\|^2\ge c\,t$ for every $t\ge1$ with a constant $c>0$, so $\sup_t\E\|w_t\|^2=\infty$, and for every compact $K\subset\R^P$ there is a finite $T$ with $\Pr(w_T\notin K)>0$. For the random-feature minibatches of Theorem~\ref{thm:G} below, $\rho(\mathcal T)\ge1$ if and only if $\eta\ge\etacrit^{\mathrm{EoSS},P}(\beta,B)$.
\end{corollary}
That bounded second moments of a linear recursion with i.i.d.\ coefficients are equivalent to $\rho(\E[F_t\otimes F_t])<1$ is standard \citep[Chapter~3]{costa2005mjls}, and it is the criterion used for SGD without momentum by \citet{ma2021linear} and \citet{mulayoff2024exact}. The corollary records it for \eqref{eq:hbP}, with the divergence statement used below. Its last conclusion has the form of the escape condition of \citet[Definition~A.1]{andreyev2026momentum}, but on this model that condition does not single out $\rho(\mathcal T)$: when $\xi_t$ has unbounded support (for instance Gaussian), $\Pr(w_1\notin K)>0$ for every compact $K$ at every $\eta$, also where $\rho(\mathcal T)<1$; with Gaussian features one minibatch has unbounded curvature, so the same holds from every start $w_0\ne0$ without additive noise; and $\rho(\mathcal T)$ does not depend on the initial parameter, so as a function $f(\theta_0)$ it is constant and the definition then asks for escape from every start; read as escape in probability, $\Pr(w_t\notin K)\to1$ for every compact $K$, it does not follow from $\rho(\mathcal T)\ge1$, since between the mean-square and the almost-sure thresholds of Proposition~\ref{prop:G3} typical trajectories decay. The content of the corollary is the moment bound. Without additive noise, $\Sigma_\xi=0$, the same eigenmatrix gives $\langle W,Q_t\rangle=\rho(\mathcal T)^t\langle W,Q_0\rangle$, so $\E\|w_t\|^2+\E\|w_{t-1}\|^2\ge\rho(\mathcal T)^t\langle W,Q_0\rangle/\tr W$ for every start with $\langle W,Q_0\rangle>0$, for instance a deterministic $w_0\notin\ker W_{11}$ with $w_{-1}=0$: at $\rho(\mathcal T)\ge1$ the second moment from such a start does not decay, and for $\rho(\mathcal T)>1$ it grows geometrically. It needs $\Sigma_\xi\succ0$: for the homogeneous recursion, the one of their Appendix~D, $\rho(\mathcal T)\ge1$ does not force escape (the start $w_0=w_{-1}=0$ stays at $0$), and the bound $\langle W,Q_t\rangle\ge tg$ is what the additive noise supplies.
\begin{proof}
Let $s_t=(w_t,w_{t-1})$, $Q_t=\E[s_ts_t^\top]$ and $G=\mathrm{blkdiag}(\eta^2\Sigma_\xi,0)$. The cross terms $\E[F_ts_t(\xi_t^\top,0)]$ vanish: $s_t$ is independent of $(F_t,\xi_t)$, and either $F_t$ is independent of $\xi_t$, which has mean $0$, or $\E[s_t]=0$ for every $t$, because $\E[w_t]$ follows the deterministic heavy-ball recursion with $H$ from $0$. So $Q_{t+1}=\mathcal T(Q_t)+G$. If $\rho(\mathcal T)<1$, $Q_t=\mathcal T^t(Q_0)+\sum_{s<t}\mathcal T^s(G)$ is bounded. The adjoint $\mathcal T^*(Y)=\E[F_t^\top YF_t]$ maps the positive semidefinite cone of $\mathbb S^{2P}$ into itself, so by the Krein--Rutman theorem in finite dimension \citep[Thm.~1.3.2]{berman1994} it has an eigenvector $W\succeq0$, $W\ne0$, with eigenvalue $\rho(\mathcal T^*)=\rho(\mathcal T)$. Its upper-left block $W_{11}$ is not zero: if it were, $W=\mathrm{blkdiag}(0,W_{22})$, and $F_t^\top WF_t=\mathrm{blkdiag}(W_{22},0)$ for every $F_t$, so $\rho W=\mathcal T^*(W)$ forces $W_{22}=0$. Hence $g=\langle W,G\rangle=\eta^2\tr(W_{11}\Sigma_\xi)>0$ and $\langle W,Q_{t+1}\rangle=\rho(\mathcal T)\langle W,Q_t\rangle+g$. For $\rho(\mathcal T)\ge1$ this gives $\langle W,Q_t\rangle\ge tg$, and $\langle W,Q_t\rangle\le\tr W\,\tr Q_t=\tr W\,(\E\|w_t\|^2+\E\|w_{t-1}\|^2)$, so the bound holds with $c=g/\tr W$. If $w_T\in K$ almost surely for every $T$, then $\E\|w_T\|^2\le\max_K\|w\|^2$ for every $T$, a contradiction. The last statement is Theorem~\ref{thm:G}, whose proof shows $\rho(\mathcal T)<1$ for $\eta<\etacrit^{\mathrm{EoSS},P}$ and the absence of a finite stationary second moment otherwise; the latter excludes $\rho(\mathcal T)<1$, since then $\sum_s\mathcal T^s(G)$ would be one.
\end{proof}
With three equiprobable non-commuting $2\times2$ minibatch Hessians and $\beta=0.9$, the exact recursion stays bounded at $0.9\,\eta^\star$ ($\rho(\mathcal T)=0.9886$), and $\langle W,Q_t\rangle\ge tg$ holds for $t\le3000$ at $\eta^\star$ and $1.02\,\eta^\star$ (\texttt{paper/checks/ms-instability-criterion.mjs}, which also matches the recursion to a Monte Carlo run).
\paragraph{The open question of \citet{andreyev2026momentum}.} For SGD without momentum, \citet{andreyev2024eoss} prove that on the quadratic approximation batch sharpness above $2/\eta$ forces divergence. For momentum, \citet[Section~3]{andreyev2026momentum} write that ``for SGDM and SGDN, we do not have an analogous theoretical result'' and, in a later paragraph, that ``developing such a theory for momentum methods, including identifying the precise instability criterion and proving that it is saturated under progressive sharpening, remains an important open question''; the first item of their future work (Section~6.2) is ``finding a valid instability criterion for SGDM and SGDN''. Their Definition~A.1 fixes an open set $U$ of parameters, where a local model of the loss is trusted, and calls $f:U\to\R$ a valid instability criterion with threshold $c$ on $U$ if $f(\theta_0)>c$ implies that the iterates leave every compact subset of $U$, equivalently that for every compact $K\subset U$ containing $\theta_0$ there is a finite $T$ with $\Pr(\theta_T\notin K\mid\theta_0)>0$; on the quadratic model of this paper $U=\R^P$. Their Appendix~D writes the homogeneous second-moment recursion of linearized SGDM, states that its mean-square stability is $\rho(T_{\mathrm{HB}})<1$, and reduces this to $\rho(T_{\mathrm{slow}}(\eta,\beta))<1$ for sufficiently small $\eta$ (their eq.~(31)); it does not show that $\rho\ge1$ forces escape in the sense of their Definition~A.1, and for the homogeneous recursion it does not (the start $w_0=w_{-1}=0$ stays at $0$). Their Theorem~4.1 keeps the leading order of $T_{\mathrm{slow}}$, ``in the noise-dominated regime and for sufficiently small $\eta$'': it is the SGD condition of \citet{ma2021linear} at $\eta_{\mathrm{eff}}=\eta/(1-\beta)$, the step-size rescaling that \citet{yuan2016momentum} establish for momentum methods with a small constant step, $\rho(I-\eta_{\mathrm{eff}}K_H+\eta_{\mathrm{eff}}^2G)<1$ in the notation of Lemma~\ref{lem:remainder}, and the dropped remainder is bounded only by $\|R\|\lesssim\eta^2\|H\|^2/(1-\beta)^3$. For heavy-ball SGD on the quadratic model this paper adds, to the linear mean-square part of their question: (i) a closed stationary equation that decides the boundary at every $\eta>0$ and $\beta\in[0,1)$, Theorem~\ref{thm:F}(2), with no small-$\eta$ and no noise-dominated assumption, and its exact relation $I-T_{\mathrm{slow}}=\mathcal S\circ\mathcal L$ to their Schur complement; (ii) the generality of their Appendix~D in the minibatch Hessians, which need not commute with each other or with $H$, without the small-$\eta$ expansion; (iii) that for $\rho(\mathcal T)\ge1$ the second moment grows at least linearly, $\E\|w_t\|^2+\E\|w_{t-1}\|^2\ge ct$ (Corollary~\ref{cor:msinstab}), whose probabilistic conclusion has the form of their Definition~A.1 but is also met below $\rho(\mathcal T)=1$ (additive noise of unbounded support meets it at every $\eta$), so it does not settle the first item of their Section~6.2; the content is the moment bound, the heavy-ball analogue of the divergence theorem of \citet{andreyev2024eoss} for batch sharpness at $\beta=0$, stated for $\rho(\mathcal T)$ and not for batch sharpness; the argument is the standard one for linear systems with i.i.d.\ coefficients \citep{costa2005mjls} and needs the additive noise $\Sigma_\xi\succ0$, which their model, an interpolating minimum, does not have (their footnote~5 states that relaxing interpolation changes the steady-state variance but not the stability criteria); (iv) their remainder exactly, $R=\frac{2\beta\eta^2}{(1-\beta)^3}H\Sigma H+R_3$, which attains the scale of their bound and is not small once $\beta\eta\|H\|\gtrsim(1-\beta)^2$ (Lemma~\ref{lem:remainder}), a ratio equal to $11.5$ on the run where they evaluate it; (v) for random-feature minibatches, the boundary as the root of a secular equation (Theorem~\ref{thm:G}), whose two limits take the values of their plateaus $2(1+\beta)/\eta$ and $2(1-\beta)/\eta$, through $\lambda_1$ and $\tr H/B$ (Corollary~\ref{cor:plateaus}). For SGD with Nesterov momentum, the exact mean-square criterion on random quadratics in any dimension is the outer-momentum criterion of \citet[Theorem~1]{sharma2026outer} with one local step. Concurrently and independently, \citet[Theorem~F.1]{cai2026llm} give the exact mean-square criterion of the homogeneous heavy-ball recursion with i.i.d.\ curvature noise that need not commute with the mean; it is the threshold of Theorem~\ref{thm:G6}, and their comparison of two momentum values (their eq.~(67)) gives the monotonicity of Corollary~\ref{cor:etaeff} without strictness. Whether training saturates the criterion is measured, on one $33$-parameter network and with the quasi-static diagnostic of Section~\ref{sec:phase5}, in Phases~5 and 5b: the runs cross the $|H|$ reading between $\eta=0.1$ and $0.15$ and train up to $5.1\times$ above it without diverging, at checkpoints where $H$ is indefinite and the linearised iteration has no mean-square stable step. On the networks with $P=3\,610$ and $12\,826$ parameters, Phase~6 brackets the threshold of the Gauss--Newton surrogate from both sides and finds training at it, within $0.81$--$1.10\,\eta$ over the window, for $\eta\ge0.02$ (Section~\ref{sec:phase6}). Theorem~\ref{thm:modal-bridge} relates the criterion to batch sharpness: at the threshold an identity puts a modal curvature between the two SGDM plateau values $2(1-\beta)/\eta^\star$ and $2(1+\beta)/\eta^\star$, a stationary batch sharpness lies above it for positive semidefinite minibatch Hessians and, in the random-feature model, reaches the plateaus in the two batch limits, and a variational test gives sufficient instability conditions through a fixed-checkpoint batch sharpness and the noise share; Theorem~\ref{thm:slow-fast-saturation} proves saturation in a quasi-stationary slow-fast model, with an equilibrium at distance $\Theta(\sigma_\xi^2/v_{\mathrm{PS}})$ below the threshold, and a secant equilibrium above the tangent threshold. Not answered here: the dynamics of training beyond the quadratic and slow-fast models, which is what bounds those runs, and a threshold on the batch sharpness of \citet[Definition~3.1]{andreyev2026momentum} alone. The criterion is also a mean-square one: between the mean-square and the almost-sure thresholds of Proposition~\ref{prop:G3} the second moment diverges while typical trajectories decay (the failed cell of Phase~1, Section~\ref{sec:phase1}).

\begin{theorem}[Secular equation of the stochastic ceiling]\label{thm:G}
Let $H_{B_t}=\frac1B\sum_{b=1}^Bz_{i_b}z_{i_b}^\top$ with $z_i=H^{1/2}\omega_i$, where the coordinates $\omega_{i,k}$ (in the eigenbasis of $H$) are independent with mean $0$, variance $1$ and $\E[\omega_{i,k}^4]=3+\kappa_{\mathrm{feat}}$, $\kappa_{\mathrm{feat}}\ge-2$ ($0$ for Gaussian features, $-2$ for Rademacher features, for which $\omega^4\equiv1$), the $B$ draws i.i.d.; let $\sigma_{\xi,k}^2=u_k^\top\Sigma_\xi u_k>0$, $\kappa_\beta=\frac{1+\beta}{1-\beta}$, $\alpha_{\mathrm{diag}}=1+\kappa_\beta\frac{1+\kappa_{\mathrm{feat}}}{B}$, $\Lambda=\mathrm{diag}(\lambda_k)$ and $\lambda=(\lambda_1,\dots,\lambda_P)^\top$. Then $\Var(u_k^\top H_{B_t}u_j)=\frac{\lambda_k\lambda_j}{B}\big(1+(1+\kappa_{\mathrm{feat}})\delta_{kj}\big)$, the off-diagonal entries of $\Sigma_0$ decouple from each other and from the diagonal, the diagonal normalized step variances $c=(C_{kk})_k$ solve
\begin{equation}\label{eq:mmatrix}
\Big[I-\tfrac{\eta}{2(1+\beta)}\big(\alpha_{\mathrm{diag}}\Lambda+\tfrac{\kappa_\beta}{B}\lambda\mathbf 1^\top\big)\Big]c=\frac{\sigma_\xi^2}{1-\beta^2},
\end{equation}
and a finite stationary second moment of $(w_t,w_{t-1})$ exists if and only if $\eta<\etacrit^{\mathrm{EoSS},P}(\beta,B)$, the unique root in $\big(0,\frac{2(1+\beta)}{\lambda_1\max(\alpha_{\mathrm{diag}},1)}\big]$ (the right end is attained only without curvature noise on the top mode, as for $P=1$ and Rademacher features) of
\begin{equation}\label{eq:secular}
\frac{\eta}{2(1-\beta)B}\sum_{k=1}^P\frac{\lambda_k}{1-\frac{\eta\lambda_k}{2(1+\beta)}\alpha_{\mathrm{diag}}}=1 .
\end{equation}
With $\mu^\star=2(1+\beta)/\etacrit^{\mathrm{EoSS},P}(\beta,B)$,
\begin{equation}\label{eq:sandwich}
\max(\alpha_{\mathrm{diag}},0)\lambda_1+\kappa_\beta\frac{\tr H}{B}\ \ge\ \mu^\star\ \ge\ \max\Big(\alpha_{\mathrm{diag}}\lambda_1+\kappa_\beta\frac{\lambda_1}{B},\ \min(\alpha_{\mathrm{diag}},0)\lambda_1+\kappa_\beta\frac{\tr H}{B}\Big),
\end{equation}
with $\alpha_{\mathrm{diag}}\lambda_1+\kappa_\beta\lambda_1/B=\lambda_1(1+\kappa_\beta(2+\kappa_{\mathrm{feat}})/B)\ge\lambda_1$. For $P=1$, \eqref{eq:secular} is the ceiling of Proposition~\ref{prop:oneD} with $\nu_B^2=\lambda_1^2(2+\kappa_{\mathrm{feat}})/B$.
\end{theorem}
\begin{corollary}[Two plateaus; the Rademacher separation]\label{cor:plateaus}
Write $\varepsilon_B=\kappa_\beta/B$ and $T(\mu)=\sum_{k\ge2}\lambda_k/(\mu-\alpha_{\mathrm{diag}}\lambda_k)$. At the ceiling, $\mu^\star=2(1+\beta)/\etacrit^{\mathrm{EoSS},P}$ satisfies exactly $\mu^\star=\alpha_{\mathrm{diag}}\lambda_1+\varepsilon_B\lambda_1/\big(1-\varepsilon_BT(\mu^\star)\big)$. If $\alpha_{\mathrm{diag}}>0$, $\lambda_1>\lambda_2$ and, with $S_\perp=\sum_{k\ge2}\lambda_k/(\lambda_1-\lambda_k)$, $\varepsilon_BS_\perp/\alpha_{\mathrm{diag}}\le\frac12$ (large batch with a spectral gap), then $\lambda_1=\frac{2(1+\beta)}{\eta}\big(1-\varepsilon_B(2+\kappa_{\mathrm{feat}})+\varrho\big)$ at $\eta=\etacrit^{\mathrm{EoSS},P}$, with $-2\varepsilon_B^2S_\perp/\alpha_{\mathrm{diag}}\le\varrho\le\varepsilon_B^2\big(2+\kappa_{\mathrm{feat}}+2\varepsilon_BS_\perp/\alpha_{\mathrm{diag}}\big)^2$. The gap is needed, not an artifact of the proof: with $\lambda_1=\lambda_2$ the top eigenvalue on those two modes is $\lambda_1(1+\varepsilon_B(3+\kappa_{\mathrm{feat}}))$, so the first-order coefficient itself changes. As $\kappa_\beta\tr H/(B\lambda_1)\to\infty$ with $\lambda_1/\tr H\to0$ (small batch, high intrinsic rank), $\tr H/B=\frac{2(1-\beta)}{\eta}(1+o(1))$ at the ceiling. The two empirical plateaus of \citet[eqs.~1--2]{andreyev2026momentum}, $2(1+\beta)/\eta$ and $2(1-\beta)/\eta$, are the values of the two limits of \eqref{eq:secular}, taken by $\lambda_1$ and by the batch trace-sharpness $\tr H/B$ respectively, not by batch sharpness itself; Theorem~\ref{thm:modal-bridge}(2) gives the two limits of a stationary batch sharpness, and the large-batch plateau of SGD with Nesterov momentum in \citet{andreyev2026momentum}, $2(1+\beta)/(\eta(1+2\beta))$, is not covered. For Rademacher features the directional curvature along $u_1$ is deterministic, $\Var(u_1^\top H_{B_t}u_1)=0$, so the one-dimensional ceiling of Proposition~\ref{prop:oneD} is $2(1+\beta)/\lambda_1$ for every $B$, while by \eqref{eq:sandwich}, with $\tr H_\perp=\tr H-\lambda_1$, $\etacrit^{\mathrm{EoSS},P}\le2(1+\beta)/(\lambda_1+\kappa_\beta\tr H_\perp/B)$ for $B\le\kappa_\beta$ and $\etacrit^{\mathrm{EoSS},P}\le2(1+\beta)/\max(\lambda_1,\kappa_\beta\tr H/B)$ for $B>\kappa_\beta$: the cross-mode couplings $\lambda_1\lambda_k/B$ lower the ceiling, and a one-dimensional probe along $u_1$ cannot see them.
\end{corollary}
The same structure, a cone-preserving second-moment operator with a rank-one global coupling, the Krein--Rutman theorem, a secular equation and bounds by $\tr H$ and $\lambda_{\max}$ with the factor $2(1-\beta)$, appears in the mean-square stability of full-batch zeroth-order methods with Gaussian two-point directions \citep{song2026zeroth}: ZO-GD, ZO-GDM, and ZO-Adam under the assumption that the preconditioner commutes with $H$; mini-batch zeroth-order methods are left there as future work. There the randomness is the search direction and the Hessian is the full-batch one. Here it is the minibatch: the sampled Hessians need not commute with each other or with $H$ (Theorem~\ref{thm:F}, Lemma~\ref{lem:russodye}), the coupling carries the feature kurtosis through $\alpha_{\mathrm{diag}}$, and the curvature-noise covariance need not be $H$ (Theorem~\ref{thm:G5}). A small-batch edge of stability governed by a trace rather than by the top Hessian eigenvalue also appears in the high-dimensional model of \citet{agarwala2024high}, through the trace of the neural tangent kernel.
\begin{proposition}[Mean-square against almost-sure stability for all $\beta\in[0,1)$; a distribution-free trace identity]\label{prop:G3}
(1) For the homogeneous Heavy-Ball recursion $x_{t+1}=(1+\beta-\eta(\lambda_1+\zeta_t))x_t-\beta x_{t-1}$ along $u_1$, with $\beta\in[0,1)$, $\E\zeta_t=0$ and $\Var\zeta_t=\nu_{11,B}^2<\infty$, the mean-square threshold is exactly $\frac{2(1+\beta)}{\lambda_1}(1+\kappa_\beta\nu_{11,B}^2/\lambda_1^2)^{-1}=\frac{2(1+\beta)}{\lambda_1}\big(1-\kappa_\beta\nu_{11,B}^2/\lambda_1^2+\varrho\big)$, where $0\le\varrho\le\big(\kappa_\beta\nu_{11,B}^2/\lambda_1^2\big)^2$. Let $u_t=2\kappa_\beta\delta+2\kappa_\beta(1+\delta)\zeta_t/\lambda_1$ at $\eta=\frac{2(1+\beta)}{\lambda_1}(1+\delta)$, and write $m_{3,B}=\E|\zeta_t|^3$. Set $\varepsilon_0(0)=\frac12$ and $\varepsilon_0(\beta)=\frac{1-\beta}{8(1+\beta)}$ for $\beta\in(0,1)$. If $|u_t|\le\varepsilon_0(\beta)$ almost surely, $m_{3,B}<\infty$ and $|\delta|\le\frac12$, then the top Lyapunov exponent $\gamma_{\mathrm{Lyap}}(\eta,\beta)=\lim_{t\to\infty}\frac1t\log\|M_t\cdots M_1\|$ of the companion matrix $M_t$ \citep{furstenberg1960products,caughey1965almost,khasminskii1967necessary,kozin1969survey,ariaratnam1979random,arnold1984formula} satisfies $\gamma_{\mathrm{Lyap}}(\eta,\beta)=2\kappa_\beta\delta-2\kappa_\beta^2\nu_{11,B}^2/\lambda_1^2+O\big(\delta^2+m_{3,B}/\lambda_1^3+\nu_{11,B}^4/\lambda_1^4\big)$, with an implied constant depending only on $\beta$. If in addition $|\zeta_t|\le c_0(\beta)\lambda_1$ almost surely for a positive $c_0(\beta)$ depending only on $\beta$, then on the range $|\delta|\le c_1(\beta)$ the exponent is negative below and positive above a window of width $O(m_{3,B}/\lambda_1^3+\nu_{11,B}^4/\lambda_1^4)$ around $\delta=\kappa_\beta\nu_{11,B}^2/\lambda_1^2$, so the almost-sure threshold there is $\etacrit^{\mathrm{a.s.}}(\beta)=\frac{2(1+\beta)}{\lambda_1}\big(1+\kappa_\beta\nu_{11,B}^2/\lambda_1^2+O(m_{3,B}/\lambda_1^3+\nu_{11,B}^4/\lambda_1^4)\big)$. At $\beta=0$, for each fixed $p>0$ there is a $c_p>0$ such that $|\zeta_t|\le c_p\lambda_1$ almost surely implies that the threshold of $\E|1-\eta(\lambda_1+\zeta_t)|^p=1$ is $\frac2{\lambda_1}\big(1-(p-1)\nu_{11,B}^2/\lambda_1^2+O_p(m_{3,B}/\lambda_1^3+\nu_{11,B}^4/\lambda_1^4)\big)$. At $\beta=0$ the support condition can be replaced by a density bound near the singular multiplier: if $|\delta|\le\frac18$, $m_{3,B}<\infty$ and the law of $1+u_t$ has a density bounded by $M_{\mathrm{sing}}$ on $[-\frac12,\frac12]$, then $\big|\gamma_{\mathrm{Lyap}}(\eta,0)-(\E u_t-\frac12\E u_t^2)\big|\le11\,\E|u_t|^3+(1+\log2)M_{\mathrm{sing}}$, and the root at $\beta=0$ is $\frac2{\lambda_1}\big(1+\nu_{11,B}^2/\lambda_1^2+O(m_{3,B}/\lambda_1^3+\nu_{11,B}^4/\lambda_1^4+M_{\mathrm{sing}})\big)$. With $\varepsilon=\nu_{11,B}^2/\lambda_1^2$, Gaussian $\zeta_t$ has $M_{\mathrm{sing}}\le(2\pi\frac{49}{16}\varepsilon)^{-1/2}e^{-1/(162\varepsilon)}$, and $\zeta_t=\lambda_1(\chi^2_B/B-1)$ has $M_{\mathrm{sing}}\le\frac{18}{7}\sqrt{B/(2\pi)}\,e^{-0.0056B}$; both are smaller than every power of $1/B$ as $B\to\infty$, but not small at moderate $B$. A finite third moment alone does not give these remainders: if the law of $u_t$ charges every neighbourhood of $-1$ with an unbounded density, then $\E\big[|\log|1+u_t||\,\mathbf 1_{|1+u_t|\le1/2}\big]/m_{3,B}$ can be arbitrarily large. For $\beta\in(0,1)$ with unbounded $\zeta_t$ the remainder is left open. If $\zeta_t=0$ almost surely, both thresholds equal $2(1+\beta)/\lambda_1$, which is the Rademacher case $\nu_{11,B}^2=0$. Rosenthal's inequality gives $\nu_{11,B}^2=O(1/B)$ and $m_{3,B}=O(B^{-3/2})$ for an average of $B$ i.i.d.\ summands with a finite third moment; those rates enter the remainders only when $|u_t|\le\varepsilon_0(\beta)$ almost surely. (2) In the setting of Theorem~\ref{thm:F} with $H_{B_t}=\frac1B\sum_bH_{i_b}$, $H_i\succeq0$ i.i.d., $\E H_i=H$ and no further structure, every stationary $\Sigma_0$ satisfies $\tr\big([2H-\frac{\eta}{1+\beta}H^2-\frac{\eta}{(1-\beta)B}(\E H_i^2-H^2)]\Sigma_0\big)=\frac{\eta}{1-\beta}\tr\Sigma_\xi$, since $\tr[H,[H,\Sigma_0]]=0$; multiplying by $H^{-1}$ before the trace gives $\tr\big([2I-\frac{\eta}{1+\beta}(H+\frac{\kappa_\beta}{B}\E[(H_i-H)H^{-1}(H_i-H)])]\Sigma_0\big)=\frac{\eta}{1-\beta}\tr(H^{-1}\Sigma_\xi)>0$ with a positive semidefinite correction, hence the distribution-free necessary condition $\etacrit^{\mathrm{EoSS},P}\le2(1+\beta)/\lambda_{\min}\big(H+\frac{\kappa_\beta}{B}\E[(H_i-H)H^{-1}(H_i-H)]\big)$ for every $B\ge1$; and if $\E[H_iXH_i]\succeq c_0\tr(HX)H$ for $X\succeq0$ and $B\ge\kappa_\beta$, $\etacrit^{\mathrm{EoSS},P}\le2(1-\beta)B/(c_0\tr H)$ (equality is possible: $H_i\in\{0,2\}$, $\beta=0$, $B=1$, $c_0=2$). Proof in Appendix~\ref{app:thmFG}.
\end{proposition}
Numerically, $M_{\mathrm{sing}}\approx0.4$ at $B=64$ and below $10^{-6}$ from $B=1024$ (\texttt{paper/checks/g3-singular-density.mjs}). For $\lambda_1=1$, $B=64$, $\beta=0$ and $\zeta_t=\chi^2_B/B-1$, the mean-square formula gives $1.9394$; the root of $\E\log|1-\eta\chi^2_B/B|=0$ is $2.0688$ by quadrature, against the first-order values $1.9375$ and $2.0625$, and a Gaussian $\zeta_t$ of the same variance has quadrature root $2.078$ (\texttt{paper/checks/g3-lyapunov-remainder.mjs}); these two roots are values of the integral.
Part~(1) is the mechanism behind the one failed cell of Phase~1 (Section~\ref{sec:phase1}): a finite sample-path harness follows typical trajectories, so its threshold can lie above the mean-square one, toward the almost-sure one, and Theorem~\ref{thm:G} is about the mean-square one. Jensen's inequality puts the almost-sure threshold at or above the mean-square one; under the support bound they differ by $\pm\kappa_\beta\nu_{11,B}^2/\lambda_1^2$ to first order.

\begin{theorem}[Secular equation for a curvature-noise covariance that is not the Hessian; flat-mode collapse]\label{thm:G5}
In the setting of Theorem~\ref{thm:F} let $H=\sum_k\lambda_ku_ku_k^\top\succ0$ and $H_{B_t}=H+\frac1B\sum_{b}(z_{i_b}z_{i_b}^\top-\Sigma_z)$ with $z_i=\Sigma_z^{1/2}\omega_i$, $\Sigma_z=\operatorname{diag}(s_1,\dots,s_P)\succeq0$ in the eigenbasis of $H$ (for instance a Gauss--Newton covariance $\Sigma_z=H_{\mathrm{GN}}\neq H$ that commutes with $H$), and $\omega_{i,k}$ independent, standardised, with $\E\omega_{i,k}^4=3+\kappa_{\mathrm{feat}}$, $\kappa_{\mathrm{feat}}\ge-2$. Write $\rho_k=s_k^2/\lambda_k$, $d_k=\lambda_k+\kappa_\beta\frac{1+\kappa_{\mathrm{feat}}}{B}\rho_k$ and $\mathcal T_{\mathrm{flat}}=\sum_k\rho_k=\tr(H^{-1}\Sigma_z^2)$ ($\Sigma_z=H$ gives $\rho_k=\lambda_k$, $d_k=\alpha_{\mathrm{diag}}\lambda_k$ and Theorem~\ref{thm:G}).
(1) In the eigenbasis, $(\mathcal V(X))_{kk}=\frac{s_k}{B}\sum_js_jX_{jj}+\frac{1+\kappa_{\mathrm{feat}}}{B}s_k^2X_{kk}$ and $(\mathcal V(X))_{jk}=\frac{s_js_k}{B}X_{kj}$ for $j\ne k$ ($=\frac{s_js_k}{B}X_{jk}$ for symmetric $X$), so the diagonal of the stationary equation closes on $\gamma_{0,k}=(\Sigma_0)_{kk}$ and reads $2\lambda_k\gamma_{0,k}-\frac{\eta}{1+\beta}\lambda_k^2\gamma_{0,k}-\frac{\eta}{1-\beta}\big(\frac{s_k}{B}\sum_js_j\gamma_{0,j}+\frac{1+\kappa_{\mathrm{feat}}}{B}s_k^2\gamma_{0,k}\big)=\frac{\eta}{1-\beta}u_k^\top\Sigma_\xi u_k$.
(2) For $\Sigma_\xi\succ0$, a finite stationary second moment exists if and only if $\eta<\eta^\star\triangleq2(1+\beta)/\lambda_{\max}(\widetilde{\mathcal M})$, $\widetilde{\mathcal M}=\operatorname{diag}(d_k)+\frac{\kappa_\beta}{B}vv^\top$, $v_k=\sqrt{\rho_k}$. With $\mathcal A=\{k:s_k>0\}$, a mode outside $\mathcal A$ is decoupled ($d_k=\lambda_k$, $v_k=0$), so $\eta^\star$ is the smaller of $\min_{k\notin\mathcal A}2(1+\beta)/\lambda_k$ and, when $\mathcal A\ne\emptyset$, the unique root in $(0,\min_{k\in\mathcal A}2(1+\beta)/\max(d_k,\lambda_k)]$ of the secular equation (whose terms with $k\notin\mathcal A$ vanish)
\[
\frac{\eta}{2(1-\beta)B}\sum_{k=1}^P\frac{s_k^2/\lambda_k}{1-\frac{\eta\lambda_k}{2(1+\beta)}-\frac{\eta(1+\kappa_{\mathrm{feat}})}{2(1-\beta)B}\frac{s_k^2}{\lambda_k}}=1.
\]
(3) With $\mu^\star=2(1+\beta)/\eta^\star$, $\max_kd_k+\frac{\kappa_\beta}{B}\mathcal T_{\mathrm{flat}}\ge\mu^\star\ge\max\big\{\max_k(\lambda_k+\kappa_\beta\frac{2+\kappa_{\mathrm{feat}}}{B}\rho_k),\ \min_kd_k+\frac{\kappa_\beta}{B}\mathcal T_{\mathrm{flat}}\big\}$; in particular $\eta^\star\le2(1+\beta)/\lambda_1$ for every $\kappa_{\mathrm{feat}}\ge-2$, and $\eta^\star\le2(1-\beta)B/\mathcal T_{\mathrm{flat}}$ whenever $\min_kd_k\ge0$ (for instance $\kappa_{\mathrm{feat}}\ge-1$); without that condition the second bound can fail. For every $\kappa_{\mathrm{feat}}\ge-2$ and $P\ge2$, $\eta^\star\le\min_{j\ne k}\frac{2(1-\beta)B\sqrt{\lambda_j\lambda_k}}{s_js_k}$, so a mode with $\lambda_k\to0$ at fixed $s_k>0$, with some $s_j>0$, $j\ne k$, drives $\eta^\star$ to zero as $O(\sqrt{\lambda_k})$, and with $P=1$ and Rademacher features there is no collapse: curvature noise in a flat direction has no restoring force to balance it, a mechanism excluded by $\Sigma_z=H$. A tether $\frac\mu2\|w-\Theta_0\|^2$ leaves $s_k$ unchanged and replaces $\lambda_k$ by $\lambda_k+\mu$, so $\rho_k(\mu)\le s_k^2/\mu$ and $\mathcal T_{\mathrm{flat}}(\mu)\le\tr(\Sigma_z^2)/\mu$. Proof in Appendix~\ref{app:thmFG}.
\end{theorem}
Examples for part (3): with Rademacher features, $\Sigma_z=H=I_2$, $B=1$ and $\beta=\frac12$, $\eta^\star=0.75$ exceeds the bound $2(1-\beta)B/\mathcal T_{\mathrm{flat}}=0.5$, which therefore needs $\min_kd_k\ge0$; with $\lambda=(1,10^{-4})$, $s=(1,1)$, $\beta=\frac12$, $B=1$ and Rademacher features, $\eta^\star=0.00998$ against the bound $0.01$. On the network checkpoints of Phase~2 the negative directions of $H$ dominate instead of flat modes (Section~\ref{sec:phase4}).

\begin{theorem}[Exact mean-square threshold in $\R^P$ for arbitrary non-commuting per-sample Hessians]\label{thm:G6}
In the setting of Theorem~\ref{thm:F}, let $H=\frac1N\sum_{i=1}^NH_i\succ0$ (or $H+\mu I\succ0$ under a tether $\mu>0$) have eigenvalues $\lambda_1\ge\dots\ge\lambda_P>0$ and top eigenspace $\mathcal U_1=\ker(H-\lambda_1I)$ with orthogonal projector $P_1$. Let $H_{B_t}=\frac1B\sum_{b=1}^BH_{i_b}$ be sampled with replacement of batch size $B$ from arbitrary symmetric per-sample Hessians $H_i\in\mathbb S^P$ (without any feature or commutation assumption), write $\Delta H_i\triangleq H_i-H$, and define $\mathcal V(X)\triangleq\E[(H_{B_t}-H)X(H_{B_t}-H)]=\frac1B\E[\Delta H_iX\Delta H_i]\ne0$. Let $\mathcal L_{\eta,\beta}(X)=(HX+XH)-\frac{\eta}{1+\beta}HXH+\frac{\eta\beta}{(1-\beta)(1-\beta^2)}[H,[H,X]]$, which multiplies entry $(j,k)$ in the eigenbasis of $H$ by $\ell_{jk}(\eta,\beta)=(\lambda_j+\lambda_k)+\eta\,b_{jk}(\beta)$, positive for $\eta<\eta_{\mathrm{det}}$ (part (1)), where $b_{jk}(\beta)\triangleq-\frac{\lambda_j\lambda_k}{1+\beta}+\frac{\beta(\lambda_j-\lambda_k)^2}{(1-\beta)(1-\beta^2)}$, and for $\eta\in(0,\eta_{\mathrm{det}})$ with $\eta_{\mathrm{det}}\triangleq\frac{2(1+\beta)}{\lambda_1}$, set $\mathcal N_{\eta,\beta}\triangleq\eta\mathcal L_{\eta,\beta}^{-1}$ and $\mathcal K_{\eta,\beta}(X)\triangleq\frac1{1-\beta}\mathcal N_{\eta,\beta}(\mathcal V(X))=\frac{\eta}{1-\beta}\mathcal L_{\eta,\beta}^{-1}(\mathcal V(X))$ on $\mathbb S^P$.
(1) \emph{(Invertibility, cone preservation, real spectrum with the inertia of $\mathcal V$, and regular splitting.)} For every $\beta\in[0,1)$ and $\eta\in(0,\eta_{\mathrm{det}})$, $\mathcal L_{\eta,\beta}$ is invertible, $[1/\ell_{jk}(\eta,\beta)]_{j,k=1}^P\succeq0$ is positive semidefinite (singular when two eigenvalues of $H$ coincide), and $\mathcal L_{\eta,\beta}^{-1}$ and $\mathcal K_{\eta,\beta}$ preserve the cone $\mathbb S_+^P$. On $\mathcal H=(\mathbb S^P,\langle X,Y\rangle_F=\tr(XY))$, $\mathcal N_{\eta,\beta}\succ_{\mathcal H}0$ is self-adjoint and positive definite, whereas $\mathcal V$ is self-adjoint on $\mathcal H$ but in general indefinite (for $\Delta H=\operatorname{diag}(1,-1)$, $\mathcal V(E_{12}+E_{21})=-(E_{12}+E_{21})$); hence $\mathcal K_{\eta,\beta}$ is similar to the symmetric operator $\frac1{1-\beta}\mathcal N_{\eta,\beta}^{1/2}\mathcal V\mathcal N_{\eta,\beta}^{1/2}$, has a real spectrum with the same inertia as $\mathcal V$ (in general including negative eigenvalues), and by the Krein--Rutman theorem \citep{krein1948linear} its spectral radius $\rho(\mathcal K_{\eta,\beta})=\lambda_{\max}(\mathcal K_{\eta,\beta})$ is its largest eigenvalue, with a PSD eigenmatrix $X_\eta\in\mathbb S_+^P\setminus\{0\}$. Splitting $\mathcal T(\eta,\beta)=\mathcal T_0+\Delta$ on $\mathbb S^{2P}$, $\rho(\mathcal T(\eta,\beta))<1\iff\rho((I-\mathcal T_0)^{-1}\Delta)<1\iff\rho(\mathcal K_{\eta,\beta})<1$ and $\rho(\mathcal T(\eta,\beta))=1\iff\rho(\mathcal K_{\eta,\beta})=1$, because $\rho((I-\mathcal T_0)^{-1}\Delta)=\rho(\mathcal K_{\eta,\beta})$ by the factorization of Theorem~\ref{thm:splitting}(1) in Section~\ref{sec:burnin}.
(2) \emph{(Variational Rayleigh quotient and strict monotonicity for all $\beta\in[0,1)$.)} Since $\mathcal K_{\eta,\beta}$ is self-adjoint in the inner product weighted by $\mathcal N_{\eta,\beta}^{-1}$ and $\rho(\mathcal K_{\eta,\beta})=\lambda_{\max}(\mathcal K_{\eta,\beta})$,
\[
\rho(\mathcal K_{\eta,\beta})
=\frac{1}{1-\beta}\max_{X\in\mathbb S^P\setminus\{0\}}\frac{\langle X,\mathcal V(X)\rangle_F}{\sum_{j,k=1}^P\big(\ell_{jk}(\eta,\beta)/\eta\big)\widetilde X_{jk}^2}.
\] where $\widetilde X=U^\top XU$ is $X$ in the eigenbasis $U=[u_1,\dots,u_P]$ of $H$ and the sum runs over all $P^2$ entries.
The numerator does not depend on $\eta$ and its maximum is strictly positive because $X=I_P$ gives $\langle I_P,\mathcal V(I_P)\rangle_F=\frac1B\E\|\Delta H_i\|_F^2>0$. Each denominator weight $\ell_{jk}(\eta,\beta)/\eta=(\lambda_j+\lambda_k)/\eta+b_{jk}(\beta)$ is strictly decreasing in $\eta$ since $\lambda_j+\lambda_k>0$ and tends to $+\infty$ as $\eta\downarrow0$; hence $\eta\mapsto\rho(\mathcal K_{\eta,\beta})$ is continuous, strictly increasing on $(0,\eta_{\mathrm{det}})$ for every $\beta\in[0,1)$, and satisfies $\lim_{\eta\downarrow0}\rho(\mathcal K_{\eta,\beta})=0$ (at $\beta=0$, $\eta\mapsto\mathcal N_{\eta,0}$ even grows in the order of completely positive maps on $\mathbb S_+^P$, since with $\mu_j=\lambda_j^{-1}$ and $x_j=\mu_j-\eta/2>0$, $[\frac{\lambda_j+\lambda_k}{\ell_{jk}(\eta,0)^2}]$ is congruent via $\operatorname{diag}(\mu)$ to $[\frac1{x_j+x_k}+\frac\eta{(x_j+x_k)^2}]=\int_0^\infty(1+\eta s)e^{-s(x_j+x_k)}\,ds\succeq0$, positive definite when the $\lambda_j$ are distinct).
(3) \emph{(Necessary and sufficient blow-up, the two rates, and uniqueness of the exact ceiling.)} Near $\eta_{\mathrm{det}}$, only the top-block weights ($\lambda_j=\lambda_k=\lambda_1$) tend to zero, as $\ell_{11}(\eta,\beta)/\eta\asymp\eta_{\mathrm{det}}-\eta$, while all other weights stay bounded away from zero. By the Rayleigh quotient in (2): if $P_1\Delta H_iP_1\ne0$ for some sample $i$ (for simple $\lambda_1$, $\nu_{11,B}^2>0$), $X=P_1$ gives $\rho(\mathcal K_{\eta,\beta})=\Theta((\eta_{\mathrm{det}}-\eta)^{-1})$; if $P_1\Delta H_iP_1=0$ for all $i$ but $\Delta H_iP_1\ne0$ for some $i$ (pure cross-mode coupling, as in the Rademacher case of Corollary~\ref{cor:plateaus}), $X=sP_1+\mathcal V(P_1)$ ($\mathcal V(P_1)\ne0$) gives $\rho(\mathcal K_{\eta,\beta})=\Theta((\eta_{\mathrm{det}}-\eta)^{-1/2})$; if $\Delta H_iP_1=0$ for all $i$, $\mathcal V$ neither reads from nor writes to the top block, $\mathcal K_{\eta,\beta}$ reduces to the complement where all weights stay bounded away from zero, and $\rho(\mathcal K_{\eta,\beta})=O(1)$ remains bounded. Thus $\lim_{\eta\uparrow\eta_{\mathrm{det}}}\rho(\mathcal K_{\eta,\beta})=+\infty\iff\Delta H_iP_1\not\equiv0$, in which case there is a unique ceiling $\etacrit^{\mathrm{EoSS},P}(\beta,B)\in(0,\eta_{\mathrm{det}})$ with $\rho(\mathcal K_{\eta,\beta})=1$, its Krein--Rutman eigenmatrix $\Sigma_0^\star\in\mathbb S_+^P\setminus\{0\}$ is the second moment that loses stability first, and with $\mathcal L_0(X)=HX+XH$ and $\mathcal B_\beta(X)=\frac1{1+\beta}HXH-\frac{\beta}{(1-\beta)(1-\beta^2)}[H,[H,X]]+\frac1{1-\beta}\mathcal V(X)$, $\etacrit^{\mathrm{EoSS},P}(\beta,B)=1/\lambda_{\max}(\mathcal L_0^{-1/2}\mathcal B_\beta\mathcal L_0^{-1/2})$. In every case the mean-square ceiling is $1/\lambda_{\max}(\mathcal L_0^{-1/2}\mathcal B_\beta\mathcal L_0^{-1/2})\le\eta_{\mathrm{det}}$; when $\Delta H_iP_1\equiv0$ it can equal $\eta_{\mathrm{det}}$. Proof in Appendix~\ref{app:thmFG}.
\end{theorem}
At $\beta=0$, $\mathcal B_0(X)=HXH+\mathcal V(X)=(1-\frac1B)HXH+\frac1B\E[H_iXH_i]$ and $\mathcal L_0=2C$ with $C=\frac12H\oplus H$, so the last formula has the form $\eta^\star=2/\lambda_{\max}(C^{-1}D)$ of \citet[Theorem~5]{mulayoff2024exact}, which they prove at interpolating minima for sampling without replacement (with $C^\dagger$), here with $p=1/B$ for sampling with replacement and $H\succ0$; Theorem~\ref{thm:G6} extends it to heavy-ball momentum, where the double commutator and the factors $1/(1\pm\beta)$ enter $\mathcal B_\beta$. Against the full second-moment operator, in seven models with non-commuting per-sample Hessians ($P\in\{3,4,5\}$, $N\in\{4,5,6,8\}$ samples, $B\in\{1,2\}$, two of them with repeated eigenvalues of $H$) at $\beta\in\{0.5,0.9\}$, $\rho(\mathcal T)$ is below $1$ at $0.999\,\eta^\star$ and above $1$ at $1.001\,\eta^\star$ in all 14 cases, with $|\rho(\mathcal T(\eta^\star))-1|\le1.1\times10^{-13}$ (\texttt{paper/checks/diagonalization-burnin.py}). Theorems~\ref{thm:F} and~\ref{thm:G6} use the minibatch law only through $\mathcal V(X)=\E[(H_{B_t}-H)X(H_{B_t}-H)]$ and the independence of $(H_{B_t},\xi_t)$ across steps. For $B$ of $N$ fixed matrices $H_1,\dots,H_N$ with mean $H$ drawn without replacement within each step, independently across steps, the covariances of the inclusion indicators and $\sum_i\Delta H_i=0$ give $\mathcal V(X)=p\,\frac1N\sum_i\Delta H_iX\Delta H_i$ with $p=\frac{N-B}{B(N-1)}$, in place of $p=\frac1B$ for draws with replacement (\texttt{paper/checks/second-moment-renewal.py}); Theorem~\ref{thm:G6} is exact for both, and at $\beta=0$ it is Theorem~5 of \citet{mulayoff2024exact} with their $p$. In Theorem~\ref{thm:G} the factor $\frac1B$ enters only through $\mathcal V=p\mathcal V_1$ with $\mathcal V_1(X)=\E[(zz^\top-H)X(zz^\top-H)]$, so a law with $\mathcal V=p\mathcal V_1$ replaces $\frac1B$ by $p$ in \eqref{eq:secular} and in $\alpha_{\mathrm{diag}}=1+p\kappa_\beta(1+\kappa_{\mathrm{feat}})$. A pool of $n$ i.i.d.\ features drawn once and then sampled without replacement does not have that form: centred at its own mean, it gives $\E_{\mathrm{pool}}[(H_i-\hat H)X(H_i-\hat H)]=\frac{n-1}n\mathcal V_1(X)$ (with $\tilde H_j=H_j-H$ independent and centred, $H_i-\hat H=\frac{n-1}n\tilde H_i-\frac1n\sum_{j\ne i}\tilde H_j$, the cross terms vanish, and the coefficient is $\frac{(n-1)^2}{n^2}+\frac{n-1}{n^2}=\frac{n-1}n$), so its noise operator averaged over pool draws is $\frac{n-B}{B(n-1)}\cdot\frac{n-1}n\mathcal V_1=\frac{1-\zeta}B\mathcal V_1$, $\zeta=B/n$. The secular root with $p=(1-\zeta)/B$ uses that average and approximates the threshold of one realised pool, which Theorem~\ref{thm:G6} gives exactly. A pool that is the whole support of a law with independent coordinates, such as the $2^P$ sign vectors of Rademacher features, has $\hat H=H$ and $\mathcal V=\frac{n-B}{B(n-1)}\mathcal V_1$ exactly, so for it \eqref{eq:secular} with $p=\frac{n-B}{B(n-1)}$ is exact.

Two further cases reduce to the theorems. (a) Curvature outside $H\succ0$: if $\lambda_{\min}(H)<0$, the deterministic iteration has a root above $1$ on the mode with $\lambda<0$, and $\mathcal T\succeq\mathcal T_0$ gives $\rho(\mathcal T)\ge\rho(\bar F)^2>1$ at every $\eta>0$, so there is no mean-square stable rate. If $H\succeq0$ is singular and every $H_i\succeq0$, as at an interpolating minimum, then $\ker H=\bigcap_i\ker H_i$ and every $H_i$ maps into $\operatorname{range}H$; the component of $w_t$ in $\ker H$ follows $w_{t+1}=(1+\beta)w_t-\beta w_{t-1}$ plus the noise there, independently of the rest. With $\Sigma_\xi$ supported on $\operatorname{range}H$ it stays bounded, and the second moment is bounded if and only if $\eta$ is below the threshold of Theorem~\ref{thm:G6} on $\operatorname{range}H$, with $\mathcal L_0$ inverted there (at $\beta=0$ the pseudo-inverse $C^\dagger$ of \citet{mulayoff2024exact}); with noise in $\ker H$ it grows linearly at every $\eta$. Without $H_i\succeq0$, curvature noise on a direction of zero curvature drives the threshold to $0$, the limit $\lambda_k\to0$ of Theorem~\ref{thm:G5}(3). (b) A fixed preconditioner $P_{\mathrm{pre}}\succ0$, $w_{t+1}=w_t+\beta(w_t-w_{t-1})-\eta P_{\mathrm{pre}}(H_{B_t}w_t+\xi_t)$, need not commute with $H$ or the $H_i$: $\tilde w_t=P_{\mathrm{pre}}^{-1/2}w_t$ follows \eqref{eq:hbP} with $\tilde H_i=P_{\mathrm{pre}}^{1/2}H_iP_{\mathrm{pre}}^{1/2}$ and $\tilde\xi_t=P_{\mathrm{pre}}^{1/2}\xi_t$, so Theorems~\ref{thm:F} and~\ref{thm:G6} apply with $\tilde H=P_{\mathrm{pre}}^{1/2}HP_{\mathrm{pre}}^{1/2}$, whose eigenvalues are those of $P_{\mathrm{pre}}H$. For Gaussian features $z_i\sim\mathcal N(0,H)$, $P_{\mathrm{pre}}^{1/2}z_i\sim\mathcal N(0,\tilde H)$ again has independent coordinates in the eigenbasis of $\tilde H$, so Theorem~\ref{thm:G} holds with the eigenvalues of $P_{\mathrm{pre}}H$; for other feature laws that independence can be lost. (c) Random reshuffling, with independent permutations between epochs of $m=N/B$ disjoint batches, is not i.i.d.\ across steps, but with $\xi_t$ of mean zero, independent across steps and independent of the permutation and of the minibatch Hessians (an assumption that excludes noise drawn from the same batch as $H_{B_t}$), the second moment of $(w_{em},w_{em-1})$ at epoch boundaries obeys $Q_{e+1}=\mathcal T_{\mathrm{RR}}(Q_e)+G_{\mathrm{RR}}$, with $\mathcal T_{\mathrm{RR}}(Q)=\E_\pi[F_{\mathrm{ep}}QF_{\mathrm{ep}}^\top]$ for the epoch product $F_{\mathrm{ep}}=F_{e,m}\cdots F_{e,1}$, i.i.d.\ across epochs, and $\sum_{r=1}^m(H_{B_{e,r}}-H)=0$ in every epoch. So $\rho(\mathcal T_{\mathrm{RR}})<1$ bounds the second moment, and for $G_{\mathrm{RR}}\succ0$ the argument of Corollary~\ref{cor:msinstab} makes it also necessary; a closed form of $\rho(\mathcal T_{\mathrm{RR}})$ like \eqref{eq:secular} is not known. Within an epoch the batch errors $\Delta_r=H_{B_{e,r}}-H$ are exchangeable and sum to $0$, so $\E_\pi[\Delta_rX\Delta_s]=\frac{m\delta_{rs}-1}{m-1}\mathcal V(X)$ with $\mathcal V$ the without-replacement operator above; for $m=2$ and $\beta=0$, $\Delta_2=-\Delta_1$ and $F_{\mathrm{ep}}=F_2F_1=(I-\eta H)^2+\eta^2[H,\Delta_1]-\eta^2\Delta_1^2$ exactly. Along $u_1$ at $\beta=0$, with batch errors $\zeta_r$, $\sum_r\zeta_r=0$ gives $\sum_{s<r}\zeta_s\zeta_r=-\frac12\sum_r\zeta_r^2$, hence $\E_\pi[F_{\mathrm{ep}}^2]=a^{2m}-m\,a^{2m-2}\eta^2\nu_{11,B}^2$ with $a=1-\eta\lambda_1$, up to terms of higher order in the batch errors, against $a^{2m}+m\,a^{2m-2}\eta^2\nu_{11,B}^2$ to the same order for $m$ independent steps: the first-order curvature-noise correction changes sign, and the one-dimensional mean-square threshold of reshuffling is $\frac2{\lambda_1}(1+\nu_{11,B}^2/\lambda_1^2)$ to first order, the first-order almost-sure threshold of Proposition~\ref{prop:G3}(1), instead of $\frac2{\lambda_1}(1-\nu_{11,B}^2/\lambda_1^2)$. On a scalar pool of $8$ values with $B=2$ and $\nu_{11,B}^2=0.00424$, exact enumeration over the permutations gives $2.00854$ against $2.00849$, and $1.99155$ against $1.99151$ for independent steps (\texttt{paper/checks/second-moment-renewal.py}).

\begin{corollary}[The effective step at the threshold increases with $\beta$]\label{cor:etaeff}
Let $H\succ0$ and let the law of the per-sample Hessians, hence $\mathcal V$, not depend on $\beta$, and let $\eta^\star(\beta)$ be the threshold of Theorem~\ref{thm:G6}. Then $\eta^\star_{\mathrm{eff}}(\beta)=\eta^\star(\beta)/(1-\beta)$ increases strictly on $\beta\in[0,1)$, and
\[
\frac1{\eta^\star(0)}-\frac{1-\beta}{\eta^\star(\beta)}\ge\frac{\beta}{1+\beta}\lambda_{\min}(H).
\]
\end{corollary}
\begin{proof}
By Theorem~\ref{thm:G6}(3), $1/\eta^\star(\beta)=\lambda_{\max}(\mathcal L_0^{-1/2}\mathcal B_\beta\mathcal L_0^{-1/2})$ at every $\beta$, also when the threshold is $\eta_{\mathrm{det}}$, so $1/\eta^\star_{\mathrm{eff}}(\beta)=\lambda_{\max}(\mathcal L_0^{-1/2}\mathcal A_\beta\mathcal L_0^{-1/2})$ with $\mathcal A_\beta=(1-\beta)\mathcal B_\beta$. The term $\mathcal V$ cancels in $\mathcal D_\beta=\mathcal B_0-\mathcal A_\beta$, $\mathcal D_\beta(X)=\frac{2\beta}{1+\beta}HXH+\frac{\beta}{1-\beta^2}[H,[H,X]]$, which multiplies entry $(j,k)$ in the eigenbasis of $H$ by $m_{jk}(\beta)=\frac{2\beta}{1+\beta}\lambda_j\lambda_k+\frac{\beta}{1-\beta^2}(\lambda_j-\lambda_k)^2$. Both coefficients increase strictly on $[0,1)$, so $m_{jk}(\beta_2)>m_{jk}(\beta_1)$ for $\beta_1<\beta_2$ since $\lambda_j\lambda_k>0$. $\mathcal L_0$ multiplies the same entry by $\lambda_j+\lambda_k$, so $\mathcal L_0^{-1/2}(\mathcal A_{\beta_1}-\mathcal A_{\beta_2})\mathcal L_0^{-1/2}$ multiplies it by $(m_{jk}(\beta_2)-m_{jk}(\beta_1))/(\lambda_j+\lambda_k)>0$; it is positive definite on the finite-dimensional space $(\mathbb S^P,\langle\cdot,\cdot\rangle_F)$, and Weyl's inequality makes $\lambda_{\max}$ strictly smaller at $\beta_2$. For the bound, $m_{jk}(\beta)-\frac{\beta}{1+\beta}(\lambda_j^2+\lambda_k^2)=\frac{\beta^2(\lambda_j-\lambda_k)^2}{1-\beta^2}\ge0$ and $\lambda_j^2+\lambda_k^2\ge\lambda_{\min}(H)(\lambda_j+\lambda_k)$, so $\mathcal L_0^{-1/2}\mathcal D_\beta\mathcal L_0^{-1/2}\succeq\frac{\beta}{1+\beta}\lambda_{\min}(H)\,I$, and Weyl's inequality gives the bound.
\end{proof}
The small-step analysis of momentum SGD of \citet{yuan2016momentum} makes stability depend on $(\eta,\beta)$ only through $\eta_{\mathrm{eff}}=\eta/(1-\beta)$, and the leading-order condition of \citet[Theorem~4.1]{andreyev2026momentum}, in the noise-dominated regime and for small $\eta$, is the SGD condition at $\eta_{\mathrm{eff}}$. At the threshold this equivalence fails, as it already does without noise, where $\eta^\star_{\mathrm{eff}}=2(1+\beta)/((1-\beta)\lambda_1)$. Written in $\eta_{\mathrm{eff}}$, \eqref{eq:secular} reads $\frac{\eta_{\mathrm{eff}}p}2\sum_k\lambda_k/\big(1-\frac12\eta_{\mathrm{eff}}\lambda_k(\kappa_\beta^{-1}+p(1+\kappa_{\mathrm{feat}}))\big)=1$, where $\beta$ enters only through $\kappa_\beta^{-1}=\frac{1-\beta}{1+\beta}$, decreasing to $0$; at $\kappa_{\mathrm{feat}}=-1$ the limit as $\beta\to1$ is $\eta^\star_{\mathrm{eff}}=2/(p\tr H)$. The secular roots of the Gaussian cells of Phase~1 at $P=64$ and $B=64$ give $\eta^\star_{\mathrm{eff}}=1.927$ at $\beta=0$ and $7.948$ at $\beta=0.9$, a factor $4.12$; \texttt{paper/checks/second-moment-renewal.py} computes these roots and, separately, checks the monotonicity and the bound on a non-commuting model with $P=3$, $N=5$ and $B=2$.

For the loss of mini-batch SGD with momentum on least squares, for large $n$ and $d$ under orthogonal invariance of the data, \citet{lee2022minibatch} write a Volterra equation for the loss whose kernel norm $\|\mathcal K\|<1$ they state, citing \citet[Proposition~7.4]{asmussen2003}, to be necessary for its solution to converge; their Theorem~1, the concentration of SGD with momentum on that solution over a fixed horizon, assumes $\gamma<(1+\Delta)/(\zeta\sigma_{\max}^2)$, which is $\eta\lambda_{\max}<1+\beta$ in the dictionary below, and in the proof of their Proposition~2 they put the norm in closed form, $\|\mathcal K\|=\frac1n\sum_j\frac{(1-\zeta)\gamma\sigma_j^2(1+\Delta)}{(1-\Delta)(1+\Delta+\Omega_j)}$ with $\Omega_j=1-\gamma\zeta\sigma_j^2+\Delta$. Their update sums the $B$ sample gradients and \eqref{eq:hbP} averages them (in their notation $\beta$ is the batch size and $\Delta$ the momentum); with $\eta=\gamma B$, $\beta=\Delta$, $\lambda_j=\sigma_j^2/n$ and $\zeta=B/n$, this is the left side of \eqref{eq:secular} with $\alpha_{\mathrm{diag}}=1$ and a factor $1-\zeta$; in the model of this note with $B=\zeta n$ the difference $\alpha_{\mathrm{diag}}-1=\kappa_\beta(1+\kappa_{\mathrm{feat}})/B$ is $O(1/n)$; their data are orthogonally invariant and the note's have independent coordinates, two classes whose common members are Gaussian, and their statement holds with high probability as $n,d\to\infty$ while this one is exact at finite $P$; what coincides is the two formulas. Their Proposition~2 gives two conditions that together imply $\|\mathcal K\|<1$, $\eta\lambda_{\max}<1+\beta$ (half the deterministic threshold) and a trace bound, which their convergence-rate analysis uses; locating the root of $\|\mathcal K\|=1$ is not its aim. In 6 of the 30 cells of Phase~1 ($P=64$, $\beta=0$, $B\in\{16,64\}$, two of them in the pool arm scored by \texttt{p4}) that root lies above $(1+\beta)/\lambda_1$, outside the hypothesis of their Theorem~1; there the identity is with their formula, not with their theorem on SGD with momentum. The two objects differ: their kernel norm governs the loss in the high-dimensional limit, \eqref{eq:secular} the second moment of the parameters at finite $P$; and the factor $1-\zeta$ is the pool average of the noise operator of sampling without replacement (paragraph after Theorem~\ref{thm:G6}), so for one realised pool \eqref{eq:secular} with it is an approximation, which the pool arm of Phase~1 uses as its prediction. The same secular equation is the loss-convergence condition $\widetilde U(1)<1$ of \citet[eq.~(21) and the paragraph after Proposition~3]{velikanov2023view} for mini-batch SGD with momentum under their spectrally expressible approximation, with their noise strength $\gamma=(N-b)/((N-1)b)$ of a finite dataset sampled without replacement replaced by $1/B$, their value for an infinite dataset and the factor with replacement here, and parameter $\tau=-(1+\kappa_{\mathrm{feat}})$. Their Section~5, where Proposition~3 sits, assumes $0<\tau\le1$ and an infinite-dimensional compact $H$ ($\lambda_k\to0$); $\tau$ lies in that range for $\kappa_{\mathrm{feat}}\in[-2,-1)$ ($\tau=1$ for Rademacher features) and outside it for $\kappa_{\mathrm{feat}}\ge-1$, including the Gaussian case $\tau=-1$ they attribute to Bordelon and Pehlevan, and the second assumption fails at finite $P$, so at finite $P$ the correspondence is formal; their $\lambda_{\mathrm{crit}}$ bound (Proposition~4) is a one-sided counterpart of the sandwich \eqref{eq:sandwich}. Theorem~\ref{thm:G} shows that the approximation is exact at finite $P$ and $B$ for features with independent coordinates of any kurtosis $\kappa_{\mathrm{feat}}\ge-2$, proves the ``if and only if'' on the full second-moment operator, and adds the reading of the root as the stochastic ceiling, the two plateaus and the Rademacher separation; Theorem~\ref{thm:G5}, with $\Sigma_z\ne H$, lies outside the spectrally expressible class. At $\beta=0$ the sandwich gives the batch-size scaling $B^\star\asymp\tr H/\lambda_1$ of \citet{ma2018power}; the component-decoupled data model of \citet{kidambi2018insufficiency} and \citet{liu2020mass} has $H_{B_t}$ commuting with $H$ and no cross-mode term, and \citet{paquette2021momentum}, at batch size $1$, describe by a Volterra equation (their Theorem~1) the loss of homogenized SGD, a diffusion approximation of a class of stochastic momentum methods that contains heavy-ball with fixed momentum after a rescaling of its parameters with the dimension, in a scaling of small steps and high dimension in which the deterministic ceiling $2(1+\beta)/\lambda_1$ does not enter; in that limit the homogenized dynamics of heavy-ball with fixed momentum and of SGD have the same loss when the SGD step is their $\gamma^{\mathrm{shb}}/\theta^{\mathrm{shb}}$ (their Theorem~5). Phase~1 of the lab measures Theorem~\ref{thm:G} (Section~\ref{sec:phase1}).
\begin{proposition}[A renewal-type equation for the expected loss at finite $P$]\label{prop:renewal}
In the setting of Theorem~\ref{thm:G}, with $\mathcal V=p\mathcal V_1$, $\kappa_{\mathrm{feat}}\ge-2$, $\Sigma_\xi$ diagonal in the eigenbasis of $H$, $w_0=w_{-1}=0$ and $\eta<\eta_{\mathrm{det}}$, let $g_k(0)=1$, $g_k(1)=1+\beta-\eta\lambda_k$ and $g_k(\tau+1)=(1+\beta-\eta\lambda_k)g_k(\tau)-\beta g_k(\tau-1)$ be the impulse response of mode $k$ of deterministic heavy-ball, $\psi_k(\tau)=\eta^2g_k(\tau)^2$, and $h_k$ the solution of the scalar convolution equation $h_k(\tau)=\psi_k(\tau)+p(1+\kappa_{\mathrm{feat}})\lambda_k^2\sum_{u<\tau}\psi_k(\tau-1-u)h_k(u)$. Then $\ell_t=\tr(H\E[w_tw_t^\top])$, twice the expected quadratic loss, solves exactly
\[
\ell_t=f_t+\sum_{s=0}^{t-1}K_{t-1-s}\,\ell_s,\qquad K_\tau=p\sum_{k=1}^P\lambda_k^2h_k(\tau),\qquad f_t=\sum_{k=1}^P\lambda_k\sigma_{\xi,k}^2\sum_{\tau<t}h_k(\tau).
\]
If $\kappa_{\mathrm{feat}}\ge-1$, then $h_k\ge0$; if moreover $\eta\lambda_k\alpha_{\mathrm{diag}}<2(1+\beta)$ for every $k$, then $\sum_{\tau\ge0}K_\tau=\frac{\eta p}{2(1-\beta)}\sum_k\frac{\lambda_k}{1-\eta\lambda_k\alpha_{\mathrm{diag}}/(2(1+\beta))}$, the left side of \eqref{eq:secular} with $\frac1B$ replaced by $p$ (as everywhere in this proposition), and $\ell_t$ is bounded if and only if this sum is below $1$. In this model with $B=\zeta n$ and $n\to\infty$, the regime of \citet{lee2022minibatch}, $\alpha_{\mathrm{diag}}-1=O(1/n)$ vanishes whatever the kurtosis, and with $p=(1-\zeta)/B$ and the dictionary of the previous paragraph their kernel norm is the left side of \eqref{eq:secular} at $\alpha_{\mathrm{diag}}=1$; at finite $P$, $\alpha_{\mathrm{diag}}=1$ holds exactly at $\kappa_{\mathrm{feat}}=-1$, where $h_k=\psi_k$ and the loss obeys a renewal equation of their form, and for $\kappa_{\mathrm{feat}}>-1$ the same with $h_k$ in place of $\psi_k$. For $\kappa_{\mathrm{feat}}\in[-2,-1)$ the equation is a discrete Volterra equation with a kernel of either sign, not a renewal equation in the probabilistic sense, but $\lambda_P\E\|w_t\|^2\le\ell_t\le\lambda_1\E\|w_t\|^2$, so Theorem~\ref{thm:G} and Corollary~\ref{cor:msinstab} give the same conclusion for every $\kappa_{\mathrm{feat}}\ge-2$: with every $\sigma_{\xi,k}^2>0$ and $\eta<2(1+\beta)/(\lambda_1\max(\alpha_{\mathrm{diag}},1))$, $\ell_t$ is bounded if and only if the left side of \eqref{eq:secular} is below $1$; and the Neumann series $h_k=\sum_{n\ge0}c_k^n\psi_k^{*(n+1)}$ converges absolutely, with $\sum_\tau h_k(\tau)=\Psi_k/(1-c_k\Psi_k)$, when $|c_k|\Psi_k<1$, that is $\eta\lambda_k(2-\alpha_{\mathrm{diag}})<2(1+\beta)$, in which case $\sum_\tau K_\tau$ is again the left side of \eqref{eq:secular}. That bound is sufficient, not necessary, for $h_k$ to be summable; where $h_k$ is not summable the left side can still be below $1$ (for $P=1$, $\beta=0$, $B=1$, $\kappa_{\mathrm{feat}}=-2$, $\lambda_1=1$ and $\eta=1.5$ it is $0.75$ while $h_1$ grows geometrically), and the boundedness of $\ell_t$ comes from Theorem~\ref{thm:G}, not from a kernel sum.
\end{proposition}
\begin{proof}
With $s_t=(w_t,w_{t-1})$ and $Q_t=\E[s_ts_t^\top]$, $Q_{t+1}=\mathcal T_0(Q_t)+\Delta(Q_t)+G$ (proof of Corollary~\ref{cor:msinstab}), so from $Q_0=0$, $Q_t=\sum_{s<t}\mathcal T_0^{t-1-s}(\Delta(Q_s)+G)$. The $(1,1)$ block of $\mathcal T_0^\tau(\operatorname{blkdiag}(Y,0))$ is $(\bar F^\tau)_{11}Y(\bar F^\tau)_{11}^\top$, with $(\bar F^\tau)_{11}=\operatorname{diag}(g_k(\tau))$ in the eigenbasis of $H$. By Theorem~\ref{thm:G}, $\mathcal V(X)_{kk}=p\lambda_k\tr(HX)+p(1+\kappa_{\mathrm{feat}})\lambda_k^2X_{kk}$, so $x_t=(\Sigma_{0,t})_{kk}$ solves $x_t=\sum_{s<t}\psi_k(t-1-s)\big(\sigma_{\xi,k}^2+p\lambda_k\ell_s+p(1+\kappa_{\mathrm{feat}})\lambda_k^2x_s\big)$, whose solution is $x_t=\sum_{s<t}h_k(t-1-s)(\sigma_{\xi,k}^2+p\lambda_k\ell_s)$; multiplying by $\lambda_k$ and summing over $k$ gives the equation. For $\kappa_{\mathrm{feat}}\ge-1$ the coefficient $c_k=p(1+\kappa_{\mathrm{feat}})\lambda_k^2$ is non-negative, so $h_k\ge0$, and $\sum_\tau h_k(\tau)=\Psi_k/(1-c_k\Psi_k)$ with $\Psi_k=\sum_\tau\psi_k(\tau)=\frac{\eta}{2(1-\beta)\lambda_k(1-\eta\lambda_k/(2(1+\beta)))}$ (Yule--Walker equations of the second-order autoregression), finite when $c_k\Psi_k<1$, which is $\eta\lambda_k\alpha_{\mathrm{diag}}<2(1+\beta)$; then $p\lambda_k^2\sum_\tau h_k(\tau)=\frac{\eta p}{2(1-\beta)}\frac{\lambda_k}{1-\eta\lambda_k\alpha_{\mathrm{diag}}/(2(1+\beta))}$. A renewal equation with a non-negative kernel and a forcing that increases to a positive limit has a bounded solution if and only if the kernel sums to less than $1$ \citep[Chapter~V]{asmussen2003}. The identity is checked against the exact second-moment recursion for $\kappa_{\mathrm{feat}}\in\{-2,-1,0,2\}$ in \texttt{paper/checks/second-moment-renewal.py}.
\end{proof}
\begin{remark}[When the propagator bound holds up to the threshold]\label{rem:renewal-rademacher}
For $\kappa_{\mathrm{feat}}\in[-2,-1)$, $|c_k|\Psi_k<1$ holds for every $k$ at $\eta=\eta^\star$ if and only if $\mu^\star>(2-\alpha_{\mathrm{diag}})\lambda_1$. Indeed $|c_k|\Psi_k<1$ is $\eta\lambda_k(2-\alpha_{\mathrm{diag}})<2(1+\beta)$, so the mode $k=1$ binds; with $f(\mu)=p\kappa_\beta\sum_k\lambda_k/(\mu-\alpha_{\mathrm{diag}}\lambda_k)$, which equals $1$ at $\mu^\star$ and decreases strictly for $\mu>\max(\alpha_{\mathrm{diag}},0)\lambda_1$, the condition is $f((2-\alpha_{\mathrm{diag}})\lambda_1)>1$, and since $1-\alpha_{\mathrm{diag}}=p\kappa_\beta|1+\kappa_{\mathrm{feat}}|$ the term $k=1$ of $f$ there is $1/(2|1+\kappa_{\mathrm{feat}}|)$, which gives $\sum_{k\ge2}\frac{p\kappa_\beta\lambda_k}{(2-\alpha_{\mathrm{diag}})\lambda_1-\alpha_{\mathrm{diag}}\lambda_k}>1-\frac1{2|1+\kappa_{\mathrm{feat}}|}$. The lower bound of \eqref{eq:sandwich}, with $\frac1B$ replaced by $p$, gives it whenever $\min(\alpha_{\mathrm{diag}},0)\lambda_1+p\kappa_\beta\tr H>(2-\alpha_{\mathrm{diag}})\lambda_1$. On the ten Rademacher cells of Phase~1 it holds in 9 and fails at $P=64$, $\beta=0$, $B=64$ (\texttt{paper/checks/second-moment-renewal.py}).
\end{remark}
\begin{theorem}[Modal interpolation identity, fixed-point divergence criteria, and Batch Sharpness at the threshold]\label{thm:modal-bridge}
In the setting of Theorem~\ref{thm:G6}, let $\eta^\star=\etacrit^{\mathrm{EoSS},P}(\beta,B)=1/\lambda_{\max}\big(\mathcal L_0^{-1/2}\mathcal B_\beta\mathcal L_0^{-1/2}\big)\in(0,\eta_{\mathrm{det}}]$ and let $\Sigma_0^\star\in\mathbb S_+^P\setminus\{0\}$ ($\tr\Sigma_0^\star=1$) satisfy $(\mathcal L_0-\eta^\star\mathcal B_\beta)(\Sigma_0^\star)=0$ (for $\eta^\star<\eta_{\mathrm{det}}$, equivalently $\mathcal K_{\eta^\star,\beta}(\Sigma_0^\star)=\Sigma_0^\star$).
\begin{enumerate}
\item \emph{(Modal interpolation law and one-sided Batch Sharpness instability criteria at a fixed checkpoint.)} With the modal curvature and curvature-noise ratios on $\Sigma_0^\star$,
\begin{equation}\label{eq:modal-av}
a_\star\triangleq\frac{\tr(H^2\Sigma_0^\star)}{\tr(H\Sigma_0^\star)}\in[\lambda_P,\lambda_1],\qquad
v_\star\triangleq\frac{\tr(\mathcal V(\Sigma_0^\star))}{\tr(H\Sigma_0^\star)}=\frac{\tr\big((\E[H_{B_t}^2]-H^2)\Sigma_0^\star\big)}{\tr(H\Sigma_0^\star)}\ge0,
\end{equation}
and the modal second-moment curvature $\mathcal C_{\mathrm{mod}}(\Sigma_0^\star)\triangleq\frac{\tr(\E[H_{B_t}^2]\Sigma_0^\star)}{\tr(H\Sigma_0^\star)}=a_\star+v_\star$, one has (extending Proposition~\ref{prop:G3}(2) and the 1D formula of \citet[eq.~20]{andreyev2026momentum} to $\Sigma_0^\star\in\mathbb S_+^P$)
\begin{equation}\label{eq:exact-P-interp}
\frac{1}{\eta^\star}=\frac{a_\star}{2(1+\beta)}+\frac{v_\star}{2(1-\beta)},\qquad
\mathcal C_{\mathrm{mod}}(\Sigma_0^\star)=\frac{2(1-\beta)}{\eta^\star}+\frac{2\beta}{1+\beta}a_\star=\frac{2(1+\beta)}{\eta^\star}-\frac{2\beta}{1-\beta}v_\star\in\Big[\frac{2(1-\beta)}{\eta^\star},\frac{2(1+\beta)}{\eta^\star}\Big].
\end{equation}
For every $S\in\mathbb S_+^P\setminus\{0\}$ commuting with $H$ ($[H,S]=0$), setting $a(S)\triangleq\frac{\tr(H^2S)}{\tr(HS)}$ and $v(S)\triangleq\frac{\tr(S^{1/2}\mathcal V(S^{1/2}))}{\tr(HS)}$ ($\bar a\triangleq a(I_P)=\frac{\tr(H^2)}{\tr H}$, $\bar v\triangleq v(I_P)=\frac{\E\|\Delta H_i\|_F^2}{B\tr H}$, $v(u_1u_1^\top)=\frac{\nu_{11,B}^2}{\lambda_1}$) gives $\frac{1}{\eta^\star}\ge\frac{a(S)}{2(1+\beta)}+\frac{v(S)}{2(1-\beta)}$, so $\frac{a(S)}{2(1+\beta)}+\frac{v(S)}{2(1-\beta)}\ge\frac1\eta$ is a valid sufficient instability criterion in the sense of \citet[Definition~A.1]{andreyev2026momentum} under the hypotheses of Corollary~\ref{cor:msinstab} ($\Sigma_\xi\succ0$ and either $\xi_t$ independent of $H_{B_t}$ or $\E[w_0]=\E[w_{-1}]=0$). In particular:
(a) for any checkpoint $\theta_0$ and any $H_B(\theta_0)\in\mathbb S^P$ with $\Pr(g_B(\theta_0)\ne0)=1$, setting $R_B(\theta_0)\triangleq\frac{g_B(\theta_0)^\top H_B(\theta_0)g_B(\theta_0)}{\|g_B(\theta_0)\|_2^2}$ and $W_B(\theta_0)\triangleq\frac{\|g_B(\theta_0)\|_2^2}{\E_B\|g_B(\theta_0)\|_2^2}$, the ratio-of-expectations $\overline{\mathrm{BS}}(\theta_0)\triangleq\frac{\E_B[g_B(\theta_0)^\top H_B(\theta_0)g_B(\theta_0)]}{\E_B\|g_B(\theta_0)\|_2^2}$ and the pointwise expectation-of-the-ratio $\mathrm{BS}(\theta_0)\triangleq\E_B[R_B(\theta_0)]$ of \citet[Definition~3.1]{andreyev2026momentum} satisfy the exact covariance identity $\overline{\mathrm{BS}}(\theta_0)-\mathrm{BS}(\theta_0)=\operatorname{Cov}_B(R_B(\theta_0),W_B(\theta_0))$, so $|\overline{\mathrm{BS}}(\theta_0)-\mathrm{BS}(\theta_0)|\le\sigma_{\mathrm{BS}}(\theta_0)\,\mathrm{CV}_{g^2}(\theta_0)$ with $\sigma_{\mathrm{BS}}(\theta_0)\triangleq(\Var_B(R_B(\theta_0)))^{1/2}$ and $\mathrm{CV}_{g^2}(\theta_0)\triangleq\frac{(\Var_B(\|g_B(\theta_0)\|_2^2))^{1/2}}{\E_B\|g_B(\theta_0)\|_2^2}$; when $\E[g_B(\theta_0)g_B(\theta_0)^\top\mid H_B]=\sigma_g^2H_B(\theta_0)$ ($\sigma_g>0$), $\overline{\mathrm{BS}}(\theta_0)=\bar a+\bar v$, and with noise share $\theta_{\mathrm{noise}}\triangleq\frac{\bar v}{\bar a+\bar v}\in[0,1]$,
\begin{equation}\label{eq:bs-fixed-instab}
\overline{\mathrm{BS}}(\theta_0)\ge\mathrm{BS}(\theta_0)-\sigma_{\mathrm{BS}}(\theta_0)\,\mathrm{CV}_{g^2}(\theta_0)\ge\frac{2(1+\beta)}{\eta\,\big(1+\frac{2\beta}{1-\beta}\theta_{\mathrm{noise}}\big)}\quad\Big(\text{in particular }\ge\frac{2(1+\beta)}{\eta}\text{ without }\theta_{\mathrm{noise}}\text{, or }\bar v\ge\frac{2(1-\beta)}{\eta}\Big)\implies\eta\ge\eta^\star\implies\rho(\mathcal T)\ge1,
\end{equation}
whose threshold is $\frac{2}{\eta}$ at $\beta=0$, $\frac{2(1+\beta)}{\eta}$ at $\theta_{\mathrm{noise}}=0$, and $\frac{2(1-\beta)}{\eta}$ at $\theta_{\mathrm{noise}}=1$; and
(b) in the random-feature model of Theorem~\ref{thm:G} ($H_B=\frac1BZZ^\top$, $Z=[z_{i_1},\dots,z_{i_B}]$), $v_\star=\frac{\tr H}{B}+\frac{1+\kappa_{\mathrm{feat}}}{B}a_\star$, $\frac{2(1+\beta)}{\eta^\star}=\alpha_{\mathrm{diag}}a_\star+\kappa_\beta\frac{\tr H}{B}$, and for every batch $B$ and every $g_B=\frac1BZr\in\operatorname{range}(Z)\setminus\{0\}$ with arbitrary $r\in\R^B$ and $u\triangleq r/\|r\|_2\in\mathbb S^{B-1}$ (requiring no independence between $r$ and $Z$, and holding in particular for $g_B(\theta)=\frac1B\sum_{b=1}^Br_{i_b}(\theta)z_{i_b}$ at any fixed $\theta$ on the rank-one Gauss--Newton surrogate $H_i=z_iz_i^\top$), the normalised Gram matrix $G_Z\triangleq(\tr H)^{-1}Z^\top Z\in\mathbb S_+^B$ satisfies the exact pathwise identity $\frac{g_B^\top H_Bg_B}{\|g_B\|_2^2}-\frac{\tr H}{B}=\frac{\tr H}{B}\big(u^\top(G_Z-I_B)u+\frac{\|(I_B-uu^\top)(G_Z-I_B)u\|_2^2}{u^\top G_Zu}\big)$ and $\big|\frac{g_B^\top H_Bg_B}{\|g_B\|_2^2}-\frac{\tr H}{B}\big|\le\frac{\tr H}{B}\|G_Z-I_B\|_{\mathrm{op}}\le\frac{\tr H}{B}\|G_Z-I_B\|_F$, with $\E\|(I_B-uu^\top)(G_Z-I_B)u\|_2^2=\big(B-1+\kappa_{\mathrm{feat}}(1-\sum_{b=1}^Bu_b^4)\big)\frac{\tr(H^2)}{(\tr H)^2}$ for every deterministic $u\in\mathbb S^{B-1}$ and $\E\|G_Z-I_B\|_F^2=B(B+1+\kappa_{\mathrm{feat}})\frac{\tr(H^2)}{(\tr H)^2}\triangleq\delta_{\mathrm{Gram}}^2\le B(B+1+\kappa_{\mathrm{feat}})\frac{\lambda_1}{\tr H}$; hence the pointwise expectation-of-the-ratio Batch Sharpness $\mathrm{BS}(\theta)=\E_B[\frac{g_B(\theta)^\top H_Bg_B(\theta)}{\|g_B(\theta)\|_2^2}]$ (with $\Pr(g_B(\theta)\ne0)=1$) satisfies $\mathrm{BS}(\theta)=\frac{\tr H}{B}\big(1+\E_B[\frac{\|(I_B-uu^\top)(G_Z-I_B)u\|_2^2}{u^\top G_Zu}]\big)\ge\frac{\tr H}{B}$ for every $B\ge1$ whenever $u(\theta)$ is independent of $G_Z$ (with equality at $B=1$ for all $\kappa_{\mathrm{feat}}\ge-2$), and for arbitrary $r(\theta)$,
\begin{equation}\label{eq:bs-pointwise-gram}
\Big|\mathrm{BS}(\theta)-\frac{\tr H}{B}\Big|\le\frac{\tr H}{B}\,\delta_{\mathrm{Gram}}\quad\forall\,\theta\qquad(\delta_{\mathrm{Gram}}=0\text{ at }B=1,\ \kappa_{\mathrm{feat}}=-2),
\end{equation}
so in the small-batch regime $\frac{B^2\lambda_1}{\tr H}\to0$, $\frac{\mathrm{BS}(\theta)}{2(1-\beta)/\eta^\star}=1+O(\delta_{\mathrm{Gram}}+\frac{B\lambda_1}{\tr H})\to1$ uniformly in $\theta$, and for $\kappa_{\mathrm{feat}}\ge-1$, $\mathrm{BS}(\theta)\ge\frac{2(1-\beta)}{\eta}(1+\delta_{\mathrm{Gram}})\implies\frac{\tr H}{B}\ge\frac{2(1-\beta)}{\eta}\implies\eta\ge\eta^\star\implies\rho(\mathcal T)\ge1$ is a valid sufficient instability criterion in the sense of \citet[Definition~A.1]{andreyev2026momentum} stated directly for $\mathrm{BS}(\theta)$.
\item \emph{(Global Krein--Rutman pole, stationary Batch Sharpness limit as $\eta\uparrow\eta^\star$, and the two SGDM plateaus.)} Equip $\mathbb S^P$ with the $\eta$-independent inner product $\langle X,Y\rangle_{\mathcal L_0}\triangleq\langle X,\mathcal L_0(Y)\rangle_F=2\tr(XHY)$, in which $\mathcal L_0^{-1}\mathcal B_\beta$ is self-adjoint with top eigenvalue $1/\eta^\star$. Assume $\Sigma_\xi\succ0$ and that $1/\eta^\star$ is a simple eigenvalue of $\mathcal L_0^{-1}\mathcal B_\beta$ with spectral gap $\Delta_{\mathcal G}\triangleq 1/\eta^\star-\lambda_2(\mathcal L_0^{-1}\mathcal B_\beta)>0$. Then for every $\eta\in(0,\eta^\star)$, the stationary covariance $\Sigma_0(\eta)$ of Theorem~\ref{thm:F}(2) decomposes orthogonally in $(\mathbb S^P,\langle\cdot,\cdot\rangle_{\mathcal L_0})$ as
\begin{equation}\label{eq:sigma0-pole}
\Sigma_0(\eta)=\frac{\eta\,\tr(\Sigma_0^\star\Sigma_\xi)}{2(1-\beta)\,\tr\big(H(\Sigma_0^\star)^2\big)\,\big(1-\eta/\eta^\star\big)}\,\Sigma_0^\star+R_\perp(\eta),\qquad
\|R_\perp(\eta)\|_{\mathcal L_0}\le\frac{\|\mathcal L_0^{-1/2}(\Sigma_\xi)\|_F}{(1-\beta)\,\Delta_{\mathcal G}}.
\end{equation}
When $\eta^\star<\eta_{\mathrm{det}}$, $1-\rho(\mathcal K_{\eta,\beta})=d_\star(1-\eta/\eta^\star)+O((\eta^\star-\eta)^2)$ with $d_\star\triangleq\eta^\star\partial_\eta\rho(\mathcal K_{\eta,\beta})\big|_{\eta=\eta^\star}=\frac{2(1-\beta)\tr(H(\Sigma_0^\star)^2)}{\eta^\star\langle\Sigma_0^\star,\mathcal V(\Sigma_0^\star)\rangle_F}>0$. Consequently $\Sigma_0(\eta)/\tr\Sigma_0(\eta)\to\Sigma_0^\star$ as $\eta\uparrow\eta^\star$ for every $\Sigma_\xi\succ0$, and on the linearised gradient $g_{B_t}(w_t)=H_{B_t}w_t+\xi_t$ with $(H_{B_t},\xi_t)$ independent of $w_t$, $\E\|H_{B_t}\|_{\mathrm{op}}^3<\infty$ and $\E[\|H_{B_t}\|_{\mathrm{op}}\|\xi_t\|_2^2]<\infty$, the stationary ratio-of-expectations Batch Sharpness satisfies
\begin{equation}\label{eq:bs-star-limit}
\lim_{\eta\uparrow\eta^\star}\overline{\mathrm{BS}}(\eta)\triangleq\lim_{\eta\uparrow\eta^\star}\frac{\E\big[g_{B_t}(w_t)^\top H_{B_t}g_{B_t}(w_t)\big]}{\E\big[\|g_{B_t}(w_t)\|_2^2\big]}=\mathrm{BS}^\star(\Sigma_0^\star)\triangleq\frac{\tr\big(\E[H_{B_t}^3]\,\Sigma_0^\star\big)}{\tr\big(\E[H_{B_t}^2]\,\Sigma_0^\star\big)}.
\end{equation}
Whenever $H_{B_t}\succeq0$ almost surely (in particular when $H_i\succeq0$ for all $i$), $\mathrm{BS}^\star(\Sigma_0^\star)\ge\mathcal C_{\mathrm{mod}}(\Sigma_0^\star)\ge\frac{2(1-\beta)}{\eta^\star}$ (and pointwise $\frac{w^\top H_{B_t}^3w}{w^\top H_{B_t}^2w}\ge\frac{w^\top H_{B_t}^2w}{w^\top H_{B_t}w}$ on $\{H_{B_t}w\ne0\}$), and if $H_{B_t}\preceq M_{\max}I_P$ almost surely, $\mathrm{BS}^\star(\Sigma_0^\star)\le M_{\max}$. In the random-feature model of Theorem~\ref{thm:G} with $\mu_6\triangleq\max_k\E[\omega_{i,k}^6]<\infty$, let $q_{\mathrm{feat}}\triangleq 1+\frac{1+\kappa_{\mathrm{feat}}}{B}$, $r_{\mathrm{feat}}\triangleq 1+\frac{3(1+\kappa_{\mathrm{feat}})}{B}+\frac{\E[\omega_{i,k}^6]-11-6\kappa_{\mathrm{feat}}}{B^2}$, and $m_{2,\star}\triangleq\frac{\tr(H^3\Sigma_0^\star)}{\tr(H\Sigma_0^\star)}$. Then $\alpha_{\mathrm{diag}}(m_{2,\star}-a_\star^2)=\frac{\kappa_\beta}{B}\big((\tr H)a_\star-\tr(H^2)\big)$ and, for identically distributed $\omega_{i,k}$,
\begin{equation}\label{eq:bs-rf-exact}
\mathrm{BS}^\star(\Sigma_0^\star)=\mathcal C_{\mathrm{mod}}(\Sigma_0^\star)+\frac{\frac{q_{\mathrm{feat}}}{B}\tr(H^2)+q_{\mathrm{feat}}^2(m_{2,\star}-a_\star^2)+(r_{\mathrm{feat}}-q_{\mathrm{feat}}^2)m_{2,\star}}{\frac{\tr H}{B}+q_{\mathrm{feat}}a_\star},
\end{equation}
so both $\mathcal C_{\mathrm{mod}}(\Sigma_0^\star)$ and $\mathrm{BS}^\star(\Sigma_0^\star)$ converge to the two SGDM plateaus of \citet{andreyev2026momentum}: $\frac{\mathrm{BS}^\star(\Sigma_0^\star)}{2(1-\beta)/\eta^\star}=1+O(\frac{B\lambda_1}{\tr H})\to1$ as $\frac{B\lambda_1}{\tr H}\to0$, and $\frac{\mathrm{BS}^\star(\Sigma_0^\star)}{2(1+\beta)/\eta^\star}=1+O(\frac{\kappa_\beta\tr H}{B\lambda_1})\to1$ as $\frac{\kappa_\beta\tr H}{B\lambda_1}\to0$.
\item \emph{(Suppression of $\mathrm{BS}^\star$ relative to $2(1-\beta)/\eta^\star$ under flat-mode curvature noise when $H_{B_t}\not\succeq0$.)} In the setting of Theorem~\ref{thm:G5} ($P\ge2$, where $H_{B_t}$ can be indefinite) with $\max_m(\lambda_m+s_m)\le M_0<\infty$ and $\widetilde\kappa_6\triangleq\max_m\frac{|\E[(\omega_{i,m}^2-1)^3]|}{2+\kappa_{\mathrm{feat}}}<\infty$ ($0/0\triangleq0$: $\widetilde\kappa_6=0$ for Rademacher, $4$ for Gaussian), $|\mathrm{BS}^\star(\Sigma_0^\star)|\le\big(3+\frac{\max(P+\kappa_{\mathrm{feat}},\,2(P-1)+\widetilde\kappa_6)}{B}\big)M_0$ for every spectrum, whereas if a mode $k$ has $\lambda_k\downarrow0$ at fixed $s_k>0$ while some $j\ne k$ has $s_j>0$, then $\eta^\star=O(\sqrt{\lambda_k})$ by Theorem~\ref{thm:G5}(3) and
\begin{equation}\label{eq:bs-collapse}
\frac{|\mathrm{BS}^\star(\Sigma_0^\star)|}{2(1-\beta)/\eta^\star}=O\big(\sqrt{\lambda_k}\big)\longrightarrow0\qquad(\lambda_k\downarrow0).
\end{equation}
\end{enumerate}
\end{theorem}
Proof in Appendix~\ref{app:thmFG}.

\begin{theorem}[Saturation in a slow-fast model: equilibrium at $\rho(\mathcal K_{\eta,\beta})=1$ against a Krylov--Bogoliubov secant equilibrium]\label{thm:slow-fast-saturation}
Fix $\eta>0$, $\beta\in[0,1)$ and $\Sigma_\xi\succ0$.
\begin{enumerate}
\item \emph{(Fast-epoch, quasi-stationary, and coupled stochastic progressive-sharpening barrier at $\rho(\mathcal K_{\eta,\beta})=1$.)} Let $y\in[y_{\min},y_{\max}]\subset\R$ parametrize $C^1$ families $H(y)\succ0$ and $\Delta H_i(y)=H_i(y)-H(y)$ in Theorem~\ref{thm:G6}, with threshold $\eta^\star(y)\triangleq\etacrit^{\mathrm{EoSS},P}(\beta,B;y)=1/\lambda_{\max}\big(\mathcal L_0(y)^{-1/2}\mathcal B_\beta(y)\mathcal L_0(y)^{-1/2}\big)$ and spectral radius $r(y)\triangleq\rho(\mathcal K_{\eta,\beta}(y))$. Suppose $\eta/\eta^\star(y)$ is strictly increasing on $[y_{\min},y_{\max}]$ and crosses $1$ at $y^\star\in(y_{\min},y_{\max})$ with $\eta<\eta_{\mathrm{det}}(y^\star)$, simple top eigenvalue $1/\eta^\star(y^\star)=1/\eta$, unit-trace eigenmatrix $\Sigma_0^\star\in\mathbb S_+^P$, and $g_0\triangleq-\partial_y(\eta/\eta^\star(y))\big|_{y=y^\star}<0$ (equivalently $g_\star\triangleq r'(y^\star)>0$). For the fast recursion \eqref{eq:hbP} at frozen $y$ and a stochastic slow gradient $g_{y,t}(w_t,y)=-v_{\mathrm{PS}}+\frac12w_t^\top\Phi_{B_t}(y)w_t+\xi_{y,t}$ with drive $v_{\mathrm{PS}}>0$, $\E[\Phi_{B_t}(y)\mid\mathcal F_{t-1}]=\Phi(y)\in\mathbb S^P$ $C^1$ in $y$ ($\Phi(y)=\partial_yH(y)$ for $\mathcal L(w,y)=\mathcal L_{\mathrm{slow}}(y)+\frac12w^\top H(y)w$) and $\E[\xi_{y,t}\mid\mathcal F_{t-1}]=0$, the $K$-step fast-epoch expected slow gradient $G_{\mathrm{slow},K}(y)\triangleq\E[\frac1K\sum_{t=1}^Kg_{y,t}(w_t,y)]$ satisfies:
(a) for every $y\ge y^\star$ where $\Phi(y)\succ0$, $G_{\mathrm{slow},K}(y)\ge -v_{\mathrm{PS}}+\frac18\lambda_{\min}(\Phi(y))\big(c(y)(K+1)-2\E\|w_0\|^2/K\big)\to+\infty$ as $K\to\infty$ with $c(y)>0$ from Corollary~\ref{cor:msinstab};
(b) for every $y\in[y_{\min},y^\star)$, $\lim_{K\to\infty}G_{\mathrm{slow},K}(y)$ equals the quasi-stationary slow drift
\begin{equation}\label{eq:slow-drift}
G_{\mathrm{slow}}(y)\triangleq -v_{\mathrm{PS}}+\frac12\tr\big(\Phi(y)\,\Sigma_0(y)\big)=-v_{\mathrm{PS}}+\frac{\bar c_\star\,\sigma_\xi^2}{1-\eta/\eta^\star(y)}+O(1)\qquad(y\uparrow y^\star),
\end{equation}
where $\Sigma_0(y)$ is the stationary covariance of Theorem~\ref{thm:F}(2), $\sigma_\xi^2\triangleq\tr(\Sigma_0^\star\Sigma_\xi)>0$, and $\bar c_\star\triangleq\frac{\eta\,\tr(\Phi(y^\star)\Sigma_0^\star)}{4(1-\beta)\,\tr(H(y^\star)(\Sigma_0^\star)^2)}$. If $\tr(\Phi(y^\star)\Sigma_0^\star)>0$, then $\lim_{y\uparrow y^\star}G_{\mathrm{slow}}(y)=+\infty$: for every $v_{\mathrm{PS}}>\frac12\tr(\Phi(y_{\min})\Sigma_0(y_{\min}))$ there is a quasi-stationary equilibrium $y_{\mathrm{eq}}\in(y_{\min},y^\star)$ with $G_{\mathrm{slow}}(y_{\mathrm{eq}})=0$, and as $\sigma_\xi^2/v_{\mathrm{PS}}\to0$ (writing $\Sigma_\xi=\sigma_\xi^2\widehat\Sigma_\xi$ with $\tr(\Sigma_0^\star\widehat\Sigma_\xi)=1$), the unique equilibrium near $y^\star$ satisfies $G'_{\mathrm{slow}}(y_{\mathrm{eq}})>0$ and, with $c_\star\triangleq d_\star\bar c_\star=\frac{\tr(\Phi(y^\star)\Sigma_0^\star)}{2\langle\Sigma_0^\star,\mathcal V_{y^\star}(\Sigma_0^\star)\rangle_F}>0$,
\begin{equation}\label{eq:saturation-law}
\frac{\eta^\star(y_{\mathrm{eq}})-\eta}{\eta}=\frac{\bar c_\star\,\sigma_\xi^2}{v_{\mathrm{PS}}}+O\!\left(\frac{\sigma_\xi^4}{v_{\mathrm{PS}}^2}\right)>0,\qquad
1-\rho\big(\mathcal K_{\eta,\beta}(y_{\mathrm{eq}})\big)=\frac{c_\star\,\sigma_\xi^2}{v_{\mathrm{PS}}}+O\!\left(\frac{\sigma_\xi^4}{v_{\mathrm{PS}}^2}\right)>0;
\end{equation}
and (c) if $\Phi(y)\succ0$ on $[y_{\min},y_{\max}]$ and $v_{\mathrm{PS}}>\frac12\tr(\Phi(y_{\min})\Sigma_0(y_{\min}))$, then for $G_{\mathrm{slow},K}$ evaluated from $w_0=w_{-1}=0$ and every finite $K>\max_{y\in[y^\star,y_{\max}]}\frac{8v_{\mathrm{PS}}}{\lambda_{\min}(\Phi(y))c(y)}$, every root of $G_{\mathrm{slow},K}$ on $[y_{\min},y_{\max}]$ lies strictly in the stable region $(y_{\min},y^\star)$; moreover, on any compact interval $[y_-,y_+]\subset(y_{\min},y^\star)$ containing a root $y_{\mathrm{eq}}$ of $G_{\mathrm{slow}}$ on which $G'_{\mathrm{slow}}(y)>0$ (such as the equilibrium near $y^\star$ in (b), or any compact subinterval containing the globally unique root $y_{\mathrm{eq}}\in(y_{\min},y^\star)$ when $y\mapsto\tr(\Phi(y)\Sigma_0(y))$ has positive derivative on $[y_{\min},y^\star)$), for all sufficiently large $K$ there is a unique root $y_{\mathrm{eq},K}\in[y_-,y_+]$ of $G_{\mathrm{slow},K}$ with $y_{\mathrm{eq},K}\to y_{\mathrm{eq}}$ as $K\to\infty$, and with $0<m_0\le\min_{y\in[y_-,y_+]}G'_{\mathrm{slow},K}(y)\le\max_{y\in[y_-,y_+]}G'_{\mathrm{slow},K}(y)\triangleq L_K<\infty$, $\Var(\frac1K\sum_{t=1}^Kg_{y,t}(w_{k,t},y_k)\mid\mathcal F_{k-1})\le s_K^2$, and any $\mathcal F_{k-1}$-measurable epoch initial state satisfying $\E[\|w_{k,0}\|^2+\|w_{k,-1}\|^2\mid\mathcal F_{k-1}]\le R_0^2$ ($R_0=0$ when each epoch resets $w_{k,0}=w_{k,-1}=0$; $R_0<\infty$ under continuous carry-over $(w_{k,0},w_{k,-1})=(w_{k-1,K},w_{k-1,K-1})$ on $[y_-,y_+]$ since $\sup_{y\in[y_-,y_+]}\rho(\mathcal T_y)<1$), there is a constant $C_+<\infty$ independent of $K$ such that the coupled stochastic epoch recursion $y_{k+1}=\Pi_{[y_-,y_+]}\big(y_k-\eta_y\frac1K\sum_{t=1}^Kg_{y,t}(w_{k,t},y_k)\big)$ for every slow step size $0<\eta_y\le1/L_K$ satisfies $\E[(y_k-y_{\mathrm{eq},K})^2]\le(1-\eta_ym_0)^k\E[(y_0-y_{\mathrm{eq},K})^2]+\frac{\eta_ys_K^2}{m_0}+\frac{2C_+^2R_0^4}{m_0^2K^2}$ for all $k\ge1$.
\item \emph{(Self-consistent Krylov--Bogoliubov secant equilibrium above the tangent threshold.)} Suppose the tangent curvature $H_{\mathrm{tan}}\succ0$ is supercritical ($\eta^\star_{\mathrm{tan}}<\eta<\eta_{\mathrm{det}}(H_{\mathrm{tan}})$), and under a Krylov--Bogoliubov secant linearisation \citep{krylov1947} parameterized by the excursion variance $\gamma=\tr\Sigma_0\in[0,\gamma_{\max}]$, the $C^1$ secant family $\big(H_{\mathrm{sec}}(\gamma),\Delta H_i^{\mathrm{sec}}(\gamma)\big)$ satisfies $\eta<\eta_{\mathrm{det}}(H_{\mathrm{sec}}(\gamma))$ and has $\eta/\eta^\star_{\mathrm{sec}}(\gamma)$ strictly decreasing from $\eta/\eta^\star_{\mathrm{tan}}>1$ at $\gamma=0$ to $\eta/\eta^\star_{\mathrm{sec}}(\gamma_{\max})<1$, crossing $1$ at $\gamma^\star\in(0,\gamma_{\max})$ with a simple top eigenvalue, unit-trace eigenmatrix $\Sigma_0^\star(\gamma^\star)$, and $\partial_\gamma\eta^\star_{\mathrm{sec}}(\gamma^\star)>0$. Then for every $\Sigma_\xi\succ0$ with $\tr\Sigma_0^{\mathrm{sec}}(\gamma_{\max})\le\gamma_{\max}$ (in particular for all small $\sigma_\xi^2=\tr(\Sigma_0^\star(\gamma^\star)\Sigma_\xi)>0$) there is a self-consistent excursion variance $\gamma_{\mathrm{eq}}\in(\gamma^\star,\gamma_{\max}]$ solving $\tr\Sigma_0^{\mathrm{sec}}(\gamma_{\mathrm{eq}})=\gamma_{\mathrm{eq}}$, with $\eta^\star_{\mathrm{sec}}(\gamma_{\mathrm{eq}})>\eta>\eta^\star_{\mathrm{tan}}$; as $\sigma_\xi^2\to0$, $\gamma_{\mathrm{eq}}$ is unique on $(\gamma^\star,\gamma_{\max}]$ and satisfies
\begin{equation}\label{eq:secant-escape-law}
1-\frac{\eta}{\eta^\star_{\mathrm{sec}}(\gamma_{\mathrm{eq}})}=\frac{\eta\,\sigma_\xi^2}{2(1-\beta)\,\gamma^\star\,\tr\big(H_{\mathrm{sec}}(\gamma^\star)\Sigma_0^\star(\gamma^\star)^2\big)}\big(1+O(\sigma_\xi^2)\big)>0.
\end{equation}
\end{enumerate}
\end{theorem}
Proof in Appendix~\ref{app:thmFG}.

\paragraph{What Theorems~\ref{thm:modal-bridge} and~\ref{thm:slow-fast-saturation} answer and what they do not.} (i) The Batch Sharpness of \citet[Definition~3.1]{andreyev2026momentum} is the expectation, over the minibatch at a fixed checkpoint, of the ratio $g_B^\top H_Bg_B/\|g_B\|^2$. Theorem~\ref{thm:modal-bridge}(2) is about a stationary ratio of expectations, averaged also over the iterate $w_t$ of the linearised recursion, and part (1)(a) about a ratio of expectations at a fixed checkpoint, which differs from the expectation of the ratio by $\operatorname{Cov}_B(R_B,W_B)$, at most $\sigma_{\mathrm{BS}}\mathrm{CV}_{g^2}$, and which equals $\bar a+\bar v$ under the hypothesis $\E[g_Bg_B^\top\mid H_B]=\sigma_g^2H_B$, which rank-one per-sample curvatures aligned with the per-sample gradients satisfy and the linearised gradient $H_Bw$ of the quadratic model in general does not. Part (1)(b) reaches the expectation of the ratio only in the random-feature model with $H_i=z_iz_i^\top$ and $g_i\propto z_i$, and is informative only when $B^2\lambda_1\ll\tr H$. (ii) In the general setting the theorem is a bridge, an identity at the threshold and the two limits of a stationary Batch Sharpness, together with a sufficient instability criterion through the variational test at matrices $S$ that commute with $H$; its Batch Sharpness form needs the noise share $\theta_{\mathrm{noise}}$ besides $\overline{\mathrm{BS}}$, so it is not a threshold on Batch Sharpness alone. (iii) Theorem~\ref{thm:slow-fast-saturation} proves saturation in a quasi-stationary slow-fast model: the curvature is frozen within an epoch, the slow gradient has the form $-v_{\mathrm{PS}}+\frac12w^\top\Phi w$, and the fast state restarts at each epoch or is carried over with bounded second moment; it is not a statement about the training dynamics. The $P=33$ runs of Section~\ref{sec:phase5}, which train above the $|H|$ reading, are outside the hypotheses of its part (1), since $H$ is indefinite at $90$--$100\,\%$ of their window checkpoints, and behave as its secant part (2) describes; whether they satisfy the hypotheses of part (2) was not tested. The variational criterion, Batch Sharpness limits and saturation law are checked numerically in \texttt{paper/checks/bridge-theorems.py}.

\section{Measurements}\label{sec:exp}

\subsection{The stochastic ceiling in \texorpdfstring{$\R^P$}{R\^P}: Gaussian against Rademacher features (Phase 1)}\label{sec:phase1}

Phase~1 measures Theorem~\ref{thm:G} on the quadratic model it is about: spectrum $\lambda_k=k^{-1/2}$, $P\in\{64,256\}$, $\beta\in\{0,0.9\}$, $B\in\{4,16,64\}$ ($B\le16$ at $P=256$), isotropic additive noise $\sigma_\xi=0.01$, and three feature families, Gaussian i.i.d.\ ($\kappa_{\mathrm{feat}}=0$), Rademacher i.i.d.\ ($\kappa_{\mathrm{feat}}=-2$, $\omega^4\equiv1$) and a finite pool of $n$ Gaussian features ($n=512$ at $P=64$, $1024$ at $P=256$) sampled without replacement, as in \citet{lee2022minibatch}, for which the prediction uses the eigenvalues of the pool covariance and the factor $1-\zeta$, $\zeta=B/n$, a prediction that replaces the noise operator of the realised pool by its average over pool draws (Section~\ref{sec:lyapunov}, after Theorem~\ref{thm:G6}). The empirical ceiling is found by bisection on $\eta$ (12 iterations, 8 streams of 1000 steps, a learning rate being declared unstable when the mean-square energy over the streams grows by more than $2.5\times$ between the second and the last quarter), 8 seeds per cell; the median over seeds of the bisection half-width is below $0.05\,\%$ in every cell. Four predictions were written before the run (\texttt{labs/mean-square-boundary/README.md}, Phase~1): \texttt{p1}, every i.i.d.\ cell inside the sandwich within $3\,\%$ (the registered text has the sandwich without the terms $\max(\alpha_{\mathrm{diag}},0)$ and $\min(\alpha_{\mathrm{diag}},0)$ of \eqref{eq:sandwich}, an interval that is empty when $\alpha_{\mathrm{diag}}<0$; \texttt{p1} is scored with \eqref{eq:sandwich}, a deviation written after the registration, and fails under either form); \texttt{p2}, the secular root within $5\,\%$ median error in every i.i.d.\ cell; \texttt{p3}, in every Rademacher cell with $\kappa_\beta\tr H/(B\lambda_1)\ge1$ the one-dimensional ceiling exceeds the empirical one by more than $2\times$ while the root stays within $5\,\%$; \texttt{p4}, in the pool arm with $\zeta\le1/8$ the root of $\|\mathcal K\|=1$ of \citet{lee2022minibatch} within $8\,\%$. Table~\ref{tab:phase1} has the 30 cells.

\begin{table}[!htbp]
\centering
\caption{Phase~1 (\texttt{raw-phase1.json}, 8 seeds per cell). $\rho_B=\kappa_\beta\tr H/(B\lambda_1)$; 1D is the ceiling of Proposition~\ref{prop:oneD} along $u_1$ ($\nu_{11}^2=\lambda_1^2(2+\kappa_{\mathrm{feat}})/B$); lower, root and upper are the sandwich \eqref{eq:sandwich} and the secular root \eqref{eq:secular} (pool: with the factor $1-\zeta$ and the pool eigenvalues, seed-averaged; this uses the noise operator averaged over pool draws, and the exact threshold of one realised pool is that of Theorem~\ref{thm:G6}); emp.\ is the median empirical ceiling; err is the median over seeds of the per-seed relative error of the root (not the error of the median ceiling), Lee the same for the root with $\alpha_{\mathrm{diag}}=1$; in is the criterion of \texttt{p1}, inside the sandwich within $3\,\%$ plus the bisection half-width; 1D/emp is the mean over seeds of the per-seed ratio of the one-dimensional ceiling to the empirical one (pool: $\nu_{11}^2=\lambda_1^2(2+\kappa_{\mathrm{feat}})(1-\zeta)/B$ with the pool $\lambda_1$, seed-averaged). The entry $2.0399$ of the Gaussian $P=64$, $\beta=0$, $B=64$ cell is censored: the bisection reached its cap $1.02\times2=2.04$ to within its resolution, and the entries marked $\dagger$ are computed against it, not against a measured ceiling.}
\label{tab:phase1}
\resizebox{\textwidth}{!}{%
\begin{tabular}{rrrlrrrrrrrrrl}
\toprule
$P$ & $\beta$ & $B$ & feat. & $\rho_B$ & 1D & lower & root & upper & emp. & err (\%) & Lee (\%) & 1D/emp & in \\
\midrule
64 & 0 & 4 & Gauss. & 3.7 & 1.3333 & 0.4081 & 0.4897 & 0.5479 & 0.4894 & 0.34 & 2.35 & 2.72 & yes \\
64 & 0 & 4 & Radem. & 3.7 & 2.0000 & 0.4545 & 0.5123 & 0.5479 & 0.5135 & 0.22 & 2.45 & 3.89 & yes \\
64 & 0 & 4 & pool & 3.4 & 1.2645 & 0.4052 & 0.4887 & 0.5533 & 0.4896 & 0.33 & 2.38 & 2.58 & yes \\
64 & 0 & 16 & Gauss. & 0.9 & 1.7778 & 1.0126 & 1.4468 & 1.7778 & 1.4862 & 2.65 & 0.47 & 1.20 & yes \\
64 & 0 & 16 & Radem. & 0.9 & 2.0000 & 1.0810 & 1.5299 & 2.0000 & 1.5544 & 1.57 & 4.26 & 1.29 & yes \\
64 & 0 & 16 & pool & 0.9 & 1.7422 & 1.0156 & 1.4341 & 1.7422 & 1.4421 & 1.00 & 1.97 & 1.20 & yes \\
64 & 0 & 64 & Gauss. & 0.2 & 1.9394 & 1.6080 & 1.9268 & 1.9394 & 2.0399 & $5.54^\dagger$ & $4.11^\dagger$ & $0.96^\dagger$ & no \\
64 & 0 & 64 & Radem. & 0.2 & 2.0000 & 1.6494 & 1.9861 & 2.0000 & 2.0004 & 0.71 & 2.22 & 1.00 & yes \\
64 & 0 & 64 & pool & 0.2 & 1.8870 & 1.6050 & 1.8778 & 1.8870 & 1.9408 & 3.82 & 2.53 & 0.97 & yes \\
64 & 0.9 & 4 & Gauss. & 69.4 & 0.3619 & 0.0506 & 0.0533 & 0.0548 & 0.0532 & 0.44 & 2.45 & 6.80 & yes \\
64 & 0.9 & 4 & Radem. & 69.4 & 3.8000 & 0.0548 & 0.0558 & 0.0579 & 0.0559 & 0.39 & 2.49 & 68.04 & yes \\
64 & 0.9 & 4 & pool & 67.4 & 0.3565 & 0.0509 & 0.0536 & 0.0552 & 0.0538 & 0.97 & 1.54 & 6.60 & yes \\
64 & 0.9 & 16 & Gauss. & 17.3 & 1.1259 & 0.1946 & 0.2103 & 0.2191 & 0.2113 & 0.59 & 1.81 & 5.30 & yes \\
64 & 0.9 & 16 & Radem. & 17.3 & 3.8000 & 0.2191 & 0.2199 & 0.2215 & 0.2196 & 0.78 & 2.06 & 17.26 & yes \\
64 & 0.9 & 16 & pool & 16.6 & 1.1002 & 0.1989 & 0.2155 & 0.2256 & 0.2170 & 1.69 & 1.55 & 5.07 & yes \\
64 & 0.9 & 64 & Gauss. & 4.3 & 2.3843 & 0.6747 & 0.7948 & 0.8766 & 0.7971 & 0.77 & 1.99 & 2.99 & yes \\
64 & 0.9 & 64 & Radem. & 4.3 & 3.8000 & 0.7543 & 0.8315 & 0.8766 & 0.8428 & 1.44 & 3.54 & 4.53 & yes \\
64 & 0.9 & 64 & pool & 4.2 & 2.4203 & 0.7454 & 0.8899 & 1.0012 & 0.9027 & 1.44 & 1.05 & 2.68 & yes \\
256 & 0 & 4 & Gauss. & 7.6 & 1.3333 & 0.2249 & 0.2531 & 0.2617 & 0.2544 & 0.53 & 0.15 & 5.24 & yes \\
256 & 0 & 4 & Radem. & 7.6 & 2.0000 & 0.2383 & 0.2565 & 0.2617 & 0.2578 & 0.49 & 1.15 & 7.76 & yes \\
256 & 0 & 4 & pool & 7.3 & 1.2734 & 0.2241 & 0.2528 & 0.2627 & 0.2534 & 0.19 & 0.68 & 5.03 & yes \\
256 & 0 & 16 & Gauss. & 1.9 & 1.7778 & 0.6727 & 0.9268 & 1.0467 & 0.9269 & 0.31 & 0.79 & 1.92 & yes \\
256 & 0 & 16 & Radem. & 1.9 & 2.0000 & 0.7022 & 0.9416 & 1.0467 & 0.9450 & 0.37 & 1.15 & 2.12 & yes \\
256 & 0 & 16 & pool & 1.8 & 1.7219 & 0.6708 & 0.9232 & 1.0624 & 0.9220 & 0.25 & 0.77 & 1.86 & yes \\
256 & 0.9 & 4 & Gauss. & 145.2 & 0.3619 & 0.0252 & 0.0260 & 0.0262 & 0.0261 & 0.68 & 0.26 & 13.86 & yes \\
256 & 0.9 & 4 & Radem. & 145.2 & 3.8000 & 0.0262 & 0.0263 & 0.0269 & 0.0264 & 0.49 & 1.13 & 143.68 & yes \\
256 & 0.9 & 4 & pool & 136.9 & 0.3422 & 0.0252 & 0.0260 & 0.0262 & 0.0262 & 0.82 & 0.49 & 13.06 & yes \\
256 & 0.9 & 16 & Gauss. & 36.3 & 1.1259 & 0.0987 & 0.1034 & 0.1047 & 0.1039 & 0.44 & 0.22 & 10.85 & yes \\
256 & 0.9 & 16 & Radem. & 36.3 & 3.8000 & 0.1047 & 0.1048 & 0.1052 & 0.1051 & 0.64 & 0.99 & 36.15 & yes \\
256 & 0.9 & 16 & pool & 35.1 & 1.1017 & 0.0999 & 0.1047 & 0.1062 & 0.1049 & 0.39 & 0.43 & 10.49 & yes \\
\bottomrule
\end{tabular}}
\end{table}

Outcomes: \texttt{p3} and \texttt{p4} hold, \texttt{p1} and \texttt{p2} fail, both on the same cell, and are recorded as failed. In 29 of the 30 cells the median empirical ceiling lies inside the sandwich within that tolerance, and strictly inside in 27: the Rademacher cell at $P=64$, $\beta=0$, $B=64$ is $0.02\,\%$ above its upper bound $2$, $1.6$ times the median bisection half-width of that cell ($0.012\,\%$) and within the tolerance of \texttt{p1}, and the pool cell at the same $(P,\beta,B)$ is $2.85\,\%$ above its seed-averaged upper bound $1.8870$, the effect of the failed cell at a smaller size. The secular root predicts the median with median error between $0.19\,\%$ and $3.82\,\%$, the upper end from that pool cell ($0.19$--$2.65\,\%$ in the other 28 cells) (the Lee root, without the kurtosis term, between $0.15\,\%$ and $4.26\,\%$); the bisection resolution is $0.05\,\%$, so these errors are the finite-horizon bias of the divergence criterion and the seed-to-seed spread, not resolution. The Rademacher separation is as predicted: at $\beta=0.9$ the one-dimensional ceiling along $u_1$, which is $2(1+\beta)/\lambda_1=3.8$ for every $B$ because $\nu_{11}^2=0$, exceeds the measured ceiling by $68\times$ ($P=64$, $B=4$), $17\times$ ($B=16$), $4.5\times$ ($B=64$), $144\times$ ($P=256$, $B=4$) and $36\times$ ($B=16$), while the secular root is within $1.5\,\%$ in all five cells; the cross-mode couplings $\lambda_1\lambda_k/B$ set the ceiling, and a probe restricted to the ray along $u_1$ cannot see them. The two plateaus are visible along $B$ at $\beta=0$: at $B=4$ the ceiling is $\tr H/B$-dominated ($\rho_B=3.7$, $7.6$) and at $B=64$ it is within $4\,\%$ of the deterministic $2/\lambda_1$.

The failed cell is $P=64$, $\beta=0$, $B=64$ with Gaussian features: the measured ceiling, $2.0399$, is above the deterministic ceiling $2/\lambda_1=2$ (the bisection cap is $1.02\times$ that value), which no mean-square ceiling can be. The diagnosis is the gap between mean-square and almost-sure stability of a random-coefficient recursion: along $u_1$ the multiplier $1-\eta(\lambda_1+\zeta_t)$ has, at $\eta=2.00$--$2.04$ and $\Var(\zeta_t)=2\lambda_1^2/B$, a negative Lyapunov exponent $\E\log|1-\eta(\lambda_1+\zeta_t)|$ while its second moment grows, so the typical trajectory decays (over the $\chi^2_{64}$ law of $\zeta_t$ the exponent is $-0.072$, $-0.051$, $-0.030$ at $\eta=2.00$, $2.02$, $2.04$ and $+0.031$ at $2.10$, \texttt{paper/checks/evidence-ranges.mjs}) and the mean square diverges through rare paths that 8 streams do not sample. For the homogeneous recursion along $u_1$ at $\beta=0$ the multipliers are i.i.d., so $\E[x_t^4]/\E[x_t^2]^2=(m_4/m_2^2)^t$ with $m_q=\E|1-\eta(\lambda_1+\zeta_t)|^q$; over the $\chi^2_{64}$ law $m_4/m_2^2=1.498$ at $\eta=1.94$, so the ratio is $e^{404}$ at $t=1000$, and the relative standard error of a mean of $x_t^2$ over 8 paths is of order $e^{202}$ (\texttt{paper/checks/second-moment-renewal.py}). Resampling targets the second moment directly. Carry $S$ unit vectors $u_{t,s}\in\R^{2P}$, apply an independent step $y_{t,s}=F_{t,s}u_{t,s}$, record $\hat\mu_t=\frac1S\sum_s\|y_{t,s}\|^2$, and draw $u_{t+1,1},\dots,u_{t+1,S}$ from $\{y_{t,s}/\|y_{t,s}\|\}$ with probabilities $\|y_{t,s}\|^2/(S\hat\mu_t)$; with $\widehat M_t=\prod_{\tau<t}\hat\mu_\tau$ and $\widehat Q_t=\widehat M_t\frac1S\sum_su_{t,s}u_{t,s}^\top$, $\E[\widehat Q_{t+1}\mid\mathcal F_t]=\mathcal T(\widehat Q_t)$, so $\E\widehat Q_t=\mathcal T^t(\widehat Q_0)$ for every $S$ and $t$ (checked by exact enumeration in the same script). Along $u_1$ at $\beta=0$ the $\hat\mu_\tau$ are i.i.d.\ with mean $m_2$ and variance $(m_4-m_2^2)/S$: the relative variance of $\widehat M_t$ is $(1+(m_4/m_2^2-1)/S)^t-1$, still $e^{60}$ at $S=8$ and $t=1000$, but $\frac1T\sum_{\tau<T}\hat\mu_\tau$ estimates $m_2$ with variance $(m_4-m_2^2)/(ST)$, and its crossing of $1$ along $u_1$ is the one-dimensional mean-square ceiling $1.9394$ of this cell. On the full operator of this cell, with $S=8$ and $K=64$ independent blocks of $2m=16$ steps (the budget of one Phase~1 cell), the ratio $\widehat R=\overline{\widehat M_{2m}}/\overline{\widehat M_m}$ of block means tends as $K\to\infty$ to $\tr\mathcal T^{2m}(Q_0)/\tr\mathcal T^m(Q_0)$, which by the exact recursion crosses $1$ at $1.926417$ for $m=8$ (\texttt{paper/checks/second-moment-renewal.py}) and at the secular root as $m$ grows. A 12-step bisection of $\widehat R=1$ on 8 seeds gives a median of $1.9330$ (range $1.9200$--$1.9400$), $0.32\,\%$ above the secular root $1.9268$, where the sample paths of Phase~1 gave $2.003$ with the cap raised; the three predictions written before the run (median within $1\,\%$ of the root, every seed within $3\,\%$, median below $2.003$) hold (\texttt{paper/checks/phase1-resampling-estimator.py}, \texttt{phase1-resampling-estimator.json}). The Rademacher cell at the same $(P,\beta,B)$ has $\nu_{11}^2=0$, no such gap, and matches the root within $0.71\,\%$; the Gaussian cell at $B=16$ shows the same effect at a smaller size ($2.65\,\%$, the empirical value above the root). Proposition~\ref{prop:G3}(1) makes this quantitative: along $u_1$ the mean-square and almost-sure thresholds sit at $\frac2{\lambda_1}(1\mp\nu_{11,B}^2/\lambda_1^2)$ to first order ($1.9375$ and $2.0625$ here; exactly $1.9394$ and, by quadrature over the $\chi^2_{64}$ density, $2.0688$, \texttt{paper/checks/evidence-ranges.mjs}), and the measured $2.0399$ is the bisection cap $1.02\times2=2.04$ to within the bisection resolution, a censored observation above the mean-square ceiling. A post-hoc rerun of this cell with the cap raised to $2.3$, not a pre-registered prediction (\texttt{evidence/mean-square-boundary/phase1-diag-as-vs-ms.json}, written by \texttt{run-phase1-cap-diagnostic.ts}), gives a sample-path threshold of median $2.003$ (range $1.963$--$2.067$) for Gaussian and $2.000$ ($1.976$--$2.013$) for Rademacher features. For Gaussian features the median lies between the mean-square ceiling and the one-dimensional almost-sure root, whose order Proposition~\ref{prop:G3}(1) gives, near the middle of that interval; a sample-path harness of this length does not estimate the almost-sure root; for Rademacher features both thresholds are $2$. Theorem~\ref{thm:G} is a statement about the second moment; where the multiplicative noise along $u_1$ dominates and $B$ is large, a sample-path harness measures a threshold above the mean-square one (here between the mean-square ceiling and the one-dimensional almost-sure root), and the resampling estimator above, which targets $\E\|w_t\|^2$ itself, measures the theorem there.

\paragraph{The secular root against the exact second-moment operator (Phase 3).}
Phase~1 compares Theorem~\ref{thm:G} with sample paths, whose threshold can lie above the mean-square one where the multiplicative noise along $u_1$ is large. Phase~3 compares it with the mean-square threshold itself: the spectral radius of the second-moment operator $\mathcal T$, computed by the matrix-free power iteration from the identity of Lemma~\ref{lem:russodye}, with one application of $\mathcal T$ in $O(P^2)$ operations in the eigenbasis of $H$ for random-feature minibatches (checked against the exact expectation by enumeration on a non-symmetric input; the operator uses the closed form of $\mathcal V$ of Theorem~\ref{thm:G}, so this phase tests the diagonal closure, the secular reduction and the code, and the enumeration tests that closed form). On the spectrum $\lambda_k=k^{-1/2}$ of Phase~1, with $P\in\{64,256,1024\}$, $\beta\in\{0,0.9\}$, $B\in\{4,16,64\}$ and Gaussian and Rademacher features (36 cells), $\rho(\mathcal T)$ is below $1$ at $(1-10^{-3})\eta^\star$ and above $1$ at $(1+10^{-3})\eta^\star$ in every cell, $\eta^\star$ the secular root; the crossing located from $\rho(\mathcal T(\eta^\star))-1$ and the local slope is within $6.11\times10^{-10}$ (relative) of $\eta^\star$ in every cell (six cells at $P=64$, $\beta=0.9$, $B\in\{4,16,64\}$ for both feature families have \texttt{allConverged} false, the conjunction of four convergence flags; a rerun, \texttt{check-phase3-convergence.ts}, shows that the flag that fails is the agreement test of the Gelfand estimate on the antisymmetric matrices, while the three symmetric iterations, which locate the crossing, converge in $248$--$275$ steps from a cold start (the run, warm-started from the iterate at $(1-10^{-3})\eta^\star$, recorded $175$--$208$ at the root); these cells have crossing relative errors between $3.2\times10^{-14}$ and $5.1\times10^{-13}$ and $|\rho(\mathcal T(\eta^\star))-1|\le3.8\times10^{-14}$), and the radius on the antisymmetric matrices is $0.65$--$0.99$ at $\eta^\star$, below the symmetric radius $1$, as Lemma~\ref{lem:russodye} requires. Three predictions were written before the full run (\texttt{raw-phase3.json}, field \texttt{predictionTexts}): \texttt{q1}, in every cell the radius is below $1$ at $(1-10^{-3})\eta^\star$ and above $1$ at $(1+10^{-3})\eta^\star$; \texttt{q2}, the crossing within $10^{-8}$ (relative) of $\eta^\star$; \texttt{q3}, the radius on the antisymmetric matrices at $\eta^\star$ strictly below $1$. All three hold. They were written after a four-cell prototype at $P\le256$, which the raw file records; after them a $P=16$ debugging run, in which the ratio of norms did not converge on the antisymmetric part at $\beta=0.9$, led to replacing the antisymmetric estimator by the Gelfand estimate over a fixed horizon before the full run, so \texttt{q3} was seen at $P=16$ first; the $P=1024$ cells, beyond the reach of the Phase~1 simulation, were not run before them (\texttt{raw-phase3.json}, field \texttt{provenanceNote}).

\subsection{Trace and the ceiling in \texorpdfstring{$\R^P$}{R\^P} on two networks (Phase 2)}\label{sec:phase2}

Phase~2 asks whether the random-feature model of Theorem~\ref{thm:G} describes real per-sample Hessians. The checkpoints are two networks of the companion manuscript: a $64\text{--}48\text{--}10$ network ($P=3\,610$) and a $256\text{--}48\text{--}10$ network ($P=12\,826$), both with $\tanh$ hidden units, squared loss and $N=64$ samples of a synthetic teacher--student regression with ten outputs, with $\lambda_1=98.677$ and $148.353$ at $\Theta_0$ from analytical Hessian-vector products \citep{pearlmutter1994}; $\etacrit^{\mathrm{quad}}=2(1+\beta)/\lambda_1$ is the deterministic threshold. Predictions were written before the run (\texttt{raw-phase2.json}, field \texttt{predictionTexts}); of the three reported here \texttt{q1} and \texttt{q2} hold and \texttt{q3} fails. The others concern the estimator of the companion manuscript.

\paragraph{Trace and noise-rank ratio (\texttt{q1}).} The trace is exact, $\sum_ie_i^\top He_i$ from $P$ analytical Hessian-vector products: $\tr H=601.74$ at $P=3\,610$ and $901.72$ at $P=12\,826$, so $\tr H/\lambda_1=6.10$ and $6.08$ and $\rho_B=7.24$ and $7.22$ at $\beta=0.9$, $B=16$. One hundred shifted power iterations give Rayleigh quotients $0.012$ and $-0.072$, which are upper bounds on $\lambda_{\min}$: at $P=12\,826$ the Hessian is indefinite ($\lambda_{\min}\le-0.072=-4.9\times10^{-4}\lambda_1$); at $P=3\,610$ the value certifies nothing about the sign, and the dense spectrum of Section~\ref{sec:phase4} has $\lambda_{\min}=-0.0833$ with $1\,528$ negative eigenvalues. Both Hessians are indefinite, so the lower ends of the trace brackets below are approximations (Remark~\ref{rem:tracebracket}; the guaranteed lower end at $P=3\,610$ is $0.00411$). The effective feature kurtosis read from $\nu_{11}$ ($\nu_{11}^2=\lambda_1^2(2+\kappa_{\mathrm{feat}})/B$) is $-1.16$ and $-0.32$, $\alpha_{\mathrm{diag}}=0.81$ and $1.80$. With $\lambda_1$, $\lambda_2$ and $\tr H$ alone the secular root of Theorem~\ref{thm:G} is bracketed (Appendix~\ref{app:thmFG}, trace bracket): $[0.00492,0.00521]$ and $[0.00308,0.00338]$, which is $3.7\times$ and $2.5\times$ below the untethered ($\mu=0$) one-dimensional ceiling $\etacrit^{\mathrm{EoSS}}(0.9,16)=2(1+\beta)\lambda_1/(\lambda_1^2+\kappa_\beta\nu_B^2)=0.01928$ and $0.008564$.

\paragraph{Heavy-ball SGD with real minibatches in $\R^P$ (\texttt{q2}, \texttt{q3}).} The iteration is $w_{t+1}=w_t+\beta(w_t-w_{t-1})-\eta\,(g_{B_t}(w_t)-g_{B_t}(\Theta_0))$ with minibatches of $16$ of the $64$ samples drawn with replacement: the curvature noise is the one of the network and the mean gradient is removed, so $\Theta_0$ is a fixed point and the question is its linear stability under that noise. Every stream starts at distance $10^{-6}$ in the range of $H$ (a Gaussian start would put part of its energy in the null space of $H$, where heavy-ball keeps it constant; at $P=12\,826$, where the input dimension $256$ exceeds the $64$ samples, first-layer weights orthogonal to every input span a null space of every $H_i$ of dimension at least $48\cdot192=9\,216$), and the statistic is the one of Phase~1: the log ratio of the energy in the fourth and second quarters of a run, summed over 8 streams of $1\,000$ steps, bisected over $\eta$ in 14 steps from $1.1\times$ the quadratic threshold, 4 seeds. The median values are $0.002707$ at $P=3\,610$ and $4.90\times10^{-5}$ at $P=12\,826$ (bisection half-widths $0.05\,\%$ and $1.8\,\%$): $7.1\times$ and $175\times$ below the one-dimensional ceiling and $14\times$ and $523\times$ below the quadratic threshold. \texttt{q2} (a factor of at least $1.5$ below the one-dimensional ceiling) holds; \texttt{q3} (inside the trace bracket of the random-feature model within $10\,\%$) fails, by $1.8\times$ at $P=3\,610$ and $63\times$ at $P=12\,826$.

\paragraph{Growth grid (post hoc, no prediction registered).} On a fixed grid of learning rates, 4 streams, $1\,000$ and $4\,000$ steps, the log growth at $P=3\,610$ is negative at both lengths for $0.002$--$0.01\times\etacrit^{\mathrm{quad}}$ ($-0.99$ to $-0.22$), is $-0.88$ then $+0.69$ at $0.02\times$, $-0.42$ then $+5.8$ at $0.05\times$, and $+2.7$ at $0.1\times$ with the excursion guard tripped at $4\,000$ steps; at $P=12\,826$ it is $+1.0$ then $+10.2$ at $0.002\times$ and $+6.0$ at $0.005\times$ with the guard tripped at $4\,000$ steps. A diffusive process gives the same window ratio at any length; the growth here increases with the length, so it is exponential, and the bisection value is the learning rate at which it reaches the ratio $2.5$ within $1\,000$ steps, not a boundary. This finite-horizon statistic changes sign between $0.01\times$ and $0.02\times$ the quadratic threshold at $P=3\,610$ ($3.9\times10^{-4}$ to $7.7\times10^{-4}$, $25$--$50\times$ below the one-dimensional ceiling) and below $0.002\times$ at $P=12\,826$ (under $5.2\times10^{-5}$, more than $167\times$ below). It is not a mean-square stability boundary: the mean Hessian has negative eigenvalues at both checkpoints ($1\,528$ at $P=3\,610$, Section~\ref{sec:phase4}; $\lambda_{\min}\le-0.072$ at $P=12\,826$), so the mean of the untethered iteration has a root above $1$ at every $\eta>0$ and no learning rate is mean-square stable; the grid measures how fast energy escapes along the negative directions within $1\,000$--$4\,000$ steps.

\paragraph{What this means for the ceiling.} The one-dimensional ceiling of Proposition~\ref{prop:oneD} and the secular root of Theorem~\ref{thm:G} are statements about models in which the curvature noise entering a mode is proportional to that mode's curvature ($z=H^{1/2}\omega$ puts noise $\lambda_k\omega_k^2$ on mode $k$ and none on a flat mode). On these networks the per-sample Hessians $H_i$ inject curvature noise into directions where $H$ itself is nearly flat; Theorem~\ref{thm:G5} makes this quantitative: with a curvature-noise covariance $\Sigma_z\ne H$ the ceiling is at most $2(1-\beta)B/\tr(H^{-1}\Sigma_z^2)$ when every $d_k\ge0$, and for every $\kappa_{\mathrm{feat}}\ge-2$ and $P\ge2$ a mode with $\lambda_k\to0$ at fixed noise $s_k$ drives it to zero at least as fast as $\sqrt{\lambda_k}$, since $\eta^\star\le2(1-\beta)B\sqrt{\lambda_j\lambda_k}/(s_js_k)$ for any other mode $j$; such a direction has little restoring force and receives energy at every learning rate. At these two checkpoints, however, the mean Hessian also has negative directions, and Phase~4 (Section~\ref{sec:phase4}) shows at $P=3\,610$ that once a tether lifts them the exact threshold is $1.23\times$ below the one-dimensional ceiling, with no flat-mode collapse down to a smallest tethered eigenvalue of $\frac1{128}$: the growth measured here at $25$--$50\times$ and more than $167\times$ below that ceiling comes from the negative directions, not from curvature noise on positive flat modes. The excursions of actual training at the learning rates practitioners use come from the nonlinearity (the loss is bounded below and the mean gradient is not zero), not from a linear ceiling. In $\R^P$ the linear cap at these checkpoints is one to two orders of magnitude below the one-dimensional ceiling. Theorem~\ref{thm:G} is unaffected: Phase~1 measures it on the model it is about, and Phase~2 measures where the model stops applying.

\begin{table}[!htbp]
\centering\small
\caption{Phase~2 (\texttt{raw-phase2.json}; $\beta=0.9$, $B=16$, 4 seeds, 8 streams, $1\,000$ steps). 1D is the ceiling of Proposition~\ref{prop:oneD} along $u_1$; the bracket is the trace bracket of the secular root (Appendix~\ref{app:thmFG}) with $\kappa_{\mathrm{feat}}$ read from $\nu_{11}$; emp.\ is the median learning rate at which the energy of the linearised minibatch iteration grows by the Phase~1 ratio within the run, which is not a stability boundary (see the text).}
\label{tab:phase2}
\resizebox{\textwidth}{!}{%
\begin{tabular}{rrrrrrrrrr}
\toprule
$P$ & $\tr H$ & $\tr H/\lambda_1$ & $\rho_B$ & $\lambda_{\min}$ (upper bd.) & $\etacrit^{\mathrm{quad}}$ & 1D & bracket & emp. & 1D/emp \\
\midrule
3\,610 & 601.74 & 6.10 & 7.24 & 0.012 & 0.038509 & 0.019278 & [0.00492, 0.00521] & 0.002707 & 7.1 \\
12\,826 & 901.72 & 6.08 & 7.22 & $-0.072$ & 0.025615 & 0.008564 & [0.00308, 0.00338] & 0.0000490 & 175 \\
\bottomrule
\end{tabular}}
\end{table}

\subsection{The exact threshold of Theorem~\ref{thm:G6} on tethered network Hessians (Phase 4)}\label{sec:phase4}

Phase~2 measured an iteration for which the theory has no finite target: the full eigendecomposition of the mean Hessian of the $P=3\,610$ network gives $\lambda_{\min}=-0.0833$ with $1\,528$ of its $3\,610$ eigenvalues negative (prediction \texttt{t2} below), and the $P=33$ network of the companion manuscript (6 inputs, 4 $\tanh$ hidden units, one output, squared loss, $N=64$ samples of a synthetic teacher--student regression) has $\lambda_{\min}=-0.0080$ with two negative eigenvalues. A tether $\frac\mu2\|w-\Theta_0\|^2$ with $\mu=g+\max(0,-\lambda_{\min}(H))$ makes $H+\mu I\succ0$ with smallest eigenvalue $g$ and leaves $\Delta H_i$, hence $\mathcal V$, unchanged, and Theorem~\ref{thm:G6}(3) then gives one finite threshold $\eta^\star(\mu)=1/\lambda_{\max}(\mathcal L_0^{-1/2}\mathcal B_\beta\mathcal L_0^{-1/2})$. Phase~4 computes it by Lanczos from the real per-sample Hessians: at $P=33$ dense matrices from central differences of per-sample gradients ($\epsilon=10^{-5}$, symmetrised), $\beta\in\{0,0.9\}$; at $P=3\,610$ exact per-sample Hessian-vector products \citep{pearlmutter1994}, $\beta=0.9$; $B=16$ with replacement, $g\in\{128,32,8,2,\frac12,\frac18,\frac1{32},\frac1{128}\}$. Six predictions were written before the run (\texttt{raw-phase4.json}); \texttt{t1}--\texttt{t4} and \texttt{t6} hold and \texttt{t5} fails.

\paragraph{The formula against the exact operator (\texttt{t1}).} In all 16 cells at $P=33$ the spectral radius of the second-moment operator $\mathcal T$ of the tethered iteration is below $1$ at $(1-10^{-3})\eta^\star$ (largest $0.99988515$) and above $1$ at $(1+10^{-3})\eta^\star$ (smallest $1.00011489$), and a floating-point Collatz--Wielandt bracket of the last power iterate decides each side. A post-hoc run with the certified bracket of Section~\ref{sec:phase6} (\texttt{run-phase4-certified-bracket-p33.ts}, output \texttt{phase4-certified-bracket-p33.json}; run on 1 November 2026 after the expectation, both values inside the bracket and $\eta_{\mathrm{low}}/\eta_{\mathrm{up}}\ge0.98$, was written in the header of the runner and committed in the lab) gives in all 16 cells an interval $[\eta_{\mathrm{low}},\eta_{\mathrm{up}}]$ that contains the stored $\eta^\star$ and a dense reference, with $\eta_{\mathrm{low}}/\eta_{\mathrm{up}}\ge0.99931$; both offsets $(1\mp10^{-3})\eta^\star$ lie outside it, so through Theorem~\ref{thm:G6} each side of \texttt{t1} is decided with every rounding error bounded, for the computed per-sample Hessians (central differences) and not for the exact network Hessians. At $P=3\,610$, $g=2$, the power-iteration radii are $0.99982427$ and $1.00017600$ ($439$ and $237$ steps, converged to $10^{-12}$), but the last iterate is not positive definite in floating point, at $\pm10^{-3}$ and at the fallback offset $\pm10^{-2}$, so the bracket does not apply (the implementation checks \texttt{c2} and \texttt{c3} fail): at that size the crossing rests on the power iteration alone.

\paragraph{Sample paths (\texttt{t3}, \texttt{t4}).} With the statistic of Phases~1 and 2, the median sample-path ceiling of the linearised tethered iteration lies at $1.0008$--$1.190\,\eta^\star$ at $P=33$ and at $1.081$--$1.257\,\eta^\star$ at $P=3\,610$, never below the mean-square threshold, as Jensen's inequality requires of the almost-sure one (Proposition~\ref{prop:G3}). The network iteration started at distance $10^{-6}$ from the checkpoint gives the same median ceiling as the linearised one to the printed digits in 23 of 24 cells and $1.0053\times$ it in the remaining one ($P=3\,610$, $g=32$).

\paragraph{The tether sweep (\texttt{t5}, \texttt{t6}).} The fraction $f^\star=\eta^\star(\lambda_1+\mu)/(2(1+\beta))$ of the deterministic threshold rises with $\mu$ on both networks: from $0.891$ to $0.9998$ ($P=33$, $\beta=0$), from $0.171$ to $0.9965$ ($P=33$, $\beta=0.9$) and from $0.4065$ to $0.835$ ($P=3\,610$). $\eta^\star$ itself has an interior maximum, at $g=8$ ($P=33$, $\beta=0.9$, $\eta^\star=0.1941$) and at $g=32$ ($P=3\,610$, $\eta^\star=0.01788$). The flat-mode collapse of Theorem~\ref{thm:G5} was predicted at $P=3\,610$ between $g=\frac1{32}$ and $\frac1{128}$ as a log--log slope of $\eta^\star$ against $g$ in $[0.4,1.1]$; $\eta^\star$ goes from $0.0156894$ to $0.0156408$, a slope of $0.0022$, and \texttt{t5} fails. On this grid the critical eigenmatrix lies mostly on the $175$ rigid modes (eigenvalues above $2|\lambda_{\min}|$) and on the positive near-flat block ($64\,\%$ and $27\,\%$ of its weight at $g=\frac1{32}$, $6\,\%$ on the negative directions), and a diagonal upper bound computed after the run places the collapse below $g\approx6\times10^{-4}$, outside the grid. At the smallest tether the threshold is $1.23\times$ below the one-dimensional ceiling ($0.019278$) at $P=3\,610$ and $2.11\times$ below it at $P=33$, $\beta=0.9$. In Phase~2 the untethered sample-path ceiling at $P=3\,610$ was $0.002707$, $7.1\times$ below that ceiling; with the $1\,528$ negative directions lifted by $\mu\ge0.0833$ it is $0.0156$. On this checkpoint, then, most of the gap of Phase~2 comes from the negative directions of $H$, which a mean-square analysis of a fixed point does not cover, and not from curvature noise on positive flat modes.

\begin{table}[!htbp]
\centering
\caption{Phase~4 (\texttt{raw-phase4.json}; $B=16$). $g$ is the smallest eigenvalue of $H+\mu I$; 1D is the one-dimensional ceiling of Proposition~\ref{prop:oneD} of the tethered iteration; $\eta^\star$ is Theorem~\ref{thm:G6}(3) by Lanczos; $f^\star=\eta^\star/\etacrit^{(\mu)}(\beta)$; lin.\ and net.\ are the median sample-path ceilings of the linearised and of the network iteration over $\eta^\star$.}
\label{tab:phase4}
\resizebox{\textwidth}{!}{%
\begin{tabular}{rrrrrrrrr}
\toprule
$P$ & $\beta$ & $g$ & $\etacrit^{(\mu)}(\beta)$ & 1D & $\eta^\star$ & $f^\star$ & lin./$\eta^\star$ & net./$\eta^\star$ \\
\midrule
33 & 0 & 128 & 0.0149509 & 0.0149484 & 0.014948331 & 0.9998 & 1.001 & 1.001 \\
33 & 0 & 2 & 0.257374 & 0.244788 & 0.24286101 & 0.9436 & 1.107 & 1.107 \\
33 & 0 & $\frac1{128}$ & 0.346103 & 0.316663 & 0.30841041 & 0.8911 & 1.124 & 1.124 \\
33 & 0.9 & 128 & 0.0284068 & 0.0283135 & 0.028307600 & 0.9965 & 1.003 & 1.003 \\
33 & 0.9 & 8 & 0.275946 & 0.210477 & 0.19406782 & 0.7033 & 1.104 & 1.104 \\
33 & 0.9 & 2 & 0.489010 & 0.247371 & 0.18081066 & 0.3697 & 1.190 & 1.190 \\
33 & 0.9 & $\frac1{128}$ & 0.657597 & 0.237703 & 0.11242063 & 0.1710 & 1.183 & 1.183 \\
3\,610 & 0.9 & 128 & 0.0167578 & 0.0140951 & 0.013993390 & 0.8350 & 1.081 & 1.081 \\
3\,610 & 0.9 & 32 & 0.0290608 & 0.0185325 & 0.017880637 & 0.6153 & 1.196 & 1.202 \\
3\,610 & 0.9 & 2 & 0.0377132 & 0.0192735 & 0.016781922 & 0.4450 & 1.257 & 1.257 \\
3\,610 & 0.9 & $\frac1{128}$ & 0.0384739 & 0.0192781 & 0.015640812 & 0.4065 & 1.204 & 1.204 \\
\bottomrule
\end{tabular}}
\end{table}

\subsection{Does training saturate the threshold? (Phases 5 and 5b)}\label{sec:phase5}

Theorem~\ref{thm:G6} is a statement about a fixed point. Phases~5 and 5b evaluate it at the iterates of a training run, the first of the questions left open in Section~\ref{sec:lyapunov}. The $P=33$ network is trained from a fresh initialisation ($W_1\sim N(0,1/d_{\mathrm{in}})$, $W_2\sim N(0,1/d_{\mathrm{hidden}})$, zero biases) by heavy-ball SGD $\theta_{t+1}=\theta_t+\beta(\theta_t-\theta_{t-1})-\eta(g_{B_t}(\theta_t)+\mu_{\mathrm{train}}\theta_t)$ with $\beta=0.9$, $B=16$ drawn with replacement, $20\,000$ steps, $\eta\in\{0.01,0.03,0.1,0.15,0.2,0.3,0.4,0.5\}$ and weight decay $\mu_{\mathrm{train}}\in\{0,10^{-3}\}$, three seeds per cell, 48 runs. A run is divergent if its full loss becomes non-finite or exceeds $10^3$; none did. Every $200$ steps the Hessian $H$ of the training objective and the $64$ per-sample Hessians are computed at $\theta_t$, and Theorem~\ref{thm:G6} is evaluated with mean curvature $|H|$ (the eigenvalues of $H$ replaced by their absolute values) and noise $\Delta H_i=H_i-H$; we write $\eta^\star_{|H|}(\theta_t)$. $H$ is indefinite at most checkpoints of the runs with $\eta\ge0.15$, so $|H|$ is a choice of the diagnostic, not the mean curvature of the iteration. The diagnostic is quasi-static: given $\theta_t$ the minibatch Hessian has the law of Theorem~\ref{thm:G6}, but the drift of $\theta_t$ is outside the theorem. Per run the statistic is the median $m_\star$ of $\eta^\star_{|H|}(\theta_t)$ over the $51$ checkpoints of steps $10\,000$--$20\,000$, and $r=m_\star/\eta$; $r<1$ means that the run trains above the mean-square threshold of its own iterates.

\paragraph{Phase 5: registered, then read as exploratory.} Phase~5 (seeds 1--3) registered five hypotheses and six predictions before any threshold was computed: H1, that $\eta^\star_{|H|}(\theta_t)$ descends to $\eta$ and stays there ($r\in[0.8,1.25]$ in every run whose $\lambda_1$ has fallen below $0.8$ times that of the control), against H2, that $r<0.5$ throughout, and H3, that $\eta^\star=0$ because $H+\mu_{\mathrm{train}}I$ is indefinite, with H4 and H5 on the mean-square and almost-sure thresholds at the final point. The registration was committed before any threshold was computed; a pre-commitment file with the hashes of the registration, the evaluator, the runner, the code package and the job list was written at 23:25:29 UTC on 28 October 2026, after the 48 runs had been launched (23:15:35) and when 22 of them had finished, and the author declares that no run output had been read at that time. As computed, neither H1 nor H2 holds at either $\mu_{\mathrm{train}}$; \texttt{p4} holds and \texttt{p1}, \texttt{p2}, \texttt{p3}, \texttt{p5}, \texttt{p6} fail. An implementation check failed as well: at $(1\mp10^{-3})\eta^\star_{|H|}(\theta_T)$ the Collatz--Wielandt bracket of the second-moment operator leaves at least one side undecided in 12 of 48 runs (21 of 96 sides). The diagnosis, written after the runs, is a resolution limit: on the 9 undecided lower sides at $\mu_{\mathrm{train}}=0$, near-flat modes put the Perron eigenvalue of the operator $1.46\times10^{-9}$--$9.38\times10^{-6}$ below $1$ (by inverse iteration), and its eigenmatrix is too ill-conditioned for a double-precision test to separate the radius from $1$; on no side does a certificate contradict the side. The registration stops the reading of the predictions at a failed check, and the author decided not to read the Phase~5 predictions as outcomes. Phase~5 is reported as run and is used as the exploratory run of Phase~5b.

\paragraph{Phase 5b: a within-dataset replication.} Phase~5b repeats Phase~5 without any change to the network, the data, the training, the grid or the criteria, on three new seeds (1401--1403), and tests six predictions \texttt{q1}--\texttt{q6} written after reading Phase~5 and restricted to the cells where its three seeds agreed. The initial point and the minibatch stream depend on the seed only, so Phase~5b tests whether the structure of Phase~5 reproduces for new $(\theta_0,\{B_t\})$ streams on the same data and teacher; it is not a generalisation across datasets. A threshold read by a prediction with Lanczos relative residual above $10^{-10}$, a failed check of the almost-sure threshold or of the batch-sharpness identity, or a contradicting certificate suspends the predictions that read it, and Phase~5b counts as a replication only if all six hold and none is suspended. The registration, the evaluator and the hashes of the code package and job list were committed at 02:22 UTC on 29 October, before the first run started (02:33). All 48 runs finished, no check suspends a prediction, and the six hold (Table~\ref{tab:phase5b}). The certificate check at the final point (\texttt{c5b}, the reading of Phase~5 that accepts either certificate) fails on one run ($\mu_{\mathrm{train}}=0$, $\eta=0.4$, seed 1402: radii $1-6.3\times10^{-7}$ and $1-5.5\times10^{-7}$ at the two sides after $200\,000$ power steps), the resolution limit of Phase~5; the other 94 sides are decided, and no prediction reads that check. By the power-iteration check \texttt{c5} of Phase~5, 12 runs (21 sides) are undecided.

\begin{table}[!htbp]
\centering
\caption{Phase~5b (\texttt{raw-phase5b.json}, field \texttt{phase5b}; seeds 1401--1403, $\beta=0.9$, $B=16$). Each prediction was written before the run and holds or fails as stated, with no tolerance; the last column gives the observed value nearest to each bound. $m_\star$ and $r$ as in the text; $n_{\mathrm{sat}}$ counts hidden units with $s_h(\theta_t)<10^{-4}$; $e_{\mathrm{GN}}$ and $g_{\mathrm{GN}}$ are the window medians of $|\eta^\star_{\mathrm{GN}}/\eta^\star_{|H|}-1|$ and of $\eta^\star_{\mathrm{GN}}/\eta^\star_{|H|}$, with the Gauss--Newton per-sample Hessians in place of the full ones; $b$ is the largest window median of batch sharpness (expectation of the ratio and ratio of expectations) and of $\lambda_1$, each times $\eta/(2(1+\beta))$; $f_-$ is the fraction of the window with $\lambda_{\min}(H)<0$.}
\label{tab:phase5b}
\small
\begin{tabular}{@{}lp{0.50\textwidth}lp{0.27\textwidth}@{}}
\toprule
 & prediction (domain) & outcome & observed \\
\midrule
\texttt{q1} & $\mu_{\mathrm{train}}=10^{-3}$, every $\eta$: $m_\star\in[0.07,0.18]$, $\max_\eta m_\star/\min_\eta m_\star\le2.5$ per seed & holds & $m_\star$ $0.088$--$0.154$; ratio $\le1.69$ \\
\texttt{q2} & $\mu_{\mathrm{train}}=10^{-3}$: $r(0.1)>1.15$, $r(0.15)\in[0.80,1.15]$, $r(0.3)<0.45$, $r$ decreasing over $0.15$--$0.4$ & holds & $1.285$; $0.840$--$0.963$; $\le0.378$ \\
\texttt{q3} & $\mu_{\mathrm{train}}=0$, $\eta=0.5$: median $n_{\mathrm{sat}}\ge1$, $r\in[1.2,1.5]$; $n_{\mathrm{sat}}=0$ at $\eta\le0.2$ and median $0$ at $10^{-3}$ & holds & $n_{\mathrm{sat}}=4$, $r$ $1.356$--$1.360$ \\
\texttt{q4} & $e_{\mathrm{GN}}<0.03$ at $\eta\le0.03$, $<0.12$ at $\eta=0.1$; $g_{\mathrm{GN}}\ge1.5$ at $10^{-3}$, $\eta\ge0.3$ & holds & $0.0235$; $0.0918$; $\ge1.670$ \\
\texttt{q5} & $\eta\ge0.15$: $b<0.38$; $\eta\le0.3$: $b<0.25$ & holds & $\le0.353$; $\le0.193$ \\
\texttt{q6} & $\eta\ge0.15$: $f_-\ge0.85$ & holds & $\ge46/51=0.902$ \\
\bottomrule
\end{tabular}
\end{table}

\paragraph{What the replication measures.} With weight decay $10^{-3}$, the threshold at the iterates does not follow the learning rate: $m_\star$ stays between $0.088$ and $0.154$ over a fifty-fold range of $\eta$, while the window median of $\lambda_1$ falls from $5.82$--$5.84$ at $\eta=0.01$ to $2.36$--$2.51$ at $\eta=0.5$. The runs cross it between $\eta=0.1$ ($r\approx1.29$) and $\eta=0.15$ ($r=0.84$--$0.96$) and above it train at $2.6$--$2.9\times$ ($\eta=0.3$), $4.0$--$4.5\times$ ($\eta=0.4$) and $4.1$--$5.1\times$ ($\eta=0.5$) the $|H|$ reading of their own iterates without diverging, so H1 does not describe these runs: the linear mean-square criterion does not act as a cap that training saturates, and what bounds the iterates beyond it is outside the linearisation. Not diverging is not training well: at $\mu_{\mathrm{train}}=10^{-3}$ the window median of the data loss is $0.0025$--$0.0028$ at $\eta\le0.03$, $0.05$ at $\eta=0.2$, $0.23$--$0.37$ at $\eta=0.3$ and $1.55$--$1.71$ at $\eta=0.5$. Without weight decay, at $\eta=0.5$ four of the hidden units saturate and the threshold returns above the learning rate ($r=1.36$). The Gauss--Newton form of the criterion agrees with the full one to $2.4\,\%$ at small learning rates and overestimates it by at least $1.67\times$ in the regime above the threshold, which is where a Gauss--Newton estimate of a safe step would be used. The one-dimensional quantities stay off their deterministic edge ($b\le0.353$), so in these runs neither batch sharpness nor $\lambda_1$ locates the boundary that the operator of Theorem~\ref{thm:G6} gives, and the Hessian has a negative eigenvalue at $90$--$100\,\%$ of the window checkpoints once $\eta\ge0.15$. Three bounds held by small margins: $r(0.15)$ by $0.040$, $e_{\mathrm{GN}}$ by $0.0065$ and $b$ by $0.027$. These are measurements on one network with $33$ parameters, one dataset of $64$ samples and three new seeds; whether the same structure holds on other data and at larger $P$ is not tested here.

\subsection{A certified bracket of the Gauss--Newton threshold (Phase 6)}\label{sec:phase6}

The object is the heavy-ball iteration whose per-sample curvature is the Gauss--Newton matrix $J_i^\top J_i$ plus the tether $\mu_{\mathrm{train}}I$, at a frozen iterate $\theta_t$, restricted to the row space of the stacked Jacobian ($R=N d_{\mathrm{out}}=640$ dimensions at $P=3\,610$ and $12\,826$). Each $J_i^\top J_i$ leaves that row space invariant and vanishes on its complement, where at $\mu_{\mathrm{train}}=10^{-3}$ the iteration is a noiseless scalar heavy ball with threshold $2(1+\beta)/\mu_{\mathrm{train}}=3\,800$; the row-space threshold is therefore the threshold $\eta^\star_{\mathrm{GN}}$ of the whole Gauss--Newton-plus-tether iteration. It is a surrogate: at the same checkpoints the Hessian of the loss is indefinite and the second moment of the exact linearisation grows (measured in the pilot runs below, not reported here), and nothing below changes that.

\paragraph{The bracket.} Both sides come from Theorem~\ref{thm:G6}, $\eta^\star=1/\lambda_{\max}(G)$ with $G=L_0^{-1/2}B_\beta L_0^{-1/2}$, and are computed in double precision with every rounding error bounded \citep[Chapter~3]{higham2002accuracy}. \emph{Upper side:} any unit vector $y$ gives $\langle y,Gy\rangle\le\lambda_{\max}(G)$, so the Rayleigh quotient $q$ of the top Ritz vector of a short Lanczos run on $G$ (with full reorthogonalisation; $30$ steps, then $60$ and $120$ when the Ritz residual or the bracket is wide), minus a forward error bound $\delta_q$, gives $\eta_{\mathrm{up}}=1/(q-\delta_q)\ge\eta^\star$; no convergence is needed. \emph{Lower side:} the operator $K(\eta)$ of Theorem~\ref{thm:G6} preserves the cone $\mathbb S_+^R$, so a witness $X\succ0$ with $X-K(\eta)(X)\succ0$ gives $\rho(K(\eta))<1$ \citep[Theorem~4.2]{vandergraft1968cones} and hence $\eta<\eta^\star$. The witnesses are the positive part of $L_0^{-1/2}(y)$ plus $\varepsilon I$ for $\varepsilon=10^{-1},\dots,10^{-9}$; positive definiteness of $X-K(\eta)(X)$ in exact arithmetic is proved by a floating-point Cholesky factorisation of the computed matrix shifted by its forming error and by Rump's bound \citep{rump2006posdef}, and $\eta_{\mathrm{low}}$ is the largest rate that passes in a $50$-step bisection below $\eta_{\mathrm{up}}$. A rate at which every witness fails gives no lower bound, never a number. Checkpoints are regenerated by rerunning training, and the stored loss trace is compared bit for bit before any bracket is computed.

\paragraph{Stage 1: $P=33$, against the dense exact value.} On every window checkpoint of the $48$ Phase~5b runs (Gauss--Newton at $\mu_{\mathrm{train}}=10^{-3}$, at $\mu_{\mathrm{train}}=0$ the same object after removing the modes with $|\lambda+\mu_{\mathrm{train}}|<10^{-8}\lambda_{\mathrm{top}}$, a numerical kernel, and $|H|$; $9\,600$ brackets over all $4\,800$ checkpoints, $2\,448$ of them in the window), the regenerated runs match the stored loss traces bit for bit, no bracket excludes the dense exact threshold, and every bracket has $\eta_{\mathrm{low}}/\eta_{\mathrm{up}}\ge0.98$ (smallest $0.9934$ for Gauss--Newton, $0.980$ for $|H|$); at $\mu_{\mathrm{train}}=0$, $94$ of the stored Lanczos values of Phase~5b lie outside the bracket, all where that Lanczos had not converged. The window medians of $\eta^\star_{\mathrm{GN}}/\eta$ at $\mu_{\mathrm{train}}=10^{-3}$ fall through $1$ as $\eta$ grows, per seed $11.9$--$12.0$ at $\eta=0.01$, $1.29$--$1.30$ at $0.1$, $0.99$--$1.44$ at $0.15$, $0.85$--$0.92$ at $0.2$, $0.68$--$0.86$ at $0.3$ and $0.48$--$0.76$ at $0.4$--$0.5$: on this network the Gauss--Newton threshold does not cap training at the largest rates.

\paragraph{Stage 2: $P=3\,610$ and $12\,826$.} The checkpoints come from exploratory pilot runs of the companion lab: the two networks of Section~\ref{sec:phase2} trained as in Section~\ref{sec:phase5} (heavy-ball, $\beta=0.9$, $B=16$ with replacement, $20\,000$ steps) on datasets drawn with pilot seeds, with $\tanh$, softsign $z/(1+|z|)$ and leaky ReLU of slope $0.1$ hidden units, which recorded at each window checkpoint the threshold $\eta^\star_{\mathrm{GN},64}$ of the Gauss--Newton iteration restricted to its top $64$ modes, an upper bound of $\eta^\star_{\mathrm{GN}}$ because a restriction to a subspace can only raise the threshold. Seed~$0$ of every pilot cell with $\mu_{\mathrm{train}}=10^{-3}$ and a finite run ($30$ runs; $\eta$ from $0.001$ to $0.1$ at $P=3\,610$ and to $0.05$ at $P=12\,826$, leaky ReLU only up to $0.01$), at the $11$ window checkpoints $10\,000,11\,000,\dots,20\,000$. The checkpoint list was committed before the stage~2 batch was run; the reference is the same certified bracket with Lanczos run to a target Ritz residual of $10^{-10}$ (reached at $232$ of the $330$ checkpoints within $200$ steps; largest residual $1.6\times10^{-6}$). All $30$ regenerations matched. At all $330$ checkpoints the bracket and the reference overlap with no tolerance, the two upper bounds agree to $7\times10^{-6}$ relative, and $\eta_{\mathrm{low}}/\eta_{\mathrm{up}}\ge0.982$ (median $0.9994$); the row space has full rank $640$ at every checkpoint. The top-$64$ compression of the pilot overstated the threshold by $1.7$--$8.9\,\%$ ($\eta_{\mathrm{up}}/\eta^\star_{\mathrm{GN},64}=0.919$--$0.983$).

Reading each checkpoint against the band $[0.8\,\eta,1.25\,\eta]$ of the Phase~5 registration (\emph{inside} when $0.8\,\eta\le\eta_{\mathrm{low}}$ and $\eta_{\mathrm{up}}\le1.25\,\eta$, \emph{below} when $\eta_{\mathrm{up}}<0.8\,\eta$, \emph{above} when $\eta_{\mathrm{low}}>1.25\,\eta$): the $12$ runs at $\eta\le0.003$ read above at every checkpoint ($\eta_{\mathrm{up}}/\eta$ between $1.6$ and $27$), as expected near a minimum. Of the $18$ runs at $\eta\ge0.01$, $16$ read inside at most of their checkpoints, with window medians of $\eta^\star_{\mathrm{GN}}/\eta$ between $0.81$ and $1.10$; the other two, $\tanh$ and softsign at $P=3\,610$, $\eta=0.01$, read above ($1.46$ and $1.81$). Over these $198$ checkpoints the readings are $161$ inside, $23$ above and $14$ below, and no run reads below at more than $4$ of its $11$ checkpoints; $13$ of the $30$ runs read inside at $80\,\%$ or more of their checkpoints. Unlike $P=33$, where the ratio keeps falling to $0.48$--$0.76$, on these two networks the window median of $\eta^\star_{\mathrm{GN}}/\eta$ stays within the band ($0.81$--$1.10$) at every finite rate from $\eta=0.02$ upwards, for both bounded activations; single checkpoints go down to $0.46$.

\paragraph{What Phase 6 answers.} Three answers that the earlier phases did not give, and that we have not found in the cited literature. (i) The exact mean-square threshold of heavy-ball SGD with per-sample curvatures that do not commute can be bracketed from both sides, with every rounding error bounded for the computed Gauss--Newton factors, at training checkpoints of networks with thousands of parameters and without forming a Hessian: at $R=640$ the operator $G$ acts on a space of dimension $205\,120$, and all $330$ brackets have relative width below $1.8\,\%$ (median $0.06\,\%$). (ii) For the Gauss--Newton surrogate, the saturation half of the open question of \citet{andreyev2026momentum} has a measured, exploratory answer that depends on the network (Theorem~\ref{thm:slow-fast-saturation} gives a model in which both outcomes occur). At $P=33$ training goes past the threshold at the largest rates (window medians $0.48$--$0.76$); at $P=3\,610$ and $12\,826$, for both bounded activations and every finite rate from $\eta=0.02$, the window median of $\eta^\star_{\mathrm{GN}}/\eta$ lies in $[0.81,1.10]$, so training runs at the certified threshold of the surrogate, a quantity the pilot bounded only from above, by $\eta^\star_{\mathrm{GN},64}$. At the smallest rates ($\eta\le0.003$) the threshold sits $1.6$--$27$ times above $\eta$: saturation is reached only once $\eta$ is large enough. (iii) At training checkpoints the top-$64$ compression overstates the threshold by $1.7$--$8.9\,\%$, a gap that had not been measured. None of the three is a statement about the loss linearisation, whose second moment grows at the same checkpoints.

\paragraph{Status.} The bracket is a measurement of the surrogate at a checkpoint, certified for the stored factors and $\mathrm{fl}(\lambda_k+\mu_{\mathrm{train}})$; the distance from those factors to the Gauss--Newton operator is reported beside it (relative residuals below $9\times10^{-10}$ and $1.4\times10^{-14}$). Stages~1 and~2 are exploratory and read only pilot checkpoints and the Phase~5b runs, one seed per cell at the larger sizes; the expectations written before stage~2 are not registered predictions, and they are reported as they came out: the top-$64$ compression within $5\,\%$ ($\eta_{\mathrm{up}}/\eta^\star_{\mathrm{GN},64}\ge0.95$) at $90\,\%$ of the checkpoints holds narrowly ($90.6\,\%$); every run with $\eta\ge0.01$ inside at most of its checkpoints fails, on the two $P=3\,610$ runs at $\eta=0.01$; no run below at most of its checkpoints holds. Stage~2 passes the validity and adoption rules of the plan (every bracket overlaps its reference with no tolerance, $\eta_{\mathrm{low}}/\eta_{\mathrm{up}}\ge0.98$ at every checkpoint), so the larger network ($P=34\,186$) may be bracketed next. The band question is left to a registration on unseen seeds and datasets. A larger network ($P=34\,186$, $R=2\,560$) has not been bracketed.


\clearpage
\section{Scope and limitations}\label{sec:limits}

All statements are about the linearised dynamics around a fixed point of a quadratic model: i.i.d.\ minibatch Hessians with mean $H$, additive noise independent of the past, and stationarity. The criterion is a mean-square one; between the mean-square and the almost-sure thresholds of Proposition~\ref{prop:G3} the second moment diverges while typical trajectories decay. Not answered here: the dynamics beyond the quadratic model, which is what bounds the training runs of Section~\ref{sec:phase5} above the criterion, and a threshold on the batch sharpness of Definition~3.1 of \citet{andreyev2026momentum} alone: Theorem~\ref{thm:modal-bridge} gives an identity at the threshold and a sufficient criterion that also uses the noise share $\theta_{\mathrm{noise}}$, and Theorem~\ref{thm:slow-fast-saturation} proves saturation in a slow-fast model, not for the training dynamics. On the two networks of Section~\ref{sec:phase2} the random-feature model of Theorem~\ref{thm:G} does not hold (\texttt{q3} fails by $1.8\times$ and $63\times$), and the untethered network iteration has no mean-square stable learning rate, because the mean Hessian is indefinite at both checkpoints; its growth within $1\,000$ steps changes sign between $0.01\times$ and $0.02\times$ the deterministic threshold at $P=3\,610$ and below $0.002\times$ at $P=12\,826$. Theorem~\ref{thm:G6} was evaluated on network Hessians only with a tether that lifts the negative eigenvalues of the mean Hessian (Section~\ref{sec:phase4}); at $P=33$ the threshold is enclosed in all 16 cells by the certified bracket of Section~\ref{sec:phase6} (a post-hoc run), for the computed per-sample Hessians; at $P=3\,610$ the crossing rests on power iteration alone. Section~\ref{sec:phase5} evaluates it at training iterates with $|H|$ in place of an indefinite $H$, as a quasi-static diagnostic of a moving point, on one network with $33$ parameters and one dataset of $64$ samples; it measures a replication on new seeds, not a generalisation across data or size. Section~\ref{sec:phase6} brackets the threshold of the Gauss--Newton surrogate, not of the loss linearisation, whose second moment grows at the same checkpoints; it is exploratory, with one seed per cell at the larger sizes.

\section{Related work}

\citet{cohen2021eos} described the edge of stability for full-batch training. \citet{andreyev2024eoss} described its stochastic counterpart, in which the batch sharpness rather than the full-batch sharpness stabilizes, and \citet{andreyev2026momentum} extended it to momentum; their one-dimensional mean-square threshold is Proposition~\ref{prop:oneD}. The linear-stability account of \citet{wu2018sgd} contains its $\beta=0$ case; see also the break-even analysis of \citet{jastrzebski2020breakeven}. The exact second-moment equation of Section~\ref{sec:lyapunov} sits in the line of stability analyses of SGD with momentum on quadratics: \citet{kidambi2018insufficiency} and \citet{liu2020mass} with component-decoupled data, \citet{ma2018power} for the batch-size scaling of SGD without momentum, \citet{mulayoff2024exact} for the exact mean-square stability threshold of SGD without momentum on quadratic losses, \citet{paquette2021momentum} (batch size $1$, through homogenized SGD) and \citet{lee2022minibatch} (mini-batches) in high-dimensional limits through Volterra equations, with the exact high-dimensional dynamics of SGD on least squares by homogenization \citep{paquette2022homogenization}, and the power-law random features model analysed for one-pass SGD \citep{paquette2024phases} and for stochastic momentum at small batch \citep{ferbach2025dana}, where SGD with fixed momentum has the scaling-law exponents of SGD and dimension-adapted Nesterov acceleration improves them, and \citet{gitman2019momentum} for the stationary covariance with additive noise, and \citet{yuan2016momentum} for the equivalence of momentum SGD with SGD at step $\eta/(1-\beta)$ for small constant steps. At $\beta=0$, $B=1$ and Gaussian features, Theorem~\ref{thm:G} is the mean-square condition of the LMS filter with independent Gaussian inputs \citep{horowitz1981performance,feuer1985convergence}, whose second-moment recursion has the same diagonal closure and rank-one secular structure; the Kronecker form of that recursion is standard in adaptive filtering \citep{sayed2008adaptive}. For momentum LMS with independent Gaussian inputs, which is Theorem~\ref{thm:G} at $B=1$ and $\kappa_{\mathrm{feat}}=0$ with $\eta=2\mu$ and $\beta=\alpha$ in their notation, \citet{roy1990momentum} close the diagonal second-moment recursion (their (A6a)) and obtain from one Schur--Cohn test, $D(1)>0$, a necessary condition for mean-square convergence, $Q(1)>0$ (their (31)--(33)); evaluated from their (A10), $Q(1)>0$ reads $\sum_k\mu\lambda_k(1+\alpha)/(1-\alpha^2-2\mu\lambda_k)<1$, which is \eqref{eq:secular} in that case (their printed stability condition~3 omits the factor $1+\alpha$, which their misadjustment (36) carries). Theorem~\ref{thm:G} shows that for $\alpha\in[0,1)$ and $2\mu\lambda_k<1-\alpha^2$ this condition is also sufficient, and extends it to minibatches and to non-Gaussian independent features. \citet{tugay1989momentum} also analysed momentum LMS, and \citet{sharma1998momentum} give small-step asymptotics (an ODE limit and the asymptotic distribution) for every $\alpha\in(-1,1)$. \citet{jin2026momentum} prove $L_p$ exponential stability and tracking and regret bounds for a normalized momentum LMS with dependent, nonstationary regressors, under a sufficient bound on the step size and a momentum $\beta\le C_\beta\mu$ that vanishes with the step; they do not give a threshold. Theorem~\ref{thm:G} has the form of the convergence condition of \citet{velikanov2023view}, made exact for independent-coordinate features (with their parameter $\tau$ outside the range $0<\tau\le1$ of their convergence analysis for $\kappa_{\mathrm{feat}}\ge-1$, Gaussian features included, and their infinite-dimensional $H$ replaced by a finite $P$), and is a finite-dimensional counterpart, with a kurtosis term, of the kernel-norm condition that \citet{lee2022minibatch} state in a high-dimensional limit. Theorem~\ref{thm:F} extends the exact mean-square analysis of SGD by \citet{mulayoff2024exact}, whose form it takes at $\beta=0$ with sampling with replacement, to heavy-ball momentum with non-commuting curvature noise; the threshold of Theorem~\ref{thm:G6} is also Theorem~F.1 of \citet{cai2026llm}, obtained concurrently and independently; \citet{mulayoff2026nonlinear} show that beyond linearization a single unstable batch can make SGD diverge in expectation. Theorem~\ref{thm:F} gives the stability boundary of the Schur-complement expansion of \citet[Appendix~D]{andreyev2026momentum} exactly, and Corollary~\ref{cor:msinstab} gives, for $\rho(\mathcal T)\ge1$ and $\Sigma_\xi\succ0$, linear growth of the second moment; its probabilistic conclusion has the form of their Definition~A.1 but is also met below $\rho(\mathcal T)=1$, so it does not settle the valid instability criterion they list as open for SGD with momentum. The secular structure of Theorem~\ref{thm:G} (cone-preserving operator, rank-one coupling, Krein--Rutman) is also that of \citet{song2026zeroth} for full-batch zeroth-order methods, whose randomness is the search direction rather than the minibatch; \citet{sharma2026outer} give the exact mean-square criterion of outer Nesterov momentum on stochastic pseudo-gradients, which at one local step covers SGD with Nesterov momentum.

\section{Conclusion}\label{sec:conclusion}

For heavy-ball SGD on a quadratic with i.i.d.\ minibatch Hessians, the mean-square stability boundary is decided at every $\eta>0$ and $\beta\in[0,1)$ by the stationary equation of Theorem~\ref{thm:F}, without a small-$\eta$, noise-dominated or commutation assumption, and beyond it the second moment diverges (Corollary~\ref{cor:msinstab}). For random-feature minibatches the boundary is the root of the secular equation of Theorem~\ref{thm:G}, whose two limits take the values of the two batch-sharpness plateaus, through $\lambda_1$ and $\tr H/B$ (Corollary~\ref{cor:plateaus}); for arbitrary per-sample Hessians it is the reciprocal of the top eigenvalue of the operator of Theorem~\ref{thm:G6}; and when the curvature noise is not proportional to the curvature, a nearly flat mode with curvature noise lowers it without bound (Theorem~\ref{thm:G5}). The secular root matches the spectral radius of the exact second-moment operator within $6.11\times10^{-10}$ relative in 36 cells and the sample-path ceiling in 29 of 30 cells; on the two network checkpoints the random-feature model does not hold, and the untethered linearised iteration, with an indefinite mean Hessian, has no mean-square stable learning rate. With a tether that lifts the negative eigenvalues of the mean Hessian, the threshold of Theorem~\ref{thm:G6} computed from the real per-sample Hessians is where the spectral radius of the exact operator crosses $1$, within a relative $\pm10^{-3}$ in $\eta$, and it lies $1.23\times$ below the one-dimensional ceiling at $P=3\,610$. At the iterates of heavy-ball training of a $33$-parameter network with weight decay the threshold does not follow the learning rate; the runs cross its $|H|$ reading between $\eta=0.1$ and $0.15$ and train up to $5.1\times$ above it without diverging, so on that network the linear mean-square criterion is not a cap that training saturates. At the threshold an exact identity ties the criterion to the two batch-sharpness plateaus, and in a quasi-stationary slow-fast model saturation occurs at a distance $\Theta(\sigma_\xi^2/v_{\mathrm{PS}})$ below the threshold, while a secant equilibrium lies above the tangent one (Theorems~\ref{thm:modal-bridge} and~\ref{thm:slow-fast-saturation}). The threshold of Theorem~\ref{thm:G6} can be bracketed from both sides with every rounding error bounded for the computed factors and without forming a Hessian: on the Gauss--Newton surrogate of networks with $3\,610$ and $12\,826$ parameters ($205\,120$-dimensional operator) all $330$ brackets have relative width below $1.8\,\%$, and there, unlike at $P=33$, the window medians of the threshold stay at $0.81$--$1.10$ times the learning rate from $\eta=0.02$ upwards.

Open questions: (i) whether the structure of Section~\ref{sec:phase5} holds on other data and at larger $P$, and what bounds the iterates above the criterion; (ii) the dynamics beyond the quadratic model, where the loss is bounded below and the mean gradient is not zero, which is where the excursions of training at practical learning rates come from (Section~\ref{sec:phase2}); (iii) a threshold on the batch sharpness of \citet[Definition~3.1]{andreyev2026momentum} alone, and saturation for the training dynamics rather than for a slow-fast model; (iv) the almost-sure threshold for $\beta\in(0,1)$ with unbounded curvature noise, where $\det M_t=\beta$ removes the singular multiplier of $\beta=0$ and the proof of Proposition~\ref{prop:G3} lacks only the return of its projective coordinate after large excursions (heuristically, the steps back after a near-singular multiplier $v$ accumulate $\log v$, as the singular integral does at $\beta=0$); (v) a measurement of the mean-square boundary on networks with an estimator of $\E\|w_t\|^2$, such as the resampling estimator of Section~\ref{sec:phase1}, instead of sample paths, whose threshold can lie above the mean-square one where the multiplicative noise along $u_1$ dominates (Section~\ref{sec:phase1}). Adaptive (state-dependent) preconditioners and orthogonalised updates are not covered; a fixed preconditioner reduces to the theorems (Section~\ref{sec:lyapunov}, after Theorem~\ref{thm:G6}).

\section*{Reproducibility}
\begin{sloppypar}
Code, raw files and tests are in the repository that accompanies this draft, \url{https://github.com/apis-os/hb-sgd-mean-square}, lab \nolinkurl{labs/mean-square-boundary}, whose \nolinkurl{README.md} maps the theorem letters used in the code to the numbering of this draft. Phase~1: \nolinkurl{solve-secular-eoss-threshold.ts}, \nolinkurl{simulate-heavy-ball-feature-noise-energy.ts}, \nolinkurl{measure-empirical-eoss-ceiling.ts}, \nolinkurl{compute-symmetric-eigenvalues-jacobi.ts}, \nolinkurl{evaluate-phase1-rademacher-cells.ts}, \nolinkurl{run-phase1.ts} (8 seeds per cell, Node 26.5, 43 min on one thread), and \nolinkurl{run-phase1-cap-diagnostic.ts} for the post-hoc diagnostic of its failed cell. Phase~2: \nolinkurl{compute-exact-hessian-trace.ts}, \nolinkurl{compute-top-eigenpair-power-iteration.ts}, \nolinkurl{compute-secular-ceiling-bracket-from-trace.ts}, \nolinkurl{simulate-network-minibatch-heavy-ball-energy.ts}, \nolinkurl{measure-network-minibatch-heavy-ball-ceiling.ts}, \nolinkurl{evaluate-phase2-trace-ceiling-cells.ts}, \nolinkurl{run-phase2.ts}. Phase~3: \nolinkurl{run-phase3.ts}. Phase~4: \nolinkurl{run-phase4.ts}, \nolinkurl{evaluate-phase4-tethered-mean-square-cells.ts}, \nolinkurl{create-phase4-sample-hessian-provider.ts}, and \nolinkurl{run-phase4-certified-bracket-p33.ts} with its output \nolinkurl{phase4-certified-bracket-p33.json} for the certified brackets at $P=33$. Phase~5: \nolinkurl{run-phase5.ts} (one run per process, 48 runs, one CPU each), \nolinkurl{run-phase5-aggregate.ts}, \nolinkurl{run-phase5-c5-perron.ts} for the certificate diagnosis; Phase~5b: \nolinkurl{evaluate-phase5b-predictions.ts}, run records in \nolinkurl{evidence/mean-square-boundary/phase5b-runs/} and the job list \nolinkurl{phase5b-jobs.json}; the registration texts, the pre-commitments and package manifests of Phases~5 and 5b and their commits are in the companion lab, not in this repository. Phase~6: plan \nolinkurl{phase6-plan-certified-mean-square-bracket.md}, \nolinkurl{compute-certified-rayleigh-upper-threshold.ts}, \nolinkurl{certify-cone-collatz-wielandt-at-rate.ts}, \nolinkurl{verify-positive-definite-shifted-cholesky.ts}, \nolinkurl{compute-adaptive-certified-threshold-bracket.ts}, \nolinkurl{run-phase6-stage1.ts}, \nolinkurl{run-phase6-stage2.ts}, \nolinkurl{summarize-phase6-stage1.ts}, \nolinkurl{summarize-phase6-stage2.ts}, the checkpoint list \nolinkurl{phase6-checkpoints.json} (committed before the runs), raw records in \nolinkurl{phase6-stage1/} and \nolinkurl{phase6-stage2/}, the pilot records \nolinkurl{phase6-pilot/} that stage~2 regenerates, and the summaries \nolinkurl{phase6-stage1-summary.json} and \nolinkurl{phase6-stage2-summary.json}. The Phase~4 row workers use a WebAssembly \nolinkurl{f64x2} product kernel shipped as \nolinkurl{dense-f64-matmul.wasm} and checked bit for bit against the sequential product; the increment-stream block of the Phase~5 checkpoint records belongs to the companion manuscript and is not computed in the repository. Tests \nolinkurl{tests/phase1-secular-ceiling.test.ts}, \nolinkurl{tests/phase2-trace-ceiling.test.ts}, \nolinkurl{tests/phase3-operator-radius.test.ts}, \nolinkurl{tests/phase4-tethered-thm-g6.test.ts}, \nolinkurl{tests/phase5b-evaluator.test.ts}, \nolinkurl{tests/phase6-certified-bracket.test.ts}, \nolinkurl{tests/phase6-stage2-summary.test.ts} and \nolinkurl{tests/phase4-certified-bracket-p33.test.ts} re-execute reduced versions of Phases~1--3 against \nolinkurl{raw-phase1.json}, \nolinkurl{raw-phase2.json} and \nolinkurl{raw-phase3.json}, including the failed predictions; check the Phase~4 kernels and a resumed row, the Phase~5b evaluator on the Phase~5b run records, the regeneration of Phase~6 checkpoints from those runs and from the pilot records, the Phase~6 stage~2 summary, and the certified brackets at $P=33$ against \nolinkurl{raw-phase4.json}; \nolinkurl{tests/phase4-5-raw-files.test.ts} recomputes the predictions and checks stored in \nolinkurl{raw-phase4.json} from its cell records and the decisions stored in \nolinkurl{raw-phase5.json} and \nolinkurl{raw-phase5b.json} from the 96 run records. The identities and limits of Theorem~\ref{thm:modal-bridge} and the saturation law of Theorem~\ref{thm:slow-fast-saturation} are checked in \nolinkurl{paper/checks/bridge-theorems.py}. The numerical checks cited in the text are in \nolinkurl{paper/checks/}, \nolinkurl{paper/check-schur-factor.py} and the lab's \nolinkurl{check-phase3-convergence.ts}.
\end{sloppypar}

\bibliographystyle{plainnat}
\bibliography{refs}

\appendix
\section{Proofs of the results of Section~\ref{sec:lyapunov}}\label{app:thmFG}

\begin{proof}[Proof of Theorem~\ref{thm:F}]
Let $A=(1+\beta)I-\eta H$ and $\Sigma_1=\E[w_{t+1}w_t^\top]$. (1) Since $(H_{B_t}-H,\xi_t)$ is independent of $\mathcal F_{t-1}$ with zero mean, $\E[w_{t+1}\mid\mathcal F_{t-1}]=Aw_t-\beta w_{t-1}$; multiplying on the right by $w_t^\top$, $\Sigma_1=A\Sigma_0-\beta\Sigma_1^\top$, and transposing ($A$ symmetric), $\Sigma_1^\top=\Sigma_0A-\beta\Sigma_1$. Hence $\Sigma_1=\frac{A\Sigma_0-\beta\Sigma_0A}{1-\beta^2}$, $\Sigma_1+\Sigma_1^\top=\frac{A\Sigma_0+\Sigma_0A}{1+\beta}=2\Sigma_0-\frac{\eta}{1+\beta}(H\Sigma_0+\Sigma_0H)$, and $D_\Theta=2\Sigma_0-\Sigma_1-\Sigma_1^\top$. (2) Expand $\Sigma_0=\E[w_{t+1}w_{t+1}^\top]$ with $w_{t+1}=Aw_t-\beta w_{t-1}-\eta(H_{B_t}-H)w_t-\eta\xi_t$; the cross terms vanish because $(H_{B_t}-H,\xi_t)$ is independent of $(w_t,w_{t-1})$ and $\E[w_t]=0$ (this covers $H_{B_t}$ and $\xi_t$ from the same minibatch), leaving $\Sigma_0=A\Sigma_0A+\beta^2\Sigma_0-\beta(A\Sigma_1+\Sigma_1^\top A)+\eta^2\mathcal V(\Sigma_0)+\eta^2\Sigma_\xi$. With $\Sigma_1$ from (1) and $A^2\Sigma_0+\Sigma_0A^2-2\beta A\Sigma_0A=[A,[A,\Sigma_0]]+2(1-\beta)A\Sigma_0A$, $-\beta(A\Sigma_1+\Sigma_1^\top A)=-\frac{\beta}{1-\beta^2}[A,[A,\Sigma_0]]-\frac{2\beta}{1+\beta}A\Sigma_0A$, and $[A,[A,\Sigma_0]]=\eta^2[H,[H,\Sigma_0]]$. Collecting, $(1-\beta^2)\Sigma_0-\frac{1-\beta}{1+\beta}A\Sigma_0A=-\frac{\beta\eta^2}{1-\beta^2}[H,[H,\Sigma_0]]+\eta^2\mathcal V(\Sigma_0)+\eta^2\Sigma_\xi$; expanding $A\Sigma_0A=(1+\beta)^2\Sigma_0-\eta(1+\beta)(H\Sigma_0+\Sigma_0H)+\eta^2H\Sigma_0H$ the left side is $\eta(1-\beta)(H\Sigma_0+\Sigma_0H)-\eta^2\frac{1-\beta}{1+\beta}H\Sigma_0H$, and dividing by $\eta(1-\beta)$ gives (2). (3) Sandwich (2) between $u_j^\top$ and $u_k$ and substitute $(\Sigma_0)_{jk}=\frac{\eta(1+\beta)}{\lambda_j+\lambda_k}C_{jk}$ from (1); at $j=k$ the bracket is $1-\frac{\eta\lambda_k}{2(1+\beta)}$. For the comparison with the Schur complement of \citet{andreyev2026momentum}: in the eigenbasis of $H$, the blocks $M_{12}$, $M_{22}$ and the mean part of $M_{21}$ act on each entry $(j,k)$ separately through $A_j\otimes A_k$, $A_j=\begin{pmatrix}1-\eta\lambda_j&-\eta\beta\\\lambda_j&\beta\end{pmatrix}$, and the curvature noise enters only $M_{11}$, as $\eta^2\mathcal V$, and $M_{21}$, as $(-\eta\mathcal V,-\eta\mathcal V,\mathcal V)$ in the $xv$, $vx$, $vv$ rows. Rescaling the $vv$ row and column and adding the third row to the first two gives $\det(I-M_{22})_{jk}=(1-\beta)\big((1-\beta)^2(1+\beta)+\beta\eta(\lambda_j+\lambda_k)\big)>0$. With $\mathcal S_{jk}=\frac{(1-\beta)^2(1+\beta)}{(1-\beta)^2(1+\beta)+\beta\eta(\lambda_j+\lambda_k)}$ and $L_{jk}$ the $(j,k)$ coefficient of $\mathcal L$ at $\mathcal V=0$, the identities $\det(I-A_j\otimes A_k)=(1-\beta)^3(1+\beta)L_{jk}$ and $\eta^2+M_{12}(I-M_{22})^{-1}(-\eta,-\eta,1)^\top=\mathcal S_{jk}\eta^2/(1-\beta)^2$ are rational identities in $(\eta,\beta,\lambda_j,\lambda_k)$; each entry of $A_j\otimes A_k$ has degree at most $(2,2,1,1)$ in $(\eta,\beta,\lambda_j,\lambda_k)$, so $\det(I-M_{22})_{jk}$ has degree at most $(6,6,3,3)$ and $\operatorname{adj}(I-M_{22})_{jk}$ at most $(4,4,2,2)$; with $s_{jk}=(1-\beta)^2(1+\beta)+\beta\eta(\lambda_j+\lambda_k)$, the script checks $\det(I-M_{22})_{jk}=(1-\beta)s_{jk}$, the first identity in its Schur form $\det(I-M_{22})_{jk}(1-M_{11})-M_{12}\operatorname{adj}(I-M_{22})M_{21}=(1-\beta)^2(1+\beta)\det(I-M_{22})_{jk}L_{jk}/s_{jk}$, and the second identity; multiplied by $s_{jk}(1-\beta)^2(1-\beta^2)$ and by $s_{jk}\det(I-M_{22})_{jk}$ they become polynomial identities of degree at most $(9,15,5,5)$ and $(9,9,4,4)$. They hold in exact rational arithmetic on a $13\times22\times13\times13$ product grid (\texttt{paper/check-schur-factor.py}), which has more points than the degree in each variable, so they hold identically. With the Schur determinant formula $\det(I-A_j\otimes A_k)=\det(I-M_{22})_{jk}(1-T_{\mathrm{slow}})_{jk}$ they give $I-T_{\mathrm{slow}}=\mathcal S\circ\mathcal L$ entry by entry, the noise term included; the same script checks the operator identity with six non-commuting per-sample Hessians at $P=2$ (residual $8\times10^{-16}$) and the leading term $R=\frac{2\beta\eta^2}{(1-\beta)^3}H\Sigma H$ of the remainder on pairs with $\lambda_k/\lambda_j\in\{1/3,10^{-2},10^{-3},0\}$, where it separates from the double-commutator reading (ratio $0.999995$ at $\eta=10^{-7}$, $\beta=0.9$). \end{proof}

\begin{lemma}[The remainder of the Schur-complement expansion]\label{lem:remainder}
In the setting of Theorem~\ref{thm:F} with $H\succeq0$, let $T_{\mathrm{slow}}(\eta,\beta)$ be the Schur complement of \citet[Appendix~D]{andreyev2026momentum}, $\eta_{\mathrm{eff}}=\eta/(1-\beta)$, $K_H(X)=HX+XH$, $G(X)=\E[H_{B_t}XH_{B_t}]=HXH+\mathcal V(X)$, and define their remainder by $T_{\mathrm{slow}}=I-\eta_{\mathrm{eff}}K_H+\eta_{\mathrm{eff}}^2G+R$. Write $\mathcal L=\eta_{\mathrm{eff}}K_H+\mathcal L_2$ for $\eta/(1-\beta)$ times the left side of Theorem~\ref{thm:F}(2), so that $\mathcal L_2(X)=-\frac{\eta^2}{1-\beta^2}HXH+\frac{\eta^2\beta}{(1-\beta)^2(1-\beta^2)}[H,[H,X]]-\eta_{\mathrm{eff}}^2\mathcal V(X)$, and $c=\frac{\beta\eta}{(1-\beta)^2(1+\beta)}$. Then, exactly,
\[
R(\Sigma)=\frac{2\beta\eta^2}{(1-\beta)^3}H\Sigma H+R_3(\Sigma),\qquad R_3=c\,K_H\circ\mathcal L_2-c^2K_H^2\circ\mathcal S\circ\mathcal L,
\]
and, in the Frobenius and in the spectral norm, $\|R_3(\Sigma)\|\le2c\|H\|\,\|\mathcal L_2(\Sigma)\|+4c^2\|H\|^2\|\mathcal L(\Sigma)\|$. The order is not uniform in $\beta$: $R_3=O\big(\beta\eta^3(\|H\|^3+\|H\|\|\mathcal V\|)/(1-\beta)^4+\beta^2\eta^3\|H\|^3/(1-\beta)^5\big)$, so $R_3$ is small against the leading term when $\eta\|H\|\ll1-\beta$, $\eta\|\mathcal V\|\ll(1-\beta)\|H\|$ and $\beta\eta\|H\|\ll(1-\beta)^2$, and it is not small when $\beta\eta\|H\|\gtrsim(1-\beta)^2$. If $\lambda_k=0$ and $\Sigma$ is symmetric, the $(j,k)$ entry is $R(\Sigma)_{jk}=\eta_{\mathrm{eff}}^2(\mathcal S_{jk}-1)\mathcal V(\Sigma)_{jk}$ for every $\eta$; it vanishes when $\mathcal V(\Sigma)_{jk}=0$ (no curvature noise, or $\Sigma_z=H$ in Theorem~\ref{thm:G5}, where $\mathcal V(\Sigma)_{jk}=\lambda_j\lambda_k\Sigma_{jk}/B$), but not in the flat-mode setting of Theorem~\ref{thm:G5} with $s_k>0$, where $\mathcal V(\Sigma)_{jk}=s_js_k\Sigma_{jk}/B$.
\end{lemma}
\begin{proof}
By the identity $I-T_{\mathrm{slow}}=\mathcal S\circ\mathcal L$ with $\mathcal S^{-1}=I+cK_H$, $\mathcal S=I-cK_H\circ\mathcal S$, so $\mathcal S\circ\mathcal L=\mathcal L-cK_H\mathcal L+c^2K_H^2\mathcal S\mathcal L=\eta_{\mathrm{eff}}K_H+\mathcal L_2-c\eta_{\mathrm{eff}}K_H^2-cK_H\mathcal L_2+c^2K_H^2\mathcal S\mathcal L$. Hence $R=T_{\mathrm{slow}}-I+\eta_{\mathrm{eff}}K_H-\eta_{\mathrm{eff}}^2G=-\mathcal L_2-\eta_{\mathrm{eff}}^2(H\Sigma H+\mathcal V)+c\eta_{\mathrm{eff}}K_H^2+R_3$. The $\mathcal V$ terms cancel. With $a=c\,\eta_{\mathrm{eff}}=\frac{\eta^2\beta}{(1-\beta)^2(1-\beta^2)}$, the coefficient of the double commutator in $\mathcal L_2$, and $K_H^2(\Sigma)-[H,[H,\Sigma]]=4H\Sigma H$, the remaining terms are $\big(\frac{\eta^2}{1-\beta^2}-\eta_{\mathrm{eff}}^2+4a\big)H\Sigma H=\frac{2\beta\eta^2}{(1-\beta)^3}H\Sigma H$. For the bound, $\|K_H\|\le2\|H\|$ in both norms. In the eigenbasis $\mathcal S$ multiplies the $(j,k)$ entry by $\mathcal S_{jk}=1/(1+c\lambda_j+c\lambda_k)=1/(a_j+a_k)$ with $a_j=\frac12+c\lambda_j>0$ for $H\succeq0$: these multipliers lie in $(0,1]$, so $\|\mathcal S\|\le1$ in the Frobenius norm, and they form a positive semidefinite Cauchy matrix with diagonal at most $1$, so by Schur's theorem on Hadamard multipliers $\mathcal S$ is also a contraction in the spectral norm. The orders follow from $c=O(\beta\eta/(1-\beta)^2)$, $\|\mathcal L\|=O(\eta\|H\|/(1-\beta))$ and $\|\mathcal L_2\|=O\big(\eta^2(\|H\|^2+\|\mathcal V\|)/(1-\beta)^2+\beta\eta^2\|H\|^2/(1-\beta)^3\big)$ (the term $-\frac{\eta^2}{1-\beta^2}HXH$ and the curvature noise carry no factor $\beta$, the double commutator carries the largest power of $(1-\beta)^{-1}$), which give $c\|H\|\|\mathcal L_2\|=O\big(\beta\eta^3\|H\|(\|H\|^2+\|\mathcal V\|)/(1-\beta)^4+\beta^2\eta^3\|H\|^3/(1-\beta)^5\big)$ against the leading $2\beta\eta^2\|H\|^2/(1-\beta)^3$, a ratio of order $\eta\|H\|/(1-\beta)+\eta\|\mathcal V\|/((1-\beta)\|H\|)+\beta\eta\|H\|/(1-\beta)^2$. For the last claim, at $\lambda_k=0$ the leading term $\lambda_j\lambda_k\Sigma_{jk}$ is zero, $(K_HX)_{jk}=\lambda_jX_{jk}$, $\mathcal S_{jk}=1/(1+c\lambda_j)$, $\mathcal L_2(\Sigma)_{jk}=a\lambda_j^2\Sigma_{jk}-\eta_{\mathrm{eff}}^2\mathcal V(\Sigma)_{jk}$ and $\mathcal L(\Sigma)_{jk}=\eta_{\mathrm{eff}}\lambda_j(1+c\lambda_j)\Sigma_{jk}-\eta_{\mathrm{eff}}^2\mathcal V(\Sigma)_{jk}$ (using $a=c\eta_{\mathrm{eff}}$), so $(R_3)_{jk}=c\lambda_j\big(a\lambda_j^2\Sigma_{jk}-\eta_{\mathrm{eff}}^2\mathcal V_{jk}\big)-\frac{c^2\lambda_j^2}{1+c\lambda_j}\big(\eta_{\mathrm{eff}}\lambda_j(1+c\lambda_j)\Sigma_{jk}-\eta_{\mathrm{eff}}^2\mathcal V_{jk}\big)=-\eta_{\mathrm{eff}}^2\frac{c\lambda_j}{1+c\lambda_j}\mathcal V_{jk}=\eta_{\mathrm{eff}}^2(\mathcal S_{jk}-1)\mathcal V(\Sigma)_{jk}$; the $\Sigma_{jk}$ terms cancel because $a=c\eta_{\mathrm{eff}}$.
\end{proof}

\begin{proof}[Proof of Theorem~\ref{thm:G}]
\emph{Entries of $\mathcal V$.} In the eigenbasis $(H_{B_t}-H)_{lm}=\frac{\sqrt{\lambda_l\lambda_m}}{B}\sum_b(\omega_{b,l}\omega_{b,m}-\delta_{lm})$ and $(\mathcal V(X))_{jk}=\sum_{l,m}\E[(H_{B_t}-H)_{jl}(H_{B_t}-H)_{mk}]X_{lm}$. With independent zero-mean unit-variance coordinates, $\Cov(\omega_j\omega_l,\omega_m\omega_k)\ne0$ only when $\{j,l\}=\{m,k\}$ as multisets: for $j=k$ it is $1$ when $l=m\ne k$ and $2+\kappa_{\mathrm{feat}}$ when $l=m=k$, so $(\mathcal V(X))_{kk}=\frac{\lambda_k}{B}\tr(HX)+\frac{1+\kappa_{\mathrm{feat}}}{B}\lambda_k^2X_{kk}$; for $j\ne k$ it is $1$ only at $(l,m)=(k,j)$ (the case $l=m=j$ gives $\E[(\omega_j^2-1)\omega_j]\E[\omega_k]=0$), so $(\mathcal V(X))_{jk}=\frac{\lambda_j\lambda_k}{B}X_{kj}$ for any $X$, which is $\frac{\lambda_j\lambda_k}{B}X_{jk}$ on the symmetric matrices where $\Sigma_0$ and $\mathcal T$ live (on full matrices, e.g.\ in a basis of single entries, the transpose matters); the $B$ i.i.d.\ draws divide by $B$. \emph{Diagonal system.} The $(k,k)$ entry of Theorem~\ref{thm:F}(2), where the double commutator vanishes, reads $2\lambda_k\gamma_{0,k}-\frac{\eta}{1+\beta}\lambda_k^2\gamma_{0,k}-\frac{\eta}{1-\beta}\big(\frac{\lambda_k}{B}\sum_j\lambda_j\gamma_{0,j}+\frac{1+\kappa_{\mathrm{feat}}}{B}\lambda_k^2\gamma_{0,k}\big)=\frac{\eta}{1-\beta}\sigma_{\xi,k}^2$; substituting $\gamma_{0,j}=\frac{\eta(1+\beta)}{2\lambda_j}c_j$ and dividing by $\eta(1+\beta)$ gives \eqref{eq:mmatrix}. The matrix $\mathcal M=\alpha_{\mathrm{diag}}\Lambda+\frac{\kappa_\beta}{B}\lambda\mathbf 1^\top$ is similar to the symmetric $\widetilde{\mathcal M}=\Lambda^{-1/2}\mathcal M\Lambda^{1/2}=\alpha_{\mathrm{diag}}\Lambda+\frac{\kappa_\beta}{B}\sqrt\lambda\sqrt\lambda^\top$, whose diagonal entries $\lambda_k(1+\kappa_\beta(2+\kappa_{\mathrm{feat}})/B)\ge\lambda_k$ and off-diagonal entries $\kappa_\beta\sqrt{\lambda_j\lambda_k}/B$ are non-negative for every $\kappa_{\mathrm{feat}}\ge-2$, so $I-\frac{\eta}{2(1+\beta)}\mathcal M$ is a nonsingular M-matrix, with a positive solution for every positive right-hand side, if and only if $\frac{\eta}{2(1+\beta)}\lambda_{\max}(\widetilde{\mathcal M})<1$ \citep[Ch.~6]{berman1994}; the rank-one determinant lemma $\det(\mu I-\alpha_{\mathrm{diag}}\Lambda-\frac{\kappa_\beta}{B}\sqrt\lambda\sqrt\lambda^\top)=\det(\mu I-\alpha_{\mathrm{diag}}\Lambda)\big(1-\frac{\kappa_\beta}{B}\sum_k\frac{\lambda_k}{\mu-\alpha_{\mathrm{diag}}\lambda_k}\big)$ identifies $\lambda_{\max}(\widetilde{\mathcal M})$ as the unique root above $\max(\alpha_{\mathrm{diag}}\lambda_1,0)$ (the sum decreases to $0$ there; the top eigenvalue lies in that interval by interlacing), and $\mu=2(1+\beta)/\eta$ gives \eqref{eq:secular}.
\emph{The full operator.} Let $\mathcal T(Q)=\E[F_tQF_t^\top]$ on $\mathbb S^{2P}$ with $F_t=\begin{pmatrix}(1+\beta)I-\eta H_{B_t}&-\beta I\\I&0\end{pmatrix}$ and $G=\mathrm{blkdiag}(\eta^2\Sigma_\xi,0)$. By Lemma~\ref{lem:russodye} it suffices to work on $\mathbb S^{2P}$. A numerical check of $\rho(\mathcal T)$ is safest on the $n(n+1)/2$-dimensional symmetric basis, with $\mathcal T$ in Kraus form, or by the orbit of the identity: a formula for $\mathcal V$ that is wrong only on the antisymmetric part leaves the theorem intact but can move the radius computed on the full basis (on $\mathbb S^{2P}$ the antisymmetric part is excluded, and by Lemma~\ref{lem:russodye} its radius is at most the symmetric one). \emph{Only if:} let $Q$ be any stationary second moment of $(w_t,w_{t-1})$. Its mean $m$ is zero automatically: taking expectations in \eqref{eq:hbP} gives $m=(1+\beta)m-\eta Hm-\beta m$, so $\eta Hm=0$ and $m=0$ since $H\succ0$; Theorem~\ref{thm:F}, which assumes $\E[w_t]=0$, therefore applies. The blocks of $Q$ satisfy Theorem~\ref{thm:F} and its diagonal solves \eqref{eq:mmatrix}, $c=Nc+b$ with $N=\frac{\eta}{2(1+\beta)}\mathcal M\ge0$ entrywise and $b=\sigma_\xi^2/(1-\beta^2)>0$, so $c\ge b>0$ and $Nc=c-b<c$ componentwise; by the Collatz--Wielandt formula $\rho(N)\le\max_k(Nc)_k/c_k<1$, i.e.\ $\frac{\eta}{2(1+\beta)}\lambda_{\max}(\widetilde{\mathcal M})<1$. \emph{If:} for $\eta<\etacrit^{\mathrm{EoSS},P}$ construct $Q$ from the stationary equations with $\Sigma_\xi$ diagonal in the eigenbasis ($\rho(\mathcal T)$ does not depend on $\Sigma_\xi$). The diagonal of $\Sigma_0$ is the positive solution of \eqref{eq:mmatrix}; each off-diagonal entry satisfies $L_{jk}(\eta)(\Sigma_0)_{jk}=0$ with $L_{jk}(\eta)=(\lambda_j+\lambda_k)-\frac{\eta\lambda_j\lambda_k}{1+\beta}(1+\frac{\kappa_\beta}{B})+\frac{\eta\beta(\lambda_j-\lambda_k)^2}{(1-\beta)(1-\beta^2)}$, and $(\Sigma_0)_{jk}=0$ solves it whatever the sign of $L_{jk}$ (its positivity below the ceiling follows from $\rho(\mathcal T)<1$, each off-diagonal pair spanning an invariant subspace, and is not used). So $\Sigma_0=\mathrm{diag}(\gamma_{0,k})$, $\Sigma_1=\frac{A}{1+\beta}\Sigma_0=\mathrm{diag}(\rho_{1,k}\gamma_{0,k})$ with $\rho_{1,k}=1-\frac{\eta\lambda_k}{1+\beta}\in(-1,1)$, and $Q=\begin{pmatrix}\Sigma_0&\Sigma_1\\\Sigma_1^\top&\Sigma_0\end{pmatrix}\succ0$ (its Schur complement is $\mathrm{diag}(\gamma_{0,k}(1-\rho_{1,k}^2))$); the block equations of Theorem~\ref{thm:F} are the blocks of $Q=\mathcal T(Q)+G$, so $(I-\mathcal T)(Q)=G\succeq0$. Since $G$ is only semidefinite, perturb: $(I-\mathcal T)(Q-\varepsilon G)=(1-\varepsilon)G+\varepsilon\mathcal T(G)$ with $\mathcal T(G)=\eta^2\E\big[\binom{A_t}{I}\Sigma_\xi(A_t\ \ I)\big]$, $A_t=(1+\beta)I-\eta H_{B_t}$, whose quadratic form at $(x,y)\ne0$ is at least $\eta^2[(1-\varepsilon)x^\top\Sigma_\xi x+\varepsilon(Ax+y)^\top\Sigma_\xi(Ax+y)]>0$ by Jensen's inequality in $A_t$ (if $x=0$ then $y\ne0$); for $0<\varepsilon<\min(1,\lambda_{\min}(Q)/(\eta^2\|\Sigma_\xi\|_{\mathrm{op}}))$ also $Q_\varepsilon=Q-\varepsilon G\succ0$. A positive operator on $\mathbb S_+^{2P}$ with $Q_\varepsilon\succ0$ and $Q_\varepsilon-\mathcal T(Q_\varepsilon)\succ0$ has $\rho(\mathcal T)<1$: some $\theta<1$ gives $\mathcal T(Q_\varepsilon)\preceq\theta Q_\varepsilon$, hence $\mathcal T^n(Q_\varepsilon)\preceq\theta^nQ_\varepsilon$, and every positive semidefinite $X$ is dominated by a multiple of $Q_\varepsilon$.
\emph{Sandwich.} In the secular equation every denominator $\mu^\star-\alpha_{\mathrm{diag}}\lambda_k$ is positive and lies between $\mu^\star-\max(\alpha_{\mathrm{diag}},0)\lambda_1$ and $\mu^\star-\min(\alpha_{\mathrm{diag}},0)\lambda_1$; replacing all of them by the smaller value gives $1\le\kappa_\beta\tr H/(B(\mu^\star-\max(\alpha_{\mathrm{diag}},0)\lambda_1))$, the upper bound, replacing them by the larger value gives the second lower bound, and keeping $k=1$ only gives $\mu^\star\ge\alpha_{\mathrm{diag}}\lambda_1+\kappa_\beta\lambda_1/B$. For the large-batch statement of the corollary, write $\mu^\star=\alpha_{\mathrm{diag}}\lambda_1+\delta$, $\delta>0$ when $\alpha_{\mathrm{diag}}>0$; separating $k=1$ in \eqref{eq:secular}, $\varepsilon_B\lambda_1/\delta=1-\varepsilon_BT(\mu^\star)$, which is the exact identity. Each denominator of $T(\mu^\star)$ is $\alpha_{\mathrm{diag}}(\lambda_1-\lambda_k)+\delta\ge\alpha_{\mathrm{diag}}(\lambda_1-\lambda_k)$, so $0\le\varepsilon_BT(\mu^\star)\le\varepsilon_BS_\perp/\alpha_{\mathrm{diag}}\le\frac12$ and $\delta=\varepsilon_B\lambda_1(1+r)$ with $0\le r=\frac{\varepsilon_BT}{1-\varepsilon_BT}\le2\varepsilon_BS_\perp/\alpha_{\mathrm{diag}}$. Then $\mu^\star/\lambda_1=1+x$ with $x=\varepsilon_B(2+\kappa_{\mathrm{feat}}+r)\ge0$, and $\frac{\eta\lambda_1}{2(1+\beta)}=\frac1{1+x}=1-\varepsilon_B(2+\kappa_{\mathrm{feat}})-\varepsilon_Br+\frac{x^2}{1+x}$, so $\varrho=\frac{x^2}{1+x}-\varepsilon_Br\in[-2\varepsilon_B^2S_\perp/\alpha_{\mathrm{diag}},\,x^2]$. With $\lambda_1=\lambda_2$ the $k=2$ term of \eqref{eq:secular} is $\lambda_1/\delta$ and doubles the first-order shift. The corollary inserts the bounds (at $\kappa_{\mathrm{feat}}=-2$, $\alpha_{\mathrm{diag}}=1-\kappa_\beta/B$ and the second lower bound is $\lambda_1+\kappa_\beta\tr H_\perp/B$ for $B\le\kappa_\beta$) and $\Var(u_1^\top H_{B_t}u_1)=\lambda_1^2(2+\kappa_{\mathrm{feat}})/B$, which vanishes for Rademacher features. For $P=1$, \eqref{eq:secular} reads $\eta\lambda_1\alpha_{\mathrm{diag}}/(2(1+\beta))+\eta\lambda_1\kappa_\beta/(2(1+\beta)B)=1$, i.e.\ $\eta=2(1+\beta)\lambda_1/(\lambda_1^2+\kappa_\beta\nu_B^2)$ with $\nu_B^2=\lambda_1^2(2+\kappa_{\mathrm{feat}})/B$, which is Proposition~\ref{prop:oneD}.
\end{proof}

\begin{proof}[Proof of Proposition~\ref{prop:G3}]
(1) Write $f=\eta\lambda_1/(2(1+\beta))$, $\varphi_1=1+\beta-\eta\lambda_1$, $\varphi_2=-\beta$ and $\rho_1=\varphi_1/(1-\varphi_2)=1-2f$. An additive excitation of variance $\sigma^2>0$ independent of the past does not change the second-moment operator; with it, squaring the recursion and taking stationary expectations gives $\gamma_0=\big(\varphi_1^2+\varphi_2^2+2\varphi_1\varphi_2\rho_1+\eta^2\nu_{11,B}^2\big)\gamma_0+\eta^2\sigma^2$, with $\rho_1=\varphi_1/(1-\varphi_2)$ from the lag-one equation. The coefficient of the noise-free part equals $1-4(1-\beta^2)f(1-f)$, and $\eta^2\nu_{11,B}^2=4(1+\beta)^2f^2\nu_{11,B}^2/\lambda_1^2$, so a finite positive $\gamma_0$ exists if and only if $f\big(1+\kappa_\beta\nu_{11,B}^2/\lambda_1^2\big)<1$, which by Corollary~\ref{cor:msinstab} is the mean-square threshold of the homogeneous recursion as well. The identity $(1+x)^{-1}=1-x+x^2/(1+x)$ with $0\le x^2/(1+x)\le x^2$ for $x\ge0$ is the formula for $\varrho$.

Let $\varepsilon=\nu_{11,B}^2/\lambda_1^2$, $\mu_3=m_{3,B}/\lambda_1^3$ and $\mu_u=\E u_t=2\kappa_\beta\delta$. On $|u|\le1/2$ the third derivative of $\log(1+\xi)$ is $2/(1+\xi)^3$, hence at most $16$, and Taylor's theorem with the Lagrange remainder gives $\log|1+u|=u-\frac12 u^2+R(u)$ with $|R(u)|\le\frac83|u|^3$. At $\beta=0$ the companion factor is the scalar $-(1+u_t)$, so $\gamma_{\mathrm{Lyap}}=\E\log|1+u_t|$ as soon as $|u_t|\le1/2$ almost surely. Then $\gamma_{\mathrm{Lyap}}=\E u_t-\frac12\E u_t^2+\theta$ with $|\theta|\le\frac83\E|u_t|^3$, $\E u_t=2\delta$ and $\E u_t^2=4\delta^2+4(1+\delta)^2\varepsilon$, hence $\gamma_{\mathrm{Lyap}}=2\delta-2\varepsilon-4\delta\varepsilon-2\delta^2\varepsilon-2\delta^2+\theta$. For $|\delta|\le1/2$, $|4\delta\varepsilon|\le2\delta^2+2\varepsilon^2$ and $2\delta^2\varepsilon\le\varepsilon^2$. Also $|u_t|\le2|\delta|+3|\zeta_t|/\lambda_1$ and $(|a|+|b|)^3\le8(|a|^3+|b|^3)$, so $\E|u_t|^3\le8(8|\delta|^3+27\mu_3)$ and $|\delta|^3\le\frac12\delta^2$. This is the expansion at $\beta=0$. \emph{Unbounded $\zeta_t$ at $\beta=0$.} Write $R(u)=\log|1+u|-(u-\frac12u^2)$ and split into $E_1=\{|u|\le\frac12\}$, $E_2=\{|u|>\frac12,|1+u|>\frac12\}$ and $E_3=\{|1+u|\le\frac12\}$. On $E_1$, $|R|\le\frac83|u|^3$. On $E_2$, $|\log|1+u||\le\max(\log2,|u|)\le8\log2\,|u|^3$ and $|u-\frac12u^2|\le|u|+\frac12u^2\le5|u|^3$ (both because $|u|>\frac12$), so $|R|\le11|u|^3$. On $E_3$, $|u|\ge\frac12$ again gives $|u-\frac12u^2|\le5|u|^3$, and $\E\big[|\log|1+u_t||\mathbf 1_{E_3}\big]=\int_{-1/2}^{1/2}(-\log|v|)f_{1+u}(v)\,dv\le(1+\log2)M_{\mathrm{sing}}$, since $\int_{-1/2}^{1/2}-\log|v|\,dv=1+\log2$. Adding the three regions gives the stated inequality, and the root follows as before with the extra $M_{\mathrm{sing}}$. For Gaussian $\zeta_t$, $1+u_t$ is normal with mean $1+2\delta\in[\frac34,\frac54]$ and standard deviation $2(1+\delta)\sqrt\varepsilon\in[\frac74\sqrt\varepsilon,\frac94\sqrt\varepsilon]$, and every $|v|\le\frac12$ is at distance at least $\frac14$ from the mean. For $\zeta_t=\lambda_1(\chi^2_B/B-1)$, $1+u_t=-1+2(1+\delta)w$ with $w=\chi^2_B/B$, so $|1+u_t|\le\frac12$ forces $w\in[\frac29,\frac67]$; there Stirling's bound $\Gamma(k)\ge\sqrt{2\pi/k}(k/e)^k$ gives a density of $w$ at most $\sqrt{B/(2\pi w^2)}\exp(-\frac B2(w-1-\log w))$ with $w-1-\log w\ge\log\frac76-\frac17>0.0112$, and the change of variable divides by $2(1+\delta)\ge\frac74$.

For $\beta\in(0,1)$ assume $|u_t|\le\varepsilon\le\varepsilon_0(\beta)$ almost surely, and set $K=2\beta/(1-\beta)$. The matrix $P=\begin{pmatrix}1&1\\-1&-\beta^{-1}\end{pmatrix}$ diagonalizes $M_\star=\begin{pmatrix}-(1+\beta)&-\beta\\1&0\end{pmatrix}$ to $\operatorname{diag}(-1,-\beta)$, and $P^{-1}M_tP=\begin{pmatrix}-1-u_t&-u_t\\\beta u_t&-\beta+\beta u_t\end{pmatrix}$. In the coordinate $r_t=y_{2,t}/y_{1,t}$ one has $y_{1,t+1}/y_{1,t}=-(1+Z_t)$ and $r_{t+1}=\beta(r_t-u_t(1+r_t))/(1+Z_t)$ with $Z_t=u_t(1+r_t)$. If $|r|\le K\varepsilon$, then $|Z|\le\varepsilon(1+K\varepsilon)\le1/5$ and
\[
|r'|\le\frac{\beta\varepsilon(K+1+K\varepsilon)}{1-\varepsilon(1+K\varepsilon)}.
\]
The bound $\varepsilon\le(1-\beta)/(8(1+\beta))$ gives $K\varepsilon=\beta/(4(1+\beta))\le1/4$ and $\varepsilon(1+\beta)+K\varepsilon^2\le5(1-\beta)/32<(1-\beta)/2$. Since $\beta(K+1)/K=(1+\beta)/2$, this is $\beta(K+1+K\varepsilon)\le K\big(1-\varepsilon(1+K\varepsilon)\big)$. The ball is therefore forward-invariant. Starting at $r_0=0$ yields $|r_t|\le K\varepsilon\le1/4$ and $|1+Z_t|\ge4/5$. The identity $r_{t+1}-\beta r_t=-\beta Z_t(1+r_t)/(1+Z_t)$ then gives $|r_{t+1}-\beta r_t|\le2\beta|u_t|$, because $(1+|r|)^2/(4/5)\le2$. The $L^2$ triangle inequality produces $(\E r_t^2)^{1/2}\le\beta(\E r_{t-1}^2)^{1/2}+2\beta(\E u_t^2)^{1/2}$, so $\E r_t^2\le\big(2\beta/(1-\beta)\big)^2\E u_t^2$.

On $|Z|\le1/2$, $1/(1+Z)=1-Z+Z^2/(1+Z)$ with $|Z^2/(1+Z)|\le2Z^2$, and $Z=u(1+r)$, so $u(1+r)^2/(1+Z)=u(1+r)^2-u^2(1+r)^3+\psi$ with $|\psi|\le2|u|^3(1+|r|)^4$. Independence of $u_t$ and $r_t$ gives $\E r_{t+1}=\beta\E r_t-\beta\mu_u+\xi_t$ with $|\xi_t|\le C_\beta(\E u_t^2+\E|u_t|^3)$: the terms that contain a factor $r_t$ are estimated by $|\E r_t|\le C_\beta(\E u_t^2)^{1/2}$ and by Lyapunov's comparison $(\E u_t^2)^{3/2}\le\E|u_t|^3$. From $r_0=0$, $\E r_t=-\frac{\beta}{1-\beta}\mu_u+O(\beta^t|\mu_u|)+\rho_t$ with $|\rho_t|\le C_\beta(\E u_t^2+\E|u_t|^3)$.

Since $|Z_t|\le1/2$, $\E\log|1+Z_t|=\E Z_t-\frac12\E Z_t^2+O(\E|Z_t|^3)$. Independence gives $\E Z_t=\mu_u(1+\E r_t)$ and $\E Z_t^2=\E u_t^2\,\E(1+r_t)^2$. Substituting $\E u_t^2=4\kappa_\beta^2\delta^2+4\kappa_\beta^2(1+\delta)^2\varepsilon$ produces the main terms $2\kappa_\beta\delta-2\kappa_\beta^2\varepsilon$. The hypothesis $|u_t|\le\varepsilon_0(\beta)$ bounds $\varepsilon$ by a constant depending only on $\beta$, and $|\delta|\varepsilon\le\frac12(\delta^2+\varepsilon^2)$, so every cross term is $O(\delta^2+\mu_3+\varepsilon^2)$ with a constant depending only on $\beta$. The Ces\`aro mean of the transient $O(\beta^t|\delta|)$ vanishes. In the coordinates of $P$ every $M_t$ maps the closed cone $C_+=\{y:|y_2|\le K\varepsilon\,y_1\}$ into $-C_+$ (the ball $|r|\le K\varepsilon$ is forward-invariant for every realisation and $1+Z_t\ge4/5$), and $e_1$ is interior to $C_+$. For any product $A$ of such matrices and any $x$, with $s=|x_1|+|x_2|/(K\varepsilon)$ both $x+se_1$ and $se_1$ lie in $C_+$, and $c\in C_+$ satisfies $\alpha e_1-c\in C_+$ for $\alpha=c_1+|c_2|/(K\varepsilon)$, so $|(Ac)_1|\le\alpha|(Ae_1)_1|$ and $\|Ax\|\le c_{K\varepsilon}\|x\|\,|(Ae_1)_1|$ with $c_{K\varepsilon}$ depending only on $K\varepsilon$; with $1\le\|y_t\|/|y_{1,t}|\le1+|r_t|$ this gives $\frac1t\log\|M_t\cdots M_1\|=\frac1t\sum_{s<t}\log|1+Z_s|+O(1/t)$. The summands are uniformly bounded, so the expectations converge to $\gamma_{\mathrm{Lyap}}$ \citep{furstenberg1960products}. This is the expansion for $\beta\in(0,1)$.

Write $\gamma_{\mathrm{Lyap}}=2\kappa_\beta\delta-2\kappa_\beta^2\varepsilon+R$ with $|R|\le C_\beta(\delta^2+\mu_3+\varepsilon^2)$ on the $\delta$-interval where $|u_t|\le\varepsilon_0(\beta)$. If $|\zeta_t|\le c\lambda_1$ almost surely, then $|u_t|\le2\kappa_\beta|\delta|+2\kappa_\beta(1+|\delta|)c$, while $\varepsilon\le c^2$ and $\mu_3\le c^3$. Fix $A=2C_\beta(2\kappa_\beta^2+1)/\kappa_\beta$ and $\tau=A(\varepsilon^2+\mu_3)$, and take $c$ small enough that $|\kappa_\beta\varepsilon|+\tau$ lies in the interval and $\tau\le3\kappa_\beta/(4C_\beta)$. Then $2\kappa_\beta\tau\ge C_\beta(2\kappa_\beta^2\varepsilon^2+2\tau^2+\mu_3+\varepsilon^2)$, so $\gamma_{\mathrm{Lyap}}$ is nonnegative at $\delta=\kappa_\beta\varepsilon+\tau$ and nonpositive at $\delta=\kappa_\beta\varepsilon-\tau$. For $d=\delta-\kappa_\beta\varepsilon$ with $\tau\le|d|\le\kappa_\beta/(2C_\beta)$, $\delta^2\le2d^2+2\kappa_\beta^2\varepsilon^2$ gives $|R|\le C_\beta(2d^2+(2\kappa_\beta^2+1)\varepsilon^2+\mu_3)<2\kappa_\beta|d|$, so $\gamma_{\mathrm{Lyap}}$ has the sign of $d$ on that whole range and the almost-sure threshold in it satisfies $|\delta^\star-\kappa_\beta\varepsilon|\le\tau$; no continuity of the limit exponent in $\delta$ is used. At the mean-square threshold $\delta=-\kappa_\beta\varepsilon/(1+\kappa_\beta\varepsilon)$ the same expansion has main term $-2\kappa_\beta^2\varepsilon(2+\kappa_\beta\varepsilon)/(1+\kappa_\beta\varepsilon)$, which dominates $R$ for small $\varepsilon$ and gives $\gamma_{\mathrm{Lyap}}<0$. In every case Jensen's inequality yields $\gamma_{\mathrm{Lyap}}\le\frac12\log\rho(\E[M_t^{\otimes 2}])$, so $\gamma_{\mathrm{Lyap}}\le0$ at the mean-square threshold.

At $\beta=0$ and $|u|\le1/2$, $f(u)=(1+u)^p$ has third derivative $p(p-1)(p-2)(1+u)^{p-3}$, bounded on $[-1/2,1/2]$ by a constant $C_p$. Thus $\E|1+u_t|^p=1+p\E u_t+\frac{p(p-1)}2\E u_t^2+\rho$ with $|\rho|\le C_p\E|u_t|^3$. For $c_p$ small enough that $|u_t|\le1/2$ throughout a neighbourhood of $\delta=-(p-1)\varepsilon$, $\log\E|1+u_t|^p=2p\delta+2p(p-1)\varepsilon+O_p(\delta^2+\mu_3+\varepsilon^2)$, whose root is $\delta=-(p-1)\varepsilon+O_p(\mu_3+\varepsilon^2)$.

The third-moment bound does not close the estimate. At $\beta=0$, $\delta=0$, $\lambda_1=1$ and $\zeta=\pm\frac12(1-\varepsilon)$ with equal probability, $\E\log|1+2\zeta|=\frac12\log\big(\varepsilon(2-\varepsilon)\big)$ and $-2\E\zeta^2=-\frac12(1-\varepsilon)^2$, so the ratio of the difference to $\E|\zeta|^3$ tends to $-\infty$ as $\varepsilon\downarrow0$. Rosenthal's inequality supplies $\nu_{11,B}^2=O(1/B)$ and $m_{3,B}=O(B^{-3/2})$ for an average of i.i.d.\ third-moment summands, and these rates may be substituted above only under $|u_t|\le\varepsilon_0(\beta)$. For $\lambda_1=1$, $B=64$ and $\beta=0$, $2/(1+2/64)=64/33$ rounds to $1.9394$, the first-order values $2(1\mp2/64)$ are $1.9375$ and $2.0625$, and the two quadrature roots are the output of \texttt{paper/checks/g3-lyapunov-remainder.mjs}. If $\zeta_t=0$ then $u_t=2\kappa_\beta\delta$ and both thresholds equal $2(1+\beta)/\lambda_1$. (2) Take the trace of Theorem~\ref{thm:F}(2) at a learning rate with a stationary second moment: $\tr[H,[H,\Sigma_0]]=0$ by cyclicity and $\tr\mathcal V(\Sigma_0)=\frac1B\E\tr(H_i\Sigma_0H_i-H\Sigma_0H)$ because the $B$ draws are i.i.d.; under the lower bound on $\E[H_iXH_i]$ the left side is at most $\tr(H\Sigma_0)\big(2-\frac{\eta c_0\tr H}{(1-\beta)B}\big)-\eta\tr(H^2\Sigma_0)\big(\frac1{1+\beta}-\frac1{(1-\beta)B}\big)$, and for $B\ge\kappa_\beta$ the last bracket is non-negative, so a positive right side forces $2>\eta c_0\tr H/((1-\beta)B)$ for every such $\eta$, and the threshold, their supremum, satisfies the non-strict inequality. Multiplying Theorem~\ref{thm:F}(2) by $H^{-1}$ first, $\tr(H^{-1}[H,[H,\Sigma_0]])=\tr(H\Sigma_0)-2\tr(H\Sigma_0)+\tr(H\Sigma_0)=0$ and $\tr(H^{-1}\mathcal V(\Sigma_0))=\frac1B\tr(\E[(H_i-H)H^{-1}(H_i-H)]\Sigma_0)$ by cyclicity, which gives the $H^{-1}$ identity; as $\Sigma_0\succ0$ and the right side is positive, the bracket cannot be negative semidefinite, which is the $\lambda_{\min}$ condition.
\end{proof}

\begin{remark}[Trace bracket of the secular root]\label{rem:tracebracket}
With $c=\eta\alpha_{\mathrm{diag}}/(2(1+\beta))>0$ and $f(x)=x/(1-cx)$, $f(x)/x$ is increasing on $[0,1/c)$, so for $0\le\lambda_k\le\lambda_2$ ($k\ge2$), $\tr H-\lambda_1\le\sum_{k\ge2}f(\lambda_k)\le(\tr H-\lambda_1)f(\lambda_2)/\lambda_2$. The lower inequality also holds for negative $\lambda_k$ ($f(x)\ge x$ for $x<1/c$); the upper one then needs $\tr H$ replaced by $\tr H+\sum_{\lambda_k<0}|\lambda_k|$, since only the non-negative eigenvalues contribute positive terms. Both networks of Section~\ref{sec:phase2} have negative eigenvalues and the bracket there uses $\tr H$, so the upper sum, and with it the lower end of the bracket on $\eta^\star$, is formally an approximation at both sizes. At $P=3\,610$ the dense spectrum of Section~\ref{sec:phase4} has $1\,528$ negative eigenvalues, all $\ge\lambda_{\min}=-0.083330$, so the negative mass is at most $127.33$, and $\tr H+127.33$ in the upper sum gives the guaranteed lower end $0.00411$ in place of $0.00492$; at $P=12\,826$, where $\lambda_{\min}\le-0.072=-4.9\times10^{-4}\lambda_1$, the negative mass was not measured. The left side of \eqref{eq:secular} is increasing in $\eta$, so the root computed with the upper sum is a lower bound on $\eta^\star$ and the root with the lower sum an upper bound: a two-sided bracket from $\lambda_1$, $\lambda_2$ and $\tr H$ alone, inside the sandwich \eqref{eq:sandwich}, and tight when the bulk is concentrated near $\lambda_2$ or near $0$. Phase~2 uses it on the networks, where only $\lambda_1$, $\lambda_2$ and $\tr H$ are available (\texttt{compute-secular-ceiling-bracket-from-trace.ts}).
\end{remark}

\begin{proof}[Proof of Theorem~\ref{thm:G5}]
(1) In the eigenbasis $(z_iz_i^\top)_{lm}=\sqrt{s_ls_m}\,\omega_{i,l}\omega_{i,m}$ and $\mathcal V(X)=\frac1B\E[(zz^\top-\Sigma_z)X(zz^\top-\Sigma_z)]$ by independence of the $B$ draws, so $(\mathcal V(X))_{jk}=\frac1B\sum_{l,m}\sqrt{s_js_ls_ms_k}\operatorname{Cov}(\omega_j\omega_l,\omega_m\omega_k)X_{lm}$. By independence and standardisation the covariance is non-zero only when $\{j,l\}=\{m,k\}$ as multisets: for $j=k$ the pairs $l=m\ne k$ contribute $\frac{s_ks_l}{B}X_{ll}$ and $l=m=k$ contributes $\frac{2+\kappa_{\mathrm{feat}}}{B}s_k^2X_{kk}$; for $j\ne k$ only $(l,m)=(k,j)$ contributes $\frac{s_js_k}{B}X_{kj}$, equal to $\frac{s_js_k}{B}X_{jk}$ for symmetric $X$. The diagonal of $[H,[H,\Sigma_0]]$ vanishes (Theorem~\ref{thm:F}), which gives the equation. (2) Dividing row $k$ by $2\lambda_k$ and using $\frac{\eta}{2(1-\beta)}=\frac{\eta\kappa_\beta}{2(1+\beta)}$ gives $M(\eta)\gamma_0=\frac{\eta}{2(1-\beta)}\Lambda^{-1}\sigma_\xi^2$ with $M(\eta)=I-\frac{\eta}{2(1+\beta)}(\operatorname{diag}(d_k)+\frac{\kappa_\beta}{B}\Lambda^{-1}ss^\top)$, similar through $\Lambda^{1/2}$ to $I-\frac{\eta}{2(1+\beta)}\widetilde{\mathcal M}$, whose off-diagonal entries $\frac{\kappa_\beta}{B}\sqrt{\rho_j\rho_k}$ are non-negative and whose diagonal entries $\lambda_k+\kappa_\beta\frac{2+\kappa_{\mathrm{feat}}}{B}\rho_k$ are positive for every $\kappa_{\mathrm{feat}}\ge-2$; $M(\eta)$ is a non-singular M-matrix with non-negative inverse if and only if $\frac{\eta}{2(1+\beta)}\lambda_{\max}(\widetilde{\mathcal M})<1$, and the rank-one determinant lemma $\det(\mu I-\operatorname{diag}(d_k)-\frac{\kappa_\beta}{B}vv^\top)=\prod_k(\mu-d_k)\big(1-\frac{\kappa_\beta}{B}\sum_k\frac{\rho_k}{\mu-d_k}\big)$ at $\mu=2(1+\beta)/\eta$ gives the secular equation. For $k\notin\mathcal A$, $v_k=0$ and $d_k=\lambda_k$, so $\widetilde{\mathcal M}=\operatorname{diag}(\lambda_k)_{k\notin\mathcal A}\oplus\widetilde{\mathcal M}_{\mathcal A}$ and $\lambda_{\max}(\widetilde{\mathcal M})=\max\big(\max_{k\notin\mathcal A}\lambda_k,\lambda_{\max}(\widetilde{\mathcal M}_{\mathcal A})\big)$. When $\mathcal A\ne\emptyset$, every $v_k$ with $k\in\mathcal A$ is positive, so $g(\mu)=\frac{\kappa_\beta}{B}\sum_{k\in\mathcal A}\frac{\rho_k}{\mu-d_k}$ decreases strictly from $+\infty$ to $0$ on $(\max_{k\in\mathcal A}d_k,\infty)$, where the eigenvalues of $\widetilde{\mathcal M}_{\mathcal A}$ are the solutions of $g(\mu)=1$; there is exactly one, $\lambda_{\max}(\widetilde{\mathcal M}_{\mathcal A})$. Its diagonal entries $\lambda_k+\kappa_\beta\frac{2+\kappa_{\mathrm{feat}}}{B}\rho_k=d_k+\frac{\kappa_\beta}{B}\rho_k$ are at least $\lambda_k$ and at least $d_k$, so $\lambda_{\max}(\widetilde{\mathcal M}_{\mathcal A})\ge\max_{k\in\mathcal A}\max(d_k,\lambda_k)$: in $\eta=2(1+\beta)/\mu$ the root lies in $(0,\min_{k\in\mathcal A}2(1+\beta)/\max(d_k,\lambda_k)]$, is the only root there, and $\eta^\star$ is the smaller of it and $\min_{k\notin\mathcal A}2(1+\beta)/\lambda_k$. The equivalence with $\rho(\mathcal T)<1$ on $\mathbb S^{2P}$ is the argument in the proof of Theorem~\ref{thm:G} (under ``Only if'' and ``If'') with the same $Q_\varepsilon$, since the off-diagonal equations are decoupled scalar equations solved by $(\Sigma_0)_{jk}=0$ for co-diagonal $\Sigma_\xi$ and $\mu^\star\ge\widetilde{\mathcal M}_{kk}\ge\lambda_k$ gives the deterministic condition $\eta<2(1+\beta)/\lambda_k$ for every mode. (3) Weyl's inequality with $\min_kd_k\,I\preceq\operatorname{diag}(d_k)\preceq\max_kd_k\,I$ and $\|v\|^2=\mathcal T_{\mathrm{flat}}$, and the Rayleigh quotients at $e_k$, give the sandwich; dropping $\min_kd_k$ from the lower bound, which gives $\eta^\star\le2(1-\beta)B/\mathcal T_{\mathrm{flat}}$, needs $\min_kd_k\ge0$; for the flat-mode rate, the Rayleigh quotient of $\widetilde{\mathcal M}$ at $(e_j+e_k)/\sqrt2$, $j\ne k$, is $\frac12(\lambda_j+\lambda_k)+\frac{\kappa_\beta}{2B}\big((2+\kappa_{\mathrm{feat}})(\rho_j+\rho_k)+2\sqrt{\rho_j\rho_k}\big)\ge\frac{\kappa_\beta}B\sqrt{\rho_j\rho_k}=\frac{\kappa_\beta s_js_k}{B\sqrt{\lambda_j\lambda_k}}$ for every $\kappa_{\mathrm{feat}}\ge-2$, so $\mu^\star\ge\frac{\kappa_\beta s_js_k}{B\sqrt{\lambda_j\lambda_k}}$ and $\eta^\star=2(1+\beta)/\mu^\star\le\frac{2(1-\beta)B\sqrt{\lambda_j\lambda_k}}{s_js_k}$; for $P=1$ with Rademacher features $\widetilde{\mathcal M}=\lambda_1$ and $\eta^\star=2(1+\beta)/\lambda_1$; under the tether $H_i\mapsto H_i+\mu I$ and $H\mapsto H+\mu I$ leave $H_i-H$, hence $s_k$, unchanged.
\end{proof}

\begin{proof}[Proof of Theorem~\ref{thm:G6}]
(1) At $\mathcal V=0$, $\rho(A_0)<1$ on $(0,\eta_{\mathrm{det}})$ and the $(1,1)$ block of $\sum_{t\ge0}A_0^t\operatorname{blkdiag}(\eta^2\Sigma_\xi,0)(A_0^t)^\top\succeq0$ equals $\frac1{1-\beta}\mathcal N_{\eta,\beta}(\Sigma_\xi)$, so $\mathcal L_{\eta,\beta}^{-1}$ preserves $\mathbb S_+^P$ and $[1/\ell_{jk}(\eta,\beta)]=\frac1\eta\mathcal N_{\eta,\beta}(\mathbf 1\mathbf 1^\top)\succeq0$ (only semidefinite: for $H=\lambda I$ it is a multiple of $\mathbf 1\mathbf 1^\top$). In the orthonormal basis $\{E_{jj}\}_j\cup\{(E_{jk}+E_{kj})/\sqrt2\}_{j<k}$ of $\mathcal H=(\mathbb S^P,\langle\cdot,\cdot\rangle_F)$, in the eigenbasis of $H$, $\mathcal N_{\eta,\beta}$ is diagonal with entries $\eta/\ell_{jk}(\eta,\beta)>0$ (for $\eta<\eta_{\mathrm{det}}$, $\ell_{jk}\ge\lambda_j+\lambda_k-\eta\lambda_j\lambda_k/(1+\beta)>\lambda_j(1-\lambda_k/\lambda_1)+\lambda_k(1-\lambda_j/\lambda_1)\ge0$), so $\mathcal N_{\eta,\beta}\succ_{\mathcal H}0$, and $\mathcal K_{\eta,\beta}=\frac1{1-\beta}\mathcal N_{\eta,\beta}\mathcal V$ is similar via $\mathcal N_{\eta,\beta}^{-1/2}$ to the symmetric operator $\widetilde{\mathcal K}_{\eta,\beta}=\frac1{1-\beta}\mathcal N_{\eta,\beta}^{1/2}\mathcal V\mathcal N_{\eta,\beta}^{1/2}$ on $\mathcal H$, which has a real spectrum and the inertia of $\mathcal V$ (for $\Delta H=\operatorname{diag}(1,-1)$, $\mathcal V(E_{12}+E_{21})=-(E_{12}+E_{21})$); since $\mathcal K_{\eta,\beta}$ preserves $\mathbb S_+^P$, the Krein--Rutman theorem \citep{krein1948linear} gives $\rho(\mathcal K_{\eta,\beta})=\lambda_{\max}(\mathcal K_{\eta,\beta})$ with an eigenmatrix $X_\eta\in\mathbb S_+^P\setminus\{0\}$. For $\mathcal T=\mathcal T_0+\Delta$ on $\mathbb S_+^{2P}$, $(I-\mathcal T_0)^{-1}=\sum_{t\ge0}\mathcal T_0^t\ge0$ gives $\rho(\mathcal T)<1\iff\rho((I-\mathcal T_0)^{-1}\Delta)<1$ and $\rho(\mathcal T)=1\iff\rho((I-\mathcal T_0)^{-1}\Delta)=1$ \citep{berman1994}, and the compression of $(I-\mathcal T_0)^{-1}\Delta$ to the $(1,1)$ block is $\mathcal K_{\eta,\beta}$; a compression alone does not fix the spectral radius, and the equality $\rho((I-\mathcal T_0)^{-1}\Delta)=\rho(\mathcal K_{\eta,\beta})$ is the factorization $\mathcal R=\mathcal J\circ\Pi$, $\Pi\circ\mathcal J=\mathcal K_{\eta,\beta}$ of Theorem~\ref{thm:splitting}(1), whose proof uses the invertibility of $\mathcal L_{\eta,\beta}$ established here and Theorem~\ref{thm:F}(2), and not $\beta>0$. (2) Substituting $Y=\mathcal N_{\eta,\beta}^{-1/2}X$ in the Rayleigh quotient of $\widetilde{\mathcal K}_{\eta,\beta}$ gives the displayed formula for $\rho(\mathcal K_{\eta,\beta})=\lambda_{\max}(\widetilde{\mathcal K}_{\eta,\beta})$; testing $X=I_P$ gives $\langle I_P,\mathcal V(I_P)\rangle_F=\frac1B\E\|\Delta H_i\|_F^2>0$, and since each weight $\ell_{jk}(\eta,\beta)/\eta=(\lambda_j+\lambda_k)/\eta+b_{jk}(\beta)$ is strictly decreasing in $\eta$ and tends to $+\infty$ as $\eta\downarrow0$, evaluating at an optimizer $X_{\eta_1}$ for $0<\eta_1<\eta_2<\eta_{\mathrm{det}}$ gives $\rho(\mathcal K_{\eta_2,\beta})>\rho(\mathcal K_{\eta_1,\beta})$ and $\lim_{\eta\downarrow0}\rho(\mathcal K_{\eta,\beta})=0$. (3) Let $w_1(\eta)\triangleq\ell_{11}(\eta,\beta)/\eta=\frac{\lambda_1^2}{(1+\beta)\eta}(\eta_{\mathrm{det}}-\eta)\downarrow0$ while $\ell_{jk}(\eta,\beta)/\eta\ge c_0>0$ outside $\mathcal I_1\times\mathcal I_1$ ($\mathcal I_1=\{j:\lambda_j=\lambda_1\}$). If $P_1\Delta H_iP_1\ne0$ for some $i$, $X=P_1$ gives $\langle P_1,\mathcal V(P_1)\rangle_F=\frac1B\E\|P_1\Delta H_iP_1\|_F^2>0$ and $\rho(\mathcal K_{\eta,\beta})=\Theta(w_1(\eta)^{-1})$. If $P_1\Delta H_iP_1=0$ for all $i$ and $\Delta H_iP_1\ne0$ for some $i$, decomposing $X=X_{11}+X_\perp$ gives $P_1\mathcal V(X_{11})P_1=0$, so $\langle X,\mathcal V(X)\rangle_F=2\langle X_{11},\mathcal V(X_\perp)\rangle_F+\langle X_\perp,\mathcal V(X_\perp)\rangle_F$ and $\rho(\mathcal K_{\eta,\beta})=O(w_1(\eta)^{-1/2})$, while testing $X=sP_1+\mathcal V(P_1)$ ($\mathcal V(P_1)\ne0$, $\langle P_1,\mathcal V(P_1)\rangle_F=0$) at $s=w_1(\eta)^{-1/2}$ gives numerator $2s\|\mathcal V(P_1)\|_F^2+O(1)$ and denominator $s^2\|P_1\|_F^2w_1(\eta)+\Theta(1)$ (the top block of $\mathcal V(P_1)$ is $P_1\mathcal V(P_1)P_1=0$), hence $\rho(\mathcal K_{\eta,\beta})=\Theta(w_1(\eta)^{-1/2})$. If $\Delta H_iP_1=0$ for all $i$, $\Delta H_iX\Delta H_i=\Delta H_i(I-P_1)X(I-P_1)\Delta H_i$ vanishes on $\mathcal I_1\times\mathcal I_1$, so the quotient restricts to $\widetilde X_{jk}=0$ on $\mathcal I_1\times\mathcal I_1$, where all weights are at least $c_0>0$, and $\rho(\mathcal K_{\eta,\beta})=O(1)$. Finally, $\rho(\mathcal K_{\eta,\beta})<1\iff\lambda_{\max}(\widetilde{\mathcal K}_{\eta,\beta})<1\iff\mathcal L_0-\eta\mathcal B_\beta\succ_{\mathcal H}0\iff\eta\lambda_{\max}(\mathcal L_0^{-1/2}\mathcal B_\beta\mathcal L_0^{-1/2})<1$. For $\eta\ge\eta_{\mathrm{det}}$ the iteration is not mean-square stable: $\mathcal T=\mathcal T_0+\Delta$ with $\Delta$ preserving the cone $\mathbb S_+^{2P}$, so $\rho(\mathcal T)\ge\rho(\mathcal T_0)=\rho(A_0)^2\ge1$ \citep[Ch.~1]{berman1994}. And $1/\lambda_{\max}(\mathcal L_0^{-1/2}\mathcal B_\beta\mathcal L_0^{-1/2})\le\eta_{\mathrm{det}}$: for $X=u_1u_1^\top$ the commutator term vanishes and $\langle X,\mathcal B_\beta X\rangle_F/\langle X,\mathcal L_0X\rangle_F\ge\frac{\lambda_1^2/(1+\beta)}{2\lambda_1}=1/\eta_{\mathrm{det}}$. Hence the ceiling is $1/\lambda_{\max}(\mathcal L_0^{-1/2}\mathcal B_\beta\mathcal L_0^{-1/2})$ in every case.
\end{proof}

\begin{proof}[Proof of Theorem~\ref{thm:modal-bridge}]
(1) From $(\mathcal L_0-\eta^\star\mathcal B_\beta)(\Sigma_0^\star)=0$, taking the trace and using $\tr(H\Sigma_0^\star)>0$ ($H\succ0$, $\Sigma_0^\star\in\mathbb S_+^P\setminus\{0\}$), $\tr[H,[H,\Sigma_0^\star]]=0$, $\tr(H\Sigma_0^\star H)=\tr(H^2\Sigma_0^\star)$, and $\tr(\mathcal V(\Sigma_0^\star))=\tr((\E[H_{B_t}^2]-H^2)\Sigma_0^\star)\ge0$, then dividing by $2\eta^\star\tr(H\Sigma_0^\star)$, gives the first identity in \eqref{eq:exact-P-interp}; $\lambda_PH\preceq H^2\preceq\lambda_1H$ gives $a_\star\in[\lambda_P,\lambda_1]$, and multiplying by $2(1\mp\beta)$ gives the sandwich for $a_\star+v_\star$. For any $S\in\mathbb S_+^P\setminus\{0\}$ with $[H,S]=0$, testing $\frac1{\eta^\star}=\max_{X\ne0}\frac{\langle X,\mathcal B_\beta(X)\rangle_F}{\langle X,\mathcal L_0(X)\rangle_F}$ at $X=S^{1/2}$ ($[H,X]=0$) gives $\frac1{\eta^\star}\ge\frac{a(S)}{2(1+\beta)}+\frac{v(S)}{2(1-\beta)}$; at $S=I_P$, $\frac{\bar a}{2(1+\beta)}+\frac{\bar v}{2(1-\beta)}=\frac{1+\frac{2\beta}{1-\beta}\theta_{\mathrm{noise}}}{2(1+\beta)}(\bar a+\bar v)$, and under $\E[g_B(\theta_0)g_B(\theta_0)^\top\mid H_B]=\sigma_g^2H_B(\theta_0)$ one has $\E[g_B^\top H_Bg_B]=\sigma_g^2\tr(\E[H_B^2])$ and $\E\|g_B\|_2^2=\sigma_g^2\tr H$, so $\overline{\mathrm{BS}}(\theta_0)=\bar a+\bar v$; without that hypothesis, $R_BW_B=g_B^\top H_Bg_B/\E_B\|g_B\|_2^2$ and $\E_B[W_B]=1$ give $\overline{\mathrm{BS}}(\theta_0)-\mathrm{BS}(\theta_0)=\E_B[R_BW_B]-\E_B[R_B]\E_B[W_B]=\operatorname{Cov}_B(R_B,W_B)$, and Cauchy--Schwarz gives $|\overline{\mathrm{BS}}(\theta_0)-\mathrm{BS}(\theta_0)|\le\sigma_{\mathrm{BS}}(\theta_0)\mathrm{CV}_{g^2}(\theta_0)$. If $\eta^\star\le\eta<\eta_{\mathrm{det}}$, Theorem~\ref{thm:G6}(2) gives $\rho(\mathcal K_{\eta,\beta})\ge1$ and Theorem~\ref{thm:G6}(1) $\rho(\mathcal T)\ge1$; if $\eta\ge\eta_{\mathrm{det}}$, $\mathcal T(Q)\succeq\bar FQ\bar F^\top$ with $\bar F=\E F_t$ gives $\rho(\mathcal T)\ge\rho(\bar F)^2\ge1$. Corollary~\ref{cor:msinstab} then yields \eqref{eq:bs-fixed-instab}. In Theorem~\ref{thm:G}, $\Sigma_0^\star=\operatorname{diag}(\gamma_{0,k}^\star)$ is diagonal in the eigenbasis of $H$ and $\tr(\mathcal V(\Sigma_0^\star))=\frac{\tr H}{B}\tr(H\Sigma_0^\star)+\frac{1+\kappa_{\mathrm{feat}}}{B}\tr(H^2\Sigma_0^\star)$. For (b), any $g_B=\frac1BZr\ne0$ ($u=r/\|r\|_2$, $D\triangleq G_Z-I_B$) has $v\triangleq(Z^\top Z)^{1/2}r\ne0$, $\|g_B\|_2^2=\frac1{B^2}\|v\|_2^2=\frac{\tr H}{B^2}(1+u^\top Du)$, and $g_B^\top H_Bg_B=\frac{\tr H}{B^3}v^\top G_Zv=\frac{(\tr H)^2}{B^3}(1+2u^\top Du+\|Du\|_2^2)$, so $\frac{u^\top G_Z^2u}{u^\top G_Zu}-1=u^\top Du+\frac{\|Du\|_2^2-(u^\top Du)^2}{u^\top G_Zu}=u^\top Du+\frac{\|(I_B-uu^\top)Du\|_2^2}{u^\top G_Zu}$ (giving $\mathrm{BS}(\theta)\ge\frac{\tr H}{B}$ when $u(\theta)$ is independent of $G_Z$, since $\E_B[D]=0$) and $\big|\frac{g_B^\top H_Bg_B}{\|g_B\|_2^2}-\frac{\tr H}{B}\big|\le\frac{\tr H}{B}\|D\|_{\mathrm{op}}\le\frac{\tr H}{B}\|D\|_F$; since $\E[D_{bb}^2]=\frac{(2+\kappa_{\mathrm{feat}})\tr(H^2)}{(\tr H)^2}$ and $\E[D_{bb'}^2]=\frac{\tr(H^2)}{(\tr H)^2}$ for $b\ne b'$ are mutually orthogonal in $L^2$, $\E\|Du\|_2^2=(B+1+\kappa_{\mathrm{feat}})\frac{\tr(H^2)}{(\tr H)^2}$, $\E[(u^\top Du)^2]=(2+\kappa_{\mathrm{feat}}\sum_{b=1}^Bu_b^4)\frac{\tr(H^2)}{(\tr H)^2}$, and $\E\|D\|_F\le(\E\|D\|_F^2)^{1/2}=\delta_{\mathrm{Gram}}$, giving \eqref{eq:bs-pointwise-gram} and, via $\frac{2(1+\beta)}{\eta^\star}\ge\kappa_\beta\frac{\tr H}{B}$ when $\kappa_{\mathrm{feat}}\ge-1$ (Theorem~\ref{thm:G}), the divergence criterion.

(2) Theorem~\ref{thm:F}(2) reads $(\mathcal L_0-\eta\mathcal B_\beta)(\Sigma_0(\eta))=\frac{\eta}{1-\beta}\Sigma_\xi$, or $(I-\eta\mathcal L_0^{-1}\mathcal B_\beta)(\Sigma_0(\eta))=\frac{\eta}{1-\beta}\mathcal L_0^{-1}(\Sigma_\xi)$. Since $\mathcal L_0^{-1}\mathcal B_\beta$ is self-adjoint in $\langle\cdot,\cdot\rangle_{\mathcal L_0}$ with simple top eigenpair $(1/\eta^\star,\Sigma_0^\star)$, $\langle\Sigma_0^\star,\mathcal L_0^{-1}(\Sigma_\xi)\rangle_{\mathcal L_0}=\tr(\Sigma_0^\star\Sigma_\xi)$, and $\|\Sigma_0^\star\|_{\mathcal L_0}^2=2\tr(H(\Sigma_0^\star)^2)>0$, orthogonal projection onto $\operatorname{span}(\Sigma_0^\star)$ and $(\Sigma_0^\star)^{\perp_{\mathcal L_0}}$ gives \eqref{eq:sigma0-pole}, where for every eigenvalue $\mu_m\le 1/\eta^\star-\Delta_{\mathcal G}$ ($m\ge2$), $1-\eta\mu_m=(1-\eta/\eta^\star)+\eta\Delta_{\mathcal G}\ge\eta\Delta_{\mathcal G}$. For $\eta^\star<\eta_{\mathrm{det}}$, differentiating $\mathcal L_{\eta,\beta}(\Sigma_0^\star(\eta))=\frac{\eta}{(1-\beta)r(\eta)}\mathcal V(\Sigma_0^\star(\eta))$ at $\eta=\eta^\star$ ($r(\eta^\star)=1$) against $\Sigma_0^\star$ and using $\eta^\star\partial_\eta\mathcal L_{\eta,\beta}\big|_{\eta^\star}(\Sigma_0^\star)=\mathcal L_{\eta^\star,\beta}(\Sigma_0^\star)-\mathcal L_0(\Sigma_0^\star)$ gives $d_\star$. Since $\E[w_t]=0$ and $(H_{B_t},\xi_t)$ is independent of $w_t$, $\E[w_t^\top H_{B_t}^2\xi_t]=0$ and $\E[w_t^\top H_{B_t}\xi_t]=0$, so $\E[g_{B_t}^\top H_{B_t}g_{B_t}]=\tr(\E[H_{B_t}^3]\Sigma_0(\eta))+\E[\xi_t^\top H_{B_t}\xi_t]$ and $\E\|g_{B_t}\|_2^2=\tr(\E[H_{B_t}^2]\Sigma_0(\eta))+\tr\Sigma_\xi$; dividing by $\tr\Sigma_0(\eta)\to\infty$ gives \eqref{eq:bs-star-limit}. When $H_{B_t}\succeq0$ a.s., Cauchy--Schwarz on $(\Sigma_0^\star)^{1/2}H_{B_t}^{1/2}$ and $(\Sigma_0^\star)^{1/2}H_{B_t}^{3/2}$ gives, for each realisation, $\tr(H_{B_t}^2\Sigma_0^\star)\le\tr(H_{B_t}\Sigma_0^\star)^{1/2}\tr(H_{B_t}^3\Sigma_0^\star)^{1/2}$, and Cauchy--Schwarz in the probability space then gives $(\tr(\E[H_{B_t}^2]\Sigma_0^\star))^2\le\tr(H\Sigma_0^\star)\tr(\E[H_{B_t}^3]\Sigma_0^\star)$. In Theorem~\ref{thm:G}, expanding $H_{B_t}^3$ over the $B(B-1)(B-2)$ distinct triples, $3B(B-1)$ pairs and $B$ diagonal triples of $H_i=z_iz_i^\top$ gives $\operatorname{diag}(\E[H_{B_t}^3])=\big[(\frac{\tr H}{B})^2+\frac{q_{\mathrm{feat}}}{B}\tr(H^2)\big]\Lambda+2\frac{\tr H}{B}q_{\mathrm{feat}}\Lambda^2+r_{\mathrm{feat}}\Lambda^3$, while $(\mu^\star-\alpha_{\mathrm{diag}}\lambda_k)\gamma_{0,k}^\star=\frac{\kappa_\beta}{B}\tr(H\Sigma_0^\star)$ ($\mu^\star=2(1+\beta)/\eta^\star$) multiplied by $\lambda_k$ and $\lambda_k^2$ and summed over $k$ gives $\mu^\star-\alpha_{\mathrm{diag}}a_\star=\frac{\kappa_\beta}{B}\tr H$, $\mu^\star a_\star-\alpha_{\mathrm{diag}}m_{2,\star}=\frac{\kappa_\beta}{B}\tr(H^2)$, $\alpha_{\mathrm{diag}}(m_{2,\star}-a_\star^2)=\frac{\kappa_\beta}{B}((\tr H)a_\star-\tr(H^2))$, and \eqref{eq:bs-rf-exact}. Write $N$ and $D=\frac{\tr H}{B}+q_{\mathrm{feat}}a_\star$ for the numerator and denominator of the fraction in \eqref{eq:bs-rf-exact}. Small batch: $\frac{2(1-\beta)}{\eta^\star}=\frac{\tr H}{B}+\frac{1-\beta}{1+\beta}\alpha_{\mathrm{diag}}a_\star$ with $a_\star\le\lambda_1$ and $|\frac{1-\beta}{1+\beta}\alpha_{\mathrm{diag}}|\le1+\frac{|1+\kappa_{\mathrm{feat}}|}{B}$, so $\frac{2(1-\beta)}{\eta^\star}=\frac{\tr H}{B}(1+O(\frac{B\lambda_1}{\tr H}))$, likewise $\mathcal C_{\mathrm{mod}}=\frac{2(1-\beta)}{\eta^\star}+O(\lambda_1)$, and $D\ge\frac{\tr H}{B}$ with $N\le q_{\mathrm{feat}}\lambda_1\frac{\tr H}{B}+O(\lambda_1^2)$ gives $N/D=\frac{\tr H}{B}O(\frac{B\lambda_1}{\tr H})$. Large batch, $\varepsilon\triangleq\frac{\kappa_\beta\tr H}{B\lambda_1}\to0$: then $B\to\infty$, $\alpha_{\mathrm{diag}},q_{\mathrm{feat}}\to1$ and $r_{\mathrm{feat}}-q_{\mathrm{feat}}^2=O(1/B)$; $\mu^\star>\alpha_{\mathrm{diag}}\lambda_1$ and $\mu^\star-\alpha_{\mathrm{diag}}a_\star=\frac{\kappa_\beta}{B}\tr H$ give $a_\star\ge\lambda_1(1-\varepsilon/\alpha_{\mathrm{diag}})$ and $\mu^\star=\lambda_1(1+O(\varepsilon))$; the moment identity gives $m_{2,\star}-a_\star^2\le\frac{\kappa_\beta\tr H\,\lambda_1}{B\alpha_{\mathrm{diag}}}=\lambda_1^2\,O(\varepsilon)$; hence $N/D=\lambda_1\,O(\varepsilon+\frac1B)$ and $\mathcal C_{\mathrm{mod}}=\mu^\star-\frac{2\beta}{1-\beta}v_\star=\mu^\star(1+O(\varepsilon))$, since $v_\star=\frac{\tr H}{B}+\frac{1+\kappa_{\mathrm{feat}}}{B}a_\star$.

(3) In the eigenbasis of Theorem~\ref{thm:G5}, $\Sigma_0^\star=\operatorname{diag}(c_1^\star,\dots,c_P^\star)$ is diagonal and $(\E[H_{B_t}^2])_{mm}=\lambda_m^2+A_m+C_m$ with $A_m\triangleq\frac{s_m}{B}\sum_{q\ne m}s_q\ge0$ and $C_m\triangleq\frac{2+\kappa_{\mathrm{feat}}}{B}s_m^2\ge0$. With $\Delta_{B_t}=H_{B_t}-H$, expanding $\E[H_{B_t}^3]=H^3+2H\mathcal V(I_P)+\mathcal V(H)+\E[\Delta_{B_t}^3]$ gives $0\le(H^3+2H\mathcal V(I_P)+\mathcal V(H))_{mm}\le 3M_0(\E[H_{B_t}^2])_{mm}$ and $|(\E[\Delta_{B_t}^3])_{mm}|=\big|\frac{s_m}{B^2}\sum_{p\ne q,\,p,q\ne m}s_ps_q+\frac{2+\kappa_{\mathrm{feat}}}{B^2}s_m\sum_{q\ne m}s_q^2+\frac{2(2+\kappa_{\mathrm{feat}})}{B^2}s_m^2\sum_{q\ne m}s_q+\frac{\E[(\omega_{i,m}^2-1)^3]}{B^2}s_m^3\big|\le\frac{P+\kappa_{\mathrm{feat}}}{B}M_0A_m+\frac{2(P-1)+\widetilde\kappa_6}{B}M_0C_m\le\frac{\max(P+\kappa_{\mathrm{feat}},\,2(P-1)+\widetilde\kappa_6)}{B}M_0(\E[H_{B_t}^2])_{mm}$ for every $m$, so $|\mathrm{BS}^\star(\Sigma_0^\star)|\le\big(3+\frac{\max(P+\kappa_{\mathrm{feat}},\,2(P-1)+\widetilde\kappa_6)}{B}\big)M_0$ and $|\mathrm{BS}^\star(\Sigma_0^\star)|/(2(1-\beta)/\eta^\star)\le O(\eta^\star)=O(\sqrt{\lambda_k})\to0$ by Theorem~\ref{thm:G5}(3).
\end{proof}

\begin{proof}[Proof of Theorem~\ref{thm:slow-fast-saturation}]
(1) For $y\ge y^\star$, $\eta\ge\eta^\star(y)$ and $\rho(\mathcal T_y)\ge1$: by Theorem~\ref{thm:G6}(1)--(2) when $\eta<\eta_{\mathrm{det}}(y)$, and when $\eta\ge\eta_{\mathrm{det}}(y)$ because $\mathcal T_y\succeq\mathcal T_{0,y}$ on the cone and $\rho(\mathcal T_{0,y})=\rho(A_{0,y})^2\ge1$ (proof of Theorem~\ref{thm:G6}(3)); so Corollary~\ref{cor:msinstab} gives $\E\|w_t\|^2+\E\|w_{t-1}\|^2\ge c(y)t$ for $t\ge1$; summing over $t=1,\dots,K$ gives $\frac1K\sum_{t=1}^K\E\|w_t\|^2\ge\frac14c(y)(K+1)-\frac{\E\|w_0\|^2}{2K}$ and (a) follows from $\tr(\Phi(y)\E[w_tw_t^\top])\ge\lambda_{\min}(\Phi(y))\E\|w_t\|^2$. For $y\in[y_{\min},y^\star)$, $\rho(\mathcal T_y)<1$ gives $\E[w_tw_t^\top]\to\Sigma_0(y)$. By perturbation theory of simple isolated eigenvalues of the $C^1$ self-adjoint family $\mathcal L_0(y)^{-1/2}\mathcal B_\beta(y)\mathcal L_0(y)^{-1/2}$, $u(y)\triangleq 1-\eta/\eta^\star(y)$, $\Sigma_0^\star(y)$ and the reduced resolvent $\widetilde R_\perp(y)=\sigma_\xi^{-2}R_\perp(y)$ of \eqref{eq:sigma0-pole} are $C^1$ in a two-sided neighbourhood of $y^\star$, with $u(y^\star)=0$ and $u'(y^\star)=g_0<0$. Taking the trace of \eqref{eq:sigma0-pole} against $\frac12\Phi(y)$ gives \eqref{eq:slow-drift}. Since $G_{\mathrm{slow}}(y_{\min})<0$ and $\lim_{y\uparrow y^\star}G_{\mathrm{slow}}(y)=+\infty$, a root $y_{\mathrm{eq}}\in(y_{\min},y^\star)$ exists. With $\varepsilon\triangleq\sigma_\xi^2/v_{\mathrm{PS}}$, multiplying $G_{\mathrm{slow}}(y)=0$ by $u(y)/v_{\mathrm{PS}}$ gives $\mathcal F(u,\varepsilon)\triangleq u-\varepsilon\big[\bar c(y(u))+\frac u2\tr(\Phi(y(u))\widetilde R_\perp(y(u)))\big]=0$, which is $C^1$ in a two-sided neighbourhood of $(0,0)$ with $\mathcal F(0,0)=0$, $\partial_u\mathcal F(0,0)=1$ and $\bar c(y^\star)=\bar c_\star>0$; the implicit function theorem gives the unique root $u_{\mathrm{eq}}=\bar c_\star\varepsilon+O(\varepsilon^2)$, at which $G'_{\mathrm{slow}}(y_{\mathrm{eq}})=\frac{|g_0|\bar c_\star\sigma_\xi^2}{u_{\mathrm{eq}}^2}(1+O(u_{\mathrm{eq}}))>0$, and $\frac{\eta^\star-\eta}{\eta}=\frac{u}{1-u}$ together with $1-r(y_{\mathrm{eq}})=d_\star u_{\mathrm{eq}}+O(u_{\mathrm{eq}}^2)$ gives \eqref{eq:saturation-law}. For (c), when $w_0=w_{-1}=0$ and $\Phi\succ0$, $\E[w_tw_t^\top]\preceq\Sigma_0(y_{\min})$ for all $t\ge0$ at $y_{\min}$ (by positivity of $\mathcal T_{y_{\min}}^s$), so $G_{\mathrm{slow},K}(y_{\min})\le G_{\mathrm{slow}}(y_{\min})<0$, while $G_{\mathrm{slow},K}(y)>0$ on $[y^\star,y_{\max}]$ by (a); continuity of $G_{\mathrm{slow},K}$ on $[y_{\min},y_{\max}]$ implies that every root of $G_{\mathrm{slow},K}$ lies in $(y_{\min},y^\star)$. Since $\sup_{y\in[y_-,y_+]}\rho(\mathcal T_y)<1$, $(G_{\mathrm{slow},K},G'_{\mathrm{slow},K})\to(G_{\mathrm{slow}},G'_{\mathrm{slow}})$ uniformly on $[y_-,y_+]$ as $K\to\infty$, giving a unique root $y_{\mathrm{eq},K}\in[y_-,y_+]$ for large $K$ with $y_{\mathrm{eq},K}\to y_{\mathrm{eq}}$. By Theorem~\ref{thm:burnin}(1), for any $\mathcal F_{k-1}$-measurable initial state $Q_{k,0}$ at epoch $k$, $\E[\frac1K\sum_{t=1}^Kg_{y,t}(w_{k,t},y_k)\mid\mathcal F_{k-1}]=G_{\mathrm{slow},K}(y_k)+b_{k,K}$ with $|b_{k,K}|=\big|\frac1{2K}\sum_{t=1}^K\tr(\Phi(y_k)(\mathcal T_{y_k}^t(\E[Q_{k,0}\mid\mathcal F_{k-1}]))_{11})\big|\le\frac{C_+R_0^2}{K}$. By the mean-value theorem on $[y_-,y_+]$, $G_{\mathrm{slow},K}(y_k)=\bar G'_k(y_k-y_{\mathrm{eq},K})$ with $\bar G'_k\in[m_0,L_K]$, so for $0<\eta_y\le1/L_K$ one has $|a_k|\triangleq|y_k-y_{\mathrm{eq},K}-\eta_yG_{\mathrm{slow},K}(y_k)|\le(1-\eta_ym_0)|y_k-y_{\mathrm{eq},K}|$; applying $(a_k-\eta_yb_{k,K})^2\le(1+\eta_ym_0)a_k^2+(1+\frac1{\eta_ym_0})\eta_y^2b_{k,K}^2\le(1-\eta_ym_0)(y_k-y_{\mathrm{eq},K})^2+\frac{2\eta_y}{m_0}b_{k,K}^2$ and non-expansiveness of $\Pi_{[y_-,y_+]}$ yields $\E[(y_{k+1}-y_{\mathrm{eq},K})^2\mid\mathcal F_{k-1}]\le(1-\eta_ym_0)(y_k-y_{\mathrm{eq},K})^2+\eta_y^2s_K^2+\frac{2\eta_yC_+^2R_0^4}{m_0K^2}$ and the stated bound.

(2) Applying \eqref{eq:slow-drift} with $\Phi=2I$ and $v_{\mathrm{PS}}=\gamma$ to the secant family on $(\gamma^\star,\gamma_{\max}]$ gives $F(\gamma)\triangleq\tr\Sigma_0^{\mathrm{sec}}(\gamma)-\gamma=\frac{\eta\,\tr(\Sigma_0^\star(\gamma)\Sigma_\xi)}{2(1-\beta)\tr(H_{\mathrm{sec}}(\gamma)\Sigma_0^\star(\gamma)^2)(1-\eta/\eta^\star_{\mathrm{sec}}(\gamma))}+\tr R_\perp^{\mathrm{sec}}(\gamma)-\gamma$, which tends to $+\infty$ as $\gamma\downarrow\gamma^\star$ and is $\le0$ at $\gamma_{\max}$; solving $(1-\eta/\eta^\star_{\mathrm{sec}}(\gamma))F(\gamma)=0$ in $u=1-\eta/\eta^\star_{\mathrm{sec}}(\gamma)$ by the implicit function theorem as in (1) gives \eqref{eq:secant-escape-law}.
\end{proof}

\section{Further properties of the second-moment operator}\label{app:more}

\subsection{Momentum diagonalizes the mean-square threshold}\label{sec:diag}

We keep the setting of Theorem~\ref{thm:G6}: $H=\sum_k\lambda_ku_ku_k^\top\succ0$ with $\lambda_1\ge\dots\ge\lambda_P>0$, per-sample Hessians $H_i$ sampled with replacement in batches of size $B$, $\Delta H_i=H_i-H$, $\mathcal V(X)=\frac1B\E[\Delta H_iX\Delta H_i]$, $\kappa_\beta=\frac{1+\beta}{1-\beta}$, $\eta_{\mathrm{det}}=\frac{2(1+\beta)}{\lambda_1}$ and $f=\eta/\eta_{\mathrm{det}}\in(0,1)$. All matrices are written in the eigenbasis of $H$. Let $\mu_1>\dots>\mu_R$ be the distinct eigenvalues of $H$, $P_r$ the orthogonal projector on the eigenspace of $\mu_r$, and
\[
\Delta_{\min}\triangleq\min_{r\ne s}|\mu_r-\mu_s|\qquad(\Delta_{\min}=+\infty\text{ if }R=1),
\]
the smallest gap between \emph{distinct} eigenvalues. On $\mathcal H=(\mathbb S^P,\langle X,Y\rangle_F)$ we use the orthonormal basis $\{E_{kk}\}_k\cup\{F_{jk}=(E_{jk}+E_{kj})/\sqrt2\}_{j<k}$.

\paragraph{The commutant splitting.} Let
\[
\mathcal D\triangleq\{X\in\mathbb S^P:[H,X]=0\}=\{X:X_{jk}=0\text{ if }\lambda_j\ne\lambda_k\},\qquad\mathcal O\triangleq\mathcal D^{\perp}=\{X:X_{jk}=0\text{ if }\lambda_j=\lambda_k\}.
\]
The orthogonal projection of $\mathcal H$ onto $\mathcal D$ is the pinching $\mathcal P(X)=\sum_rP_rXP_r$, a positive map. When the spectrum is simple, $\mathcal D$ is the set of diagonal matrices. Since $\mathcal L_{\eta,\beta}$ multiplies entry $(j,k)$ by $\ell_{jk}(\eta,\beta)$, both $\mathcal L_{\eta,\beta}$ and $\mathcal N_{\eta,\beta}=\eta\mathcal L_{\eta,\beta}^{-1}$ leave $\mathcal D$ and $\mathcal O$ invariant. On $\mathcal D$ the double commutator vanishes and the weight is constant on each eigenspace block,
\begin{equation}\label{eq:diag-wD}
\ell_{jk}=\ell_r\triangleq2\mu_r\Big(1-\frac{\eta\mu_r}{2(1+\beta)}\Big)\in\big[2(1-f)\mu_r,\,2\mu_r\big]\qquad(\lambda_j=\lambda_k=\mu_r),
\end{equation}
while on $\mathcal O$
\begin{equation}\label{eq:diag-wO}
\ell_{jk}=\Big[\lambda_j+\lambda_k-\frac{\eta\lambda_j\lambda_k}{1+\beta}\Big]+\frac{\eta\beta(\lambda_j-\lambda_k)^2}{(1-\beta)^2(1+\beta)}\ \ge\ (1-f)(\lambda_j+\lambda_k)+\frac{\eta\beta\Delta_{\min}^2}{(1-\beta)^2(1+\beta)},
\end{equation}
because $\frac{\eta\lambda_j\lambda_k}{1+\beta}=\frac{2f\lambda_j\lambda_k}{\lambda_1}\le f(\lambda_j+\lambda_k)$ (as $\lambda_j\lambda_k\le\lambda_1\lambda_k$ and $\lambda_j\lambda_k\le\lambda_1\lambda_j$). Write the blocks of $\mathcal V$ (self-adjoint on $\mathcal H$) as $\mathcal V_{DD}=\mathcal P\mathcal V|_{\mathcal D}$, $\mathcal V_{OD}=(I-\mathcal P)\mathcal V|_{\mathcal D}$, $\mathcal V_{DO}=\mathcal V_{OD}^{*}$, $\mathcal V_{OO}=(I-\mathcal P)\mathcal V|_{\mathcal O}$; norms of blocks are operator norms on $\mathcal H$ and do not depend on $(\eta,\beta)$. For a simple spectrum $\mathcal V_{DD}$ is the $P\times P$ matrix
\begin{equation}\label{eq:diag-Gamma}
\Gamma_{kl}\triangleq\langle E_{kk},\mathcal V(E_{ll})\rangle_F=(\mathcal V(E_{ll}))_{kk}=\frac1B\E\big[(u_k^\top\Delta H_iu_l)^2\big]\ \ge0,
\end{equation}
symmetric and entrywise non-negative. The \emph{diagonal part} of $\mathcal K_{\eta,\beta}$ is
\[
\mathcal K^{D}_{\eta,\beta}\triangleq\frac{\eta}{1-\beta}\mathcal L_{\eta,\beta}^{-1}\mathcal P\mathcal V\Big|_{\mathcal D}\quad\text{on }\mathcal D,\qquad\text{for a simple spectrum the matrix }\frac{\eta}{1-\beta}\operatorname{diag}(1/\ell_{kk})\,\Gamma .
\]
By Theorem~\ref{thm:G6}(1), $\mathcal K_{\eta,\beta}$ is similar to $\widetilde{\mathcal K}=\frac1{1-\beta}\mathcal N^{1/2}\mathcal V\mathcal N^{1/2}$, and in a basis adapted to $\mathcal D\oplus\mathcal O$
\begin{equation}\label{eq:diag-blocks}
\widetilde{\mathcal K}=\begin{pmatrix}\mathsf A&\mathsf B\\\mathsf B^{*}&\mathsf C\end{pmatrix},\quad
\mathsf A=\tfrac1{1-\beta}\mathcal N_D^{1/2}\mathcal V_{DD}\mathcal N_D^{1/2},\ \
\mathsf B=\tfrac1{1-\beta}\mathcal N_D^{1/2}\mathcal V_{DO}\mathcal N_O^{1/2},\ \
\mathsf C=\tfrac1{1-\beta}\mathcal N_O^{1/2}\mathcal V_{OO}\mathcal N_O^{1/2},
\end{equation}
where $\mathcal N_D,\mathcal N_O$ are the restrictions of $\mathcal N_{\eta,\beta}$; $\mathsf A$ is similar to $\mathcal K^D_{\eta,\beta}$.

\paragraph{Commutator-energy coordinates.} On $\mathcal O$ let $D_\Delta$ multiply entry $(j,k)$ by $|\lambda_j-\lambda_k|>0$, and on $\mathcal D$ let $\Lambda_D^{-1/2}(X)\triangleq H^{-1/4}XH^{-1/4}$ (the matrix $\Lambda^{-1/2}$ for a simple spectrum). For $X\in\mathcal O$ put $z=D_\Delta X$, i.e.\ $z_{jk}=|\lambda_j-\lambda_k|X_{jk}$, so that $\|z\|_F^2=\|[H,X]\|_F^2$. In the Rayleigh denominator of Theorem~\ref{thm:G6}(2) the double-commutator part of the off-diagonal weight is then isotropic:
\[
\sum_{\lambda_j\ne\lambda_k}\frac{\ell_{jk}(\eta,\beta)}{\eta}X_{jk}^2=\sum_{\lambda_j\ne\lambda_k}\Big(\frac{\lambda_j+\lambda_k}{\eta}-\frac{\lambda_j\lambda_k}{1+\beta}\Big)X_{jk}^2+\frac{\beta}{(1-\beta)^2(1+\beta)}\|z\|_F^2
\]
(the factor $\eta$ sits in $\ell_{jk}$ itself: $\ell_{jk}\supset\frac{\eta\beta}{(1-\beta)^2(1+\beta)}(\lambda_j-\lambda_k)^2$). Pairs inside one eigenspace stay in $\mathcal D$. The constants below are written in these coordinates:
\begin{equation}\label{eq:diag-commcoef}
\kappa_\Delta\triangleq\big\|D_\Delta^{-1}\mathcal V_{OD}\Lambda_D^{-1/2}\big\|,\qquad v_\Delta\triangleq\lambda_{\max}\big(D_\Delta^{-1}\mathcal V_{OO}D_\Delta^{-1}\big),\qquad v'_\Delta\triangleq\big\|D_\Delta^{-1}\mathcal V_{OO}D_\Delta^{-1}\big\|,
\end{equation}
and, at a given $\eta$, $b_\Delta\triangleq\|\widehat D^{1/2}\mathcal V_{DO}D_\Delta^{-1}\|$ with $\widehat D=\mathcal L_{\eta,\beta}^{-1}|_{\mathcal D}$. Each pair $(j,k)$ is charged with its own gap; neither $\Delta_{\min}$ nor $\lambda_P$ is pulled out of the off-diagonal block.

\begin{lemma}[Eigenvalue bounds for a symmetric $2\times2$ block operator]\label{lem:diag-block}
Let $M=\begin{pmatrix}A&B\\B^*&C\end{pmatrix}$ be self-adjoint on $\mathcal D\oplus\mathcal O$, $a=\lambda_{\max}(A)$, $c=\lambda_{\max}(C)$, $b=\|B\|$. Then
\[
a\ \le\ \lambda_{\max}(M)\ \le\ \frac{a+c}2+\sqrt{\Big(\frac{a-c}2\Big)^2+b^2},
\]
and if $a>c$ then $\lambda_{\max}(M)\le a+\frac{b^2}{a-c}$; in particular $\lambda_{\max}(M)\le a+b^2/(a-\|C\|)$ whenever $a>\|C\|$.
\end{lemma}
\begin{proof}
The lower bound is the Rayleigh quotient of $M$ at $(x,0)$ with $x$ a top eigenvector of $A$ (Cauchy interlacing). Let $\mu=\lambda_{\max}(M)$ with eigenvector $(x,y)$. If $\mu\le c$, the upper bound holds because its right side is at least $\max(a,c)$. If $\mu>c$, then $(\mu-C)$ is invertible, $y=(\mu-C)^{-1}B^*x$, and $Ax+By=\mu x$ gives $\big(A+B(\mu-C)^{-1}B^*\big)x=\mu x$ with $x\ne0$ (otherwise $y=0$). Taking the Rayleigh quotient at $x$ and using $(\mu-C)^{-1}\preceq(\mu-c)^{-1}I$,
$\mu\le a+\frac{\|B^*x\|^2}{(\mu-c)\|x\|^2}\le a+\frac{b^2}{\mu-c}$, that is $(\mu-a)(\mu-c)\le b^2$, whose largest solution is the stated bound. If $a>c$, then $\mu\ge a>c$ and $\mu-a\le b^2/(\mu-c)\le b^2/(a-c)$. Finally $c\le\|C\|$.
\end{proof}

\begin{theorem}[Two-sided bound on $\rho(\mathcal K_{\eta,\beta})$ by its diagonal part]\label{thm:diag-K}
Let $\beta\in(0,1)$ and $\eta\in(0,\eta_{\mathrm{det}})$. Write $\widehat D=\mathcal L_{\eta,\beta}^{-1}|_{\mathcal D}$, $r_D\triangleq\lambda_{\max}(\widehat D^{1/2}\mathcal V_{DD}\widehat D^{1/2})$, $\widehat\omega\triangleq\max_{\lambda_j\ne\lambda_k}1/\ell_{jk}(\eta,\beta)$ and
\[
\varepsilon_{\eta,\beta}\triangleq\frac{(1-\beta)^2(1+\beta)}{\eta\beta\Delta_{\min}^2}.
\]
(1) $\mathcal K^D_{\eta,\beta}$ preserves $\mathcal D\cap\mathbb S^P_+$, and $\rho(\mathcal K^D_{\eta,\beta})=\lambda_{\max}(\mathsf A)=\frac{\eta}{1-\beta}r_D$. With $a=\rho(\mathcal K^D_{\eta,\beta})$, $b=\|\mathsf B\|$ and $c=\lambda_{\max}(\mathsf C)\le\|\mathsf C\|$,
\[
\rho(\mathcal K^D_{\eta,\beta})\le\rho(\mathcal K_{\eta,\beta})\le\frac{a+c}2+\sqrt{\Big(\frac{a-c}2\Big)^2+b^2},\qquad
\rho(\mathcal K_{\eta,\beta})\le\rho(\mathcal K^D_{\eta,\beta})+\frac{\|\mathsf B\|^2}{\rho(\mathcal K^D_{\eta,\beta})-\|\mathsf C\|}\ \text{ if }\rho(\mathcal K^D_{\eta,\beta})>\|\mathsf C\|.
\]
(2) $\widehat\omega\le\varepsilon_{\eta,\beta}$, $\|\mathsf B\|^2\le\big(\frac{\eta}{1-\beta}\big)^2\widehat\omega\,\|\widehat D^{1/2}\mathcal V_{DO}\|^2$ and $\|\mathsf C\|\le\frac{\eta}{1-\beta}\widehat\omega\,\|\mathcal V_{OO}\|$.
(3) If $r_D>\widehat\omega\|\mathcal V_{OO}\|$, the prefactors $\eta/(1-\beta)$ cancel in the relative error:
\begin{equation}\label{eq:diag-rel}
0\le\frac{\rho(\mathcal K_{\eta,\beta})-\rho(\mathcal K^D_{\eta,\beta})}{\rho(\mathcal K^D_{\eta,\beta})}\le\frac{\widehat\omega\,\|\widehat D^{1/2}\mathcal V_{DO}\|^2}{r_D\big(r_D-\widehat\omega\|\mathcal V_{OO}\|\big)} .
\end{equation}
With $\rho_V\triangleq\lambda_{\max}(\mathcal V_{DD})\ge\frac{\E\|\Delta H_i\|_F^2}{BP}>0$, $r_D\ge\rho_V/(2\lambda_1)$ and $\|\widehat D^{1/2}\mathcal V_{DO}\|^2\le\|\mathcal V_{DO}\|^2/(2(1-f)\lambda_P)$, so that whenever $2\lambda_1\varepsilon_{\eta,\beta}\|\mathcal V_{OO}\|<\rho_V$,
\[
0\le\frac{\rho(\mathcal K_{\eta,\beta})-\rho(\mathcal K^D_{\eta,\beta})}{\rho(\mathcal K^D_{\eta,\beta})}\le\frac{2\lambda_1^2\,\varepsilon_{\eta,\beta}\,\|\mathcal V_{DO}\|^2}{(1-f)\lambda_P\,\rho_V\big(\rho_V-2\lambda_1\varepsilon_{\eta,\beta}\|\mathcal V_{OO}\|\big)} .
\]
(3$'$) (Commutator-energy form.) If $\eta\beta r_D>(1-\beta)^2(1+\beta)v'_\Delta$, then
\[
0\le\frac{\rho(\mathcal K_{\eta,\beta})-\rho(\mathcal K^D_{\eta,\beta})}{\rho(\mathcal K^D_{\eta,\beta})}\le\frac{(1+\beta)(1-\beta)^2\,b_\Delta^2}{\beta\,\eta\,r_D\big(r_D-(1-\beta)^2(1+\beta)v'_\Delta/(\beta\eta)\big)} .
\]
(4) At fixed $f$, $\eta=2f(1+\beta)/\lambda_1$ and $\varepsilon_{\eta,\beta}=\frac{(1-\beta)^2\lambda_1}{2f\beta\Delta_{\min}^2}$, so as $\beta\to1$
\[
\rho(\mathcal K^D_{\eta,\beta})\ge\frac{f\kappa_\beta\rho_V}{\lambda_1^2}\to\infty,\qquad
\frac{\rho(\mathcal K_{\eta,\beta})-\rho(\mathcal K^D_{\eta,\beta})}{\rho(\mathcal K^D_{\eta,\beta})}=O\Big(\frac{(1-\beta)^2}{\beta\Delta_{\min}^2}\Big),\qquad
\rho(\mathcal K_{\eta,\beta})-\rho(\mathcal K^D_{\eta,\beta})=O\Big(\frac{1-\beta}{\beta\Delta_{\min}^2}\Big),
\]
with the constants of (3) or (3$'$); with (3$'$) the relative error is at most $\frac{(1-\beta)^2\lambda_1b_\Delta^2}{2f\beta\,r_D\,(r_D-(1-\beta)^2\lambda_1v'_\Delta/(2f\beta))}$, where $b_\Delta$ and $r_D$ stay bounded at fixed $f$ and no $\Delta_{\min}$ appears.
(5) If $\mathcal V_{DO}=0$ ($\mathcal V$ maps $\mathcal D$ into $\mathcal D$), then $\rho(\mathcal K_{\eta,\beta})=\max\big(\rho(\mathcal K^D_{\eta,\beta}),\lambda_{\max}(\mathsf C)\big)$ exactly, for every $\beta\in[0,1)$.
\end{theorem}
\begin{proof}
(1) On $\mathcal D$, $\mathcal L^{-1}$ multiplies the block $P_rXP_r$ by $1/\ell_r>0$ by \eqref{eq:diag-wD}, so $\mathcal K^D$ is a composition of the positive maps $\mathcal V$, $\mathcal P$ and a positive blockwise scaling, and preserves $\mathcal D\cap\mathbb S^P_+$, a proper cone in $\mathcal D$. It is similar to the self-adjoint $\mathsf A$ (through $\mathcal N_D^{1/2}$), so its spectrum is real, and by the Krein--Rutman theorem in finite dimension its spectral radius is an eigenvalue, hence $\rho(\mathcal K^D)=\lambda_{\max}(\mathsf A)$; $\mathsf A=\frac{\eta}{1-\beta}\widehat D^{1/2}\mathcal V_{DD}\widehat D^{1/2}$ gives the last expression. By Theorem~\ref{thm:G6}(1), $\rho(\mathcal K)=\lambda_{\max}(\widetilde{\mathcal K})$, and Lemma~\ref{lem:diag-block} applied to \eqref{eq:diag-blocks} gives the bounds.
(2) The first claim is \eqref{eq:diag-wO}. $\mathcal N_O=\eta\,\mathcal L^{-1}|_{\mathcal O}$ is diagonal in the basis $\{F_{jk}\}_{\lambda_j\ne\lambda_k}$ with entries at most $\eta\widehat\omega$, so $\|\mathcal N_O^{1/2}\|\le(\eta\widehat\omega)^{1/2}$ and $\|\mathcal N_O\|\le\eta\widehat\omega$; with $\mathcal N_D^{1/2}=\eta^{1/2}\widehat D^{1/2}$ the two bounds follow from submultiplicativity.
(3) Insert (2) and $a=\frac{\eta}{1-\beta}r_D$ in the second bound of (1) and divide by $a$. For the explicit form: $I\in\mathcal D$ and $\langle I,\mathcal V(I)\rangle_F=\frac1B\E\|\Delta H_i\|_F^2$, so $\rho_V\ge\frac1B\E\|\Delta H_i\|_F^2/\|I\|_F^2$. If $y$ is a unit top eigenvector of $\mathcal V_{DD}$ and $x=\widehat D^{-1/2}y$, the Rayleigh quotient of $\widehat D^{1/2}\mathcal V_{DD}\widehat D^{1/2}$ at $x$ is $\rho_V/\langle y,\widehat D^{-1}y\rangle\ge\rho_V\min\widehat D$, and $\min\widehat D\ge1/(2\lambda_1)$, $\max\widehat D\le1/(2(1-f)\lambda_P)$ by \eqref{eq:diag-wD}; the right side of \eqref{eq:diag-rel} is increasing in $\widehat\omega$ and $\widehat\omega\le\varepsilon_{\eta,\beta}$.
(3$'$) By \eqref{eq:diag-wO}, $\mathcal N_O$ multiplies entry $(j,k)$ by $\eta/\ell_{jk}\le q\,(\lambda_j-\lambda_k)^{-2}$ with $q=\frac{(1-\beta)^2(1+\beta)}{\beta}$. For diagonal $0\le\mathsf a\le\mathsf b$ entrywise and any $M$, $\|M\mathsf a\|=\|M\mathsf b(\mathsf a/\mathsf b)\|\le\|M\mathsf b\|$ and $\|\mathsf aM\mathsf a\|\le\|\mathsf bM\mathsf b\|$, so $\|\mathsf B\|^2\le\frac{\eta q}{(1-\beta)^2}b_\Delta^2=\frac{\eta(1+\beta)}{\beta}b_\Delta^2$ and $\|\mathsf C\|\le\frac{q}{1-\beta}v'_\Delta=\frac{1-\beta^2}{\beta}v'_\Delta$. Insert in the second bound of (1) with $a=\frac{\eta}{1-\beta}r_D$ and divide by $a$.
(4) Substitute $\eta=2f(1+\beta)/\lambda_1$ in (3) or (3$'$); $\rho(\mathcal K^D)=\frac{\eta}{1-\beta}r_D\ge\frac{2f(1+\beta)}{\lambda_1(1-\beta)}\frac{\rho_V}{2\lambda_1}$. The absolute error is the relative one times $\rho(\mathcal K^D)\le\frac{\eta}{1-\beta}\frac{\|\mathcal V_{DD}\|}{2(1-f)\lambda_P}=O(\kappa_\beta)$.
(5) With $\mathsf B=0$, $\widetilde{\mathcal K}$ is block diagonal.
\end{proof}

At fixed $f$ the stationary regime is long lost as $\beta\to1$ ($\rho(\mathcal K)\to\infty$), so (4) is a statement about the operator, not about the threshold. The threshold itself sits at $\eta=\Theta(1-\beta)$, where $\varepsilon_{\eta,\beta}=\Theta(1-\beta)$ and the relative error loses one power of $1-\beta$. The cleanest way to see it is the $\eta$-free form of Theorem~\ref{thm:G6}(3): with $\mathcal L_0(X)=HX+XH$ and $\mathcal B_\beta$ as there, let
\[
\mathcal G_\beta\triangleq\mathcal L_0^{-1/2}\mathcal B_\beta\mathcal L_0^{-1/2},\qquad
\widetilde{\mathcal V}(X)\triangleq\mathcal P\big(H^{-1/4}\mathcal V(H^{-1/4}XH^{-1/4})H^{-1/4}\big)\ \text{on }\mathcal D,\qquad
\mathcal H_D(X)\triangleq HX\ \text{on }\mathcal D .
\]
For a simple spectrum $\widetilde{\mathcal V}$ is the matrix $\widetilde\Gamma=\Lambda^{-1/2}\Gamma\Lambda^{-1/2}$, $\widetilde\Gamma_{kl}=\frac{\E[(u_k^\top\Delta H_iu_l)^2]}{B\sqrt{\lambda_k\lambda_l}}$, and $\mathcal H_D=\Lambda$. In the splitting $\mathcal D\oplus\mathcal O$,
\begin{equation}\label{eq:diag-G}
\mathcal G_\beta=\begin{pmatrix}\frac1{2(1+\beta)}\big(\mathcal H_D+\kappa_\beta\widetilde{\mathcal V}\big)&\frac1{1-\beta}\mathcal L_{0,D}^{-1/2}\mathcal V_{DO}\mathcal L_{0,O}^{-1/2}\\[2pt]
\frac1{1-\beta}\mathcal L_{0,O}^{-1/2}\mathcal V_{OD}\mathcal L_{0,D}^{-1/2}&\mathcal S_O+\frac1{1-\beta}\mathcal L_{0,O}^{-1/2}\mathcal V_{OO}\mathcal L_{0,O}^{-1/2}\end{pmatrix},
\end{equation}
where $\mathcal S_O$ multiplies entry $(j,k)$, $\lambda_j\ne\lambda_k$, by $\big[\frac{\lambda_j\lambda_k}{1+\beta}-\frac{\beta(\lambda_j-\lambda_k)^2}{(1-\beta)^2(1+\beta)}\big]/(\lambda_j+\lambda_k)$. (On $\mathcal D$, $\mathcal L_0$ multiplies block $r$ by $2\mu_r$, so $\mathcal L_{0,D}^{-1/2}=2^{-1/2}H^{-1/4}\cdot H^{-1/4}$ and $\mathcal L_{0,D}^{-1/2}\big(\frac{H\cdot H}{1+\beta}+\frac{\mathcal P\mathcal V}{1-\beta}\big)\mathcal L_{0,D}^{-1/2}=\frac{1}{2(1+\beta)}(\mathcal H_D+\kappa_\beta\widetilde{\mathcal V})$.) Write $\mathcal G_{DD},\mathcal G_{DO},\mathcal G_{OO}$ for the blocks, $g_D\triangleq\lambda_{\max}(\mathcal G_{DD})$, and define the \emph{diagonal threshold}
\begin{equation}\label{eq:diag-etaD}
\eta_D(\beta)\triangleq\frac1{g_D}=\frac{2(1+\beta)}{\lambda_{\max}\big(\mathcal H_D+\kappa_\beta\widetilde{\mathcal V}\big)}\qquad\Big(=\frac{2(1+\beta)}{\lambda_{\max}(\Lambda+\kappa_\beta\widetilde\Gamma)}\text{ for a simple spectrum}\Big).
\end{equation}

\begin{proposition}[The diagonal threshold as an M-matrix condition; reduction to Theorems~\ref{thm:G} and~\ref{thm:G5}]\label{prop:diag-M}
Let $\beta\in[0,1)$. (1) $\eta_D\le\eta_{\mathrm{det}}$ and, for $\eta\in(0,\eta_{\mathrm{det}})$, $\rho(\mathcal K^D_{\eta,\beta})<1\iff\eta<\eta_D$. For a simple spectrum the following are equivalent: $\rho(\mathcal K^D_{\eta,\beta})<1$; $\operatorname{diag}(\ell_{kk})-\frac{\eta}{1-\beta}\Gamma$ is a non-singular M-matrix; $I-\frac{\eta}{2(1+\beta)}\big(\Lambda+\kappa_\beta\widetilde\Gamma\big)$ is a non-singular M-matrix; $\eta<\eta_D$. In terms of the normalized step variances $c_k=\frac{2\lambda_k}{\eta(1+\beta)}\gamma_{0,k}$, the diagonal equations read $\big[I-\frac{\eta}{2(1+\beta)}(\Lambda+\kappa_\beta\Gamma\Lambda^{-1})\big]c=\frac{\sigma^2_{\xi}}{1-\beta^2}+(\text{terms in the off-diagonal of }\Sigma_0)$, the last vanishing when $\mathcal V_{DO}=0$.
(2) If $\widetilde\Gamma=\operatorname{diag}(\gamma_k)+gg^\top$ with $g_k>0$, then $\eta_D=2(1+\beta)/\mu_\star$, where $\mu_\star$ is the unique $\mu>\max\big(0,\max_k(\lambda_k+\kappa_\beta\gamma_k)\big)$ at which \eqref{eq:diag-secular} holds with $\eta=2(1+\beta)/\mu$:
\begin{equation}\label{eq:diag-secular}
\frac{\eta}{2(1-\beta)}\sum_{k=1}^P\frac{g_k^2}{1-\frac{\eta\lambda_k}{2(1+\beta)}-\frac{\eta\gamma_k}{2(1-\beta)}}=1 .
\end{equation}
For the random features of Theorem~\ref{thm:G}, $\Gamma_{kl}=\frac1B\big(\lambda_k\lambda_l+(1+\kappa_{\mathrm{feat}})\lambda_k^2\delta_{kl}\big)$, so $g_k^2=\lambda_k/B$, $\gamma_k=(1+\kappa_{\mathrm{feat}})\lambda_k/B$, and \eqref{eq:diag-secular} is \eqref{eq:secular}; for the features of Theorem~\ref{thm:G5}, $\Gamma_{kl}=\frac1B\big(s_ks_l+(1+\kappa_{\mathrm{feat}})s_k^2\delta_{kl}\big)$, $g_k^2=s_k^2/(B\lambda_k)$, $\gamma_k=(1+\kappa_{\mathrm{feat}})s_k^2/(B\lambda_k)$, and \eqref{eq:diag-secular} is the equation of Theorem~\ref{thm:G5}(2). In both cases $\mathcal V_{DO}=0$.
\end{proposition}
\begin{proof}
(1) $\lambda_{\max}(\mathcal H_D+\kappa_\beta\widetilde{\mathcal V})\ge\langle E_{11},(\mathcal H_D+\kappa_\beta\widetilde{\mathcal V})E_{11}\rangle_F\ge\lambda_1$ because $\widetilde{\mathcal V}$ is a positive map, so $\eta_D\le\eta_{\mathrm{det}}$. As in Theorem~\ref{thm:G6}(1), $\rho(\mathcal K^D)<1\iff\lambda_{\max}(\mathsf A)<1\iff\mathcal L_D-\frac{\eta}{1-\beta}\mathcal V_{DD}\succ0$ on $\mathcal D$; on $\mathcal D$, $\mathcal L_D-\frac{\eta}{1-\beta}\mathcal V_{DD}=\mathcal L_{0,D}^{1/2}\big(I-\eta\mathcal G_{DD}\big)\mathcal L_{0,D}^{1/2}$ (the commutator is absent on $\mathcal D$), which is positive definite iff $\eta g_D<1$. For a simple spectrum, $\mathcal K^D=\frac{\eta}{1-\beta}\operatorname{diag}(1/\ell_{kk})\Gamma\ge0$ entrywise and $\operatorname{diag}(\ell_{kk})-\frac{\eta}{1-\beta}\Gamma=\operatorname{diag}(\ell_{kk})(I-\mathcal K^D)$ is a Z-matrix that is a non-singular M-matrix iff $\rho(\mathcal K^D)<1$ \citep[Ch.~6]{berman1994}; the identity
\[
\operatorname{diag}(\ell_{kk})-\frac{\eta}{1-\beta}\Gamma=2\Lambda^{1/2}\Big[I-\frac{\eta}{2(1+\beta)}\big(\Lambda+\kappa_\beta\widetilde\Gamma\big)\Big]\Lambda^{1/2}
\]
(use $\ell_{kk}=2\lambda_k-\eta\lambda_k^2/(1+\beta)$ and $\frac{\eta\kappa_\beta}{1+\beta}=\frac{\eta}{1-\beta}$) transfers the property to the symmetric Z-matrix in brackets, which is a non-singular M-matrix iff it is positive definite. The equation for $c$ is the $(k,k)$ entry of Theorem~\ref{thm:F}(2) divided by $\eta(1+\beta)$, as in the proof of Theorem~\ref{thm:G}; $(\mathcal V(\Sigma_0))_{kk}=\sum_l\Gamma_{kl}\gamma_{0,l}+\langle E_{kk},\mathcal V_{DO}(\Sigma_0)_{\mathcal O}\rangle_F$.
(2) By the rank-one determinant lemma, $\det(\mu I-\Lambda-\kappa_\beta\operatorname{diag}(\gamma)-\kappa_\beta gg^\top)=\prod_k(\mu-\lambda_k-\kappa_\beta\gamma_k)\big(1-\kappa_\beta\sum_kg_k^2/(\mu-\lambda_k-\kappa_\beta\gamma_k)\big)$; the sum decreases from $+\infty$ to $0$ on $(\max_k(\lambda_k+\kappa_\beta\gamma_k),\infty)$ ($g_k>0$ for every $k$), and by interlacing the top eigenvalue lies there, so it is the unique root of the bracket on that interval. At $\mu=2(1+\beta)/\eta$, multiplying numerator and denominator by $\eta/(2(1+\beta))$ and using $\frac{\eta\kappa_\beta}{2(1+\beta)}=\frac{\eta}{2(1-\beta)}$ gives \eqref{eq:diag-secular}. The entries of $\Gamma$ are $(\mathcal V(E_{ll}))_{kk}$ from the formulas for $\mathcal V$ in the proofs of Theorems~\ref{thm:G} and~\ref{thm:G5}; there $(\mathcal V(X))_{jk}=\frac{s_js_k}{B}X_{jk}$ for $j\ne k$ (with $s=\lambda$ in Theorem~\ref{thm:G}), which vanishes for diagonal $X$, so $\mathcal V_{DO}=0$. With $g_k^2=\lambda_k/B$ and $\gamma_k=(1+\kappa_{\mathrm{feat}})\lambda_k/B$, the denominator of \eqref{eq:diag-secular} is $1-\frac{\eta\lambda_k}{2(1+\beta)}\alpha_{\mathrm{diag}}$.
\end{proof}

\begin{theorem}[Momentum diagonalizes the threshold]\label{thm:diag-threshold}
Let $\beta\in(0,1)$ and $\eta^\star=\etacrit^{\mathrm{EoSS},P}(\beta,B)$ be the mean-square ceiling of Theorem~\ref{thm:G6}. Then $\eta^\star=1/\lambda_{\max}(\mathcal G_\beta)$, $\eta^\star\le\eta_D$, and the gap $1/\eta^\star-1/\eta_D\ge0$ has the following upper bounds, listed from the tightest; each holds whenever its condition holds.
(1) (Eigenvector split; the tightest of the bounds.) Let $u$ be a unit top eigenvector of $\mathcal G_{DD}$, $\gamma_D\triangleq g_D-\lambda_2(\mathcal G_{DD})$ its spectral gap ($\gamma_D=+\infty$ if $\dim\mathcal D=1$), $P_\perp=I-u\langle u,\cdot\rangle$ on $\mathcal D$, and $\mathsf E\in\{\mathsf E_\Delta,\bar{\mathsf E}\}$ one of the majorants of (2) (defined when $\theta_\beta<1$, resp.\ $\bar\theta<1$). If
\begin{equation}\label{eq:diag-eigcond}
\gamma_D+\langle u,\mathsf Eu\rangle>\|P_\perp\mathsf EP_\perp\|,
\end{equation}
then
\[
0\le\frac1{\eta^\star}-\frac1{\eta_D}\le c''_\beta\triangleq\langle u,\mathsf Eu\rangle+\frac{\|P_\perp\mathsf Eu\|^2}{\gamma_D+\langle u,\mathsf Eu\rangle-\|P_\perp\mathsf EP_\perp\|}.
\]
If $\lambda_{\max}(\widetilde{\mathcal V})$ is simple with gap $\mathrm{gap}(\widetilde{\mathcal V})>0$, then $\gamma_D\ge\frac{\mathrm{gap}(\widetilde{\mathcal V})}{2(1-\beta)}-\frac{\lambda_1}{1+\beta}$, so \eqref{eq:diag-eigcond} holds for all $\beta$ close enough to $1$, and with $\mathsf E=\mathsf E_\Delta$, $c''_\beta=c_\infty+O(1-\beta)$: the leading term is exact and the bound is asymptotically sharp. The flat and weakly coupled directions of $A$ enter only through $\|P_\perp\mathsf Eu\|^2/\gamma_D=O(1-\beta)$.
(2) (Loewner majorants of the Schur term; commutator-energy form.) If $\theta_\beta\triangleq(1-\beta^2)v_\Delta/\beta<1$,
\[
0\le\frac1{\eta^\star}-\frac1{\eta_D}\le c'_\beta\triangleq\frac{1+\beta}{2\beta}\,\frac{\kappa_\Delta^2}{1-\theta_\beta}\ \xrightarrow[\beta\to1]{}\ \kappa_\Delta^2=\big\|D_\Delta^{-1}\mathcal V_{OD}\Lambda_D^{-1/2}\big\|^2 .
\]
A weighted form, defined whenever the Schur step is: with $s_{jk}=\lambda_j+\lambda_k$, $\bar c_{jk}\triangleq\frac{\beta(\lambda_j-\lambda_k)^2}{(1-\beta)^2(1+\beta)s_{jk}}+g_D-\frac{\lambda_j\lambda_k}{(1+\beta)s_{jk}}$, $\bar{\mathcal C}=\operatorname{diag}(\bar c_{jk})$ on $\mathcal O$, $\mathcal E_{OO}=\frac1{1-\beta}\mathcal L_{0,O}^{-1/2}\mathcal V_{OO}\mathcal L_{0,O}^{-1/2}$ and $\bar\theta=\lambda_{\max}(\bar{\mathcal C}^{-1/2}\mathcal E_{OO}\bar{\mathcal C}^{-1/2})<1$, $\ \frac1{\eta^\star}-\frac1{\eta_D}\le\frac{\|\bar{\mathcal C}^{-1/2}\mathcal G_{OD}\|^2}{1-\bar\theta}$. In operator form: $\mathsf E_\Delta\triangleq\alpha_\beta A^*A$ with $A=D_\Delta^{-1}\mathcal V_{OD}\Lambda_D^{-1/2}$ and $\alpha_\beta=\frac{1+\beta}{2\beta(1-\theta_\beta)}$, and $\bar{\mathsf E}\triangleq\frac1{1-\bar\theta}\mathcal G_{DO}\bar{\mathcal C}^{-1}\mathcal G_{OD}$, both positive semidefinite on $\mathcal D$, satisfy $0\preceq\mathcal G_{DO}(\mu^\star_\eta-\mathcal G_{OO})^{-1}\mathcal G_{OD}\preceq\mathsf E$ at $\mu^\star_\eta=1/\eta^\star$, and the two scalar bounds are $c'_\beta=\|\mathsf E_\Delta\|$ and $\|\bar{\mathsf E}\|$. Two further majorants: let $T=\bar{\mathcal C}^{-1/2}\mathcal E_{OO}\bar{\mathcal C}^{-1/2}$ and $T_+$ its positive part ($\mathcal V_{OO}$, hence $T$, need not be positive semidefinite, and $T$ cannot replace $T_+$). Since $\frac1{1-t}\le1+\frac{t_+}{1-\bar\theta}$ for $t\le\bar\theta<1$, and $1+\frac{t_+}{1-\bar\theta}\le\frac1{1-\bar\theta}$ when moreover $\bar\theta\ge0$, $\bar{\mathsf E}_1\triangleq\mathcal G_{DO}\bar{\mathcal C}^{-1}\mathcal G_{OD}+\frac1{1-\bar\theta}\mathcal G_{DO}\bar{\mathcal C}^{-1/2}T_+\bar{\mathcal C}^{-1/2}\mathcal G_{OD}$ satisfies $\mathcal G_{DO}(\mu^\star_\eta-\mathcal G_{OO})^{-1}\mathcal G_{OD}\preceq\bar{\mathsf E}_1$, with $\bar{\mathsf E}_1\preceq\bar{\mathsf E}$ when $\bar\theta\ge0$ (for $\bar\theta<0$ the order can fail); and whenever $g_D>\lambda_{\max}(\mathcal G_{OO})$, $\mu^\star_\eta\ge g_D$ gives $\mathcal G_{DO}(\mu^\star_\eta-\mathcal G_{OO})^{-1}\mathcal G_{OD}\preceq\mathsf E_{\mathrm{op}}\triangleq\mathcal G_{DO}(g_D-\mathcal G_{OO})^{-1}\mathcal G_{OD}$. Both may be used as $\mathsf E$ in (1); over the 42 cells of \texttt{paper/checks/diagonalization-burnin.py} (section S8) the ratio of (1) to the true gap is $1$--$17.3$ with $\bar{\mathsf E}_1$ and $1$--$18.2$ with $\mathsf E_{\mathrm{op}}$, against $1$--$17.5$ with $\bar{\mathsf E}$; the largest ratios come from a small $\gamma_D$, not from the scalar factor.
(3) (Operator form.) Whenever $g_D>\lambda_{\max}(\mathcal G_{OO})$,
\[
0\le\frac1{\eta^\star}-\frac1{\eta_D}\le\frac{\|\mathcal G_{DO}\|^2}{g_D-\lambda_{\max}(\mathcal G_{OO})} .
\]
(4) ($\Delta_{\min}$, $\lambda_P$ form.) $\|\mathcal G_{DO}\|\le\frac{\|\mathcal V_{DO}\|}{2\lambda_P(1-\beta)}$ and $g_D-\lambda_{\max}(\mathcal G_{OO})\ge\frac{\beta\Delta_{\min}^2}{2\lambda_1(1-\beta)^2(1+\beta)}-\frac{\|\mathcal V_{OO}\|}{2\lambda_P(1-\beta)}$. Hence, if $\beta\Delta_{\min}^2>(1-\beta^2)\lambda_1\|\mathcal V_{OO}\|/\lambda_P$,
\[
0\le\frac1{\eta^\star}-\frac1{\eta_D}\le c_\beta\triangleq\frac{\lambda_1(1+\beta)\|\mathcal V_{DO}\|^2}{2\lambda_P^2\big(\beta\Delta_{\min}^2-(1-\beta^2)\lambda_1\|\mathcal V_{OO}\|/\lambda_P\big)},
\]
a quantity that stays bounded as $\beta\to1$.
(5) $\frac1{\eta_D}\in\Big[\frac{\lambda_{\max}(\widetilde{\mathcal V})}{2(1-\beta)},\ \frac{\lambda_{\max}(\widetilde{\mathcal V})}{2(1-\beta)}+\frac{\lambda_1}{2(1+\beta)}\Big]$ with $\lambda_{\max}(\widetilde{\mathcal V})\ge\frac{\E\|\Delta H_i\|_F^2}{B\tr H}>0$. Consequently
\[
1\le\frac{\eta_D}{\eta^\star}\le1+\frac{2(1-\beta)c^{\mathrm{best}}_\beta}{\lambda_{\max}(\widetilde{\mathcal V})},\qquad
\eta^\star=\frac{2(1-\beta)}{\lambda_{\max}(\widetilde{\mathcal V})}\big(1+O(1-\beta)\big)\quad(\beta\to1),
\]
with $c^{\mathrm{best}}_\beta$ the least of the defined bounds (1)--(4), so the relative error of the diagonal threshold is $\frac{2(1-\beta)c_\infty}{\lambda_{\max}(\widetilde{\mathcal V})}(1+o(1))$ by Proposition~\ref{prop:diag-expansion}, and (1) attains this to first order (at most $O\big((1-\beta)\kappa_\Delta^2/\lambda_{\max}(\widetilde{\mathcal V})\big)$ by (2)) (and $O\big((1-\beta)/(\beta\Delta_{\min}^2)\big)$ in the cruder form), one power of $1-\beta$ less than the operator statement of Theorem~\ref{thm:diag-K}(4), because $\eta^\star=\Theta(1-\beta)$.
(6) If $\mathcal V_{DO}=0$, then $1/\eta^\star=\max\big(g_D,\lambda_{\max}(\mathcal G_{OO})\big)$, and $\eta^\star=\eta_D$ as soon as $\lambda_{\max}(\mathcal G_{OO})\le g_D$, in particular for all $\beta$ close enough to $1$ when $R\ge2$ (and in the settings of Theorems~\ref{thm:G} and~\ref{thm:G5} for every $\beta$, by those theorems).
\end{theorem}
\begin{proof}
For $\eta\in(0,\eta_{\mathrm{det}})$, Theorem~\ref{thm:G6} gives $\rho(\mathcal T)<1\iff\rho(\mathcal K)<1\iff\mathcal L_0-\eta\mathcal B_\beta\succ0\iff\eta\lambda_{\max}(\mathcal G_\beta)<1$ (note $\lambda_{\max}(\mathcal G_\beta)\ge g_D>0$), and $\rho(\mathcal T)\ge1$ for $\eta\ge\eta_{\mathrm{det}}$ (there $\rho(\mathcal T)\ge\rho(\mathcal T_0)=\rho(\bar F)^2\ge1$, because $\mathcal T-\mathcal T_0$ is positive, see Section~\ref{sec:burnin}, and the spectral radius is monotone on positive maps of a proper cone). Since $\lambda_{\max}(\mathcal G_\beta)\ge g_D\ge1/\eta_{\mathrm{det}}$ (Proposition~\ref{prop:diag-M}(1)), the stable set is exactly $(0,1/\lambda_{\max}(\mathcal G_\beta))$. The lower bound $1/\eta^\star\ge g_D$ is Lemma~\ref{lem:diag-block}. (3) is Lemma~\ref{lem:diag-block} applied to \eqref{eq:diag-G}. (4) $\mathcal L_{0,D}\succeq2\lambda_PI$ and $\mathcal L_{0,O}\succeq2\lambda_PI$, which gives the first bound. On $\mathcal O$, $\frac{\lambda_j\lambda_k}{\lambda_j+\lambda_k}\le\frac{\lambda_1}2$ and $\lambda_j+\lambda_k\le2\lambda_1$, so $\mathcal S_O\preceq\big(\frac{\lambda_1}{2(1+\beta)}-\frac{\beta\Delta_{\min}^2}{2\lambda_1(1-\beta)^2(1+\beta)}\big)I$, and Weyl's inequality gives $\lambda_{\max}(\mathcal G_{OO})\le\frac{\lambda_1}{2(1+\beta)}-\frac{\beta\Delta_{\min}^2}{2\lambda_1(1-\beta)^2(1+\beta)}+\frac{\|\mathcal V_{OO}\|}{2\lambda_P(1-\beta)}$, while $g_D\ge\frac{\lambda_1}{2(1+\beta)}$ by Proposition~\ref{prop:diag-M}(1). Dividing $\|\mathcal G_{DO}\|^2$ by the difference gives $c_\beta$. (2) Let $\mu=\lambda_{\max}(\mathcal G_\beta)$. Since $\mu\ge g_D\ge\frac{\lambda_1}{2(1+\beta)}\ge\frac{\lambda_j\lambda_k}{(1+\beta)s_{jk}}$, $\mu-\mathcal S_O=\bar{\mathcal C}+(\mu-g_D)I\succeq\bar{\mathcal C}\succeq\mathcal C$, where $\mathcal C$ multiplies entry $(j,k)$ by $c_{jk}=\frac{\beta(\lambda_j-\lambda_k)^2}{(1-\beta)^2(1+\beta)s_{jk}}$. If $\mathcal E_{OO}\preceq\theta\,\mathcal C'$ for $\mathcal C'\in\{\mathcal C,\bar{\mathcal C}\}$ with $\theta<1$, then $\mu-\mathcal G_{OO}\succeq\mathcal C'-\mathcal E_{OO}\succeq(1-\theta)\mathcal C'\succ0$, so $\mu>\lambda_{\max}(\mathcal G_{OO})$ and, as in the proof of Lemma~\ref{lem:diag-block}, $\mu$ is an eigenvalue of $\mathcal G_{DD}+\mathcal G_{DO}(\mu-\mathcal G_{OO})^{-1}\mathcal G_{OD}\preceq\mathcal G_{DD}+\frac{1}{1-\theta}\mathcal G_{DO}\mathcal C'^{-1}\mathcal G_{OD}$, whence $\mu\le g_D+\frac{\|\mathcal C'^{-1/2}\mathcal G_{OD}\|^2}{1-\theta}$. For $\mathcal C'=\bar{\mathcal C}$ this is the weighted form. For $\mathcal C'=\mathcal C$: $\mathcal C^{-1/2}\mathcal L_{0,O}^{-1/2}=\big(\frac{(1-\beta)^2(1+\beta)}{\beta}\big)^{1/2}D_\Delta^{-1}$, so the least admissible $\theta$ is $\frac{1}{1-\beta}\frac{(1-\beta)^2(1+\beta)}{\beta}v_\Delta=\theta_\beta$, and $\mathcal C^{-1/2}\mathcal G_{OD}=\big(\frac{1+\beta}{\beta}\big)^{1/2}D_\Delta^{-1}\mathcal V_{OD}\mathcal L_{0,D}^{-1/2}$ with $\mathcal L_{0,D}^{-1/2}=2^{-1/2}\Lambda_D^{-1/2}$, so $\|\mathcal C^{-1/2}\mathcal G_{OD}\|^2=\frac{1+\beta}{2\beta}\kappa_\Delta^2$ and $\frac{1}{1-\theta_\beta}\mathcal G_{DO}\mathcal C^{-1}\mathcal G_{OD}=\mathsf E_\Delta$. Both majorants are positive semidefinite, as is the Schur term itself ($\mu-\mathcal G_{OO}\succ0$). (1) By the proof of (2), $\mu$ is an eigenvalue of $\mathcal G_{DD}+\mathsf S$ with $0\preceq\mathsf S\preceq\mathsf E$, hence $\mu\le\lambda_{\max}(\mathcal G_{DD}+\mathsf E)$. Write $M=\mathcal G_{DD}+\mathsf E$ in the splitting $\mathcal D=\operatorname{span}(u)\oplus u^\perp$. Since $\mathcal G_{DD}u=g_Du$, the $(u,u)$ entry is $a=g_D+\langle u,\mathsf Eu\rangle$, the off-diagonal block is $P_\perp\mathsf Eu$ (as $P_\perp\mathcal G_{DD}u=0$), and the $u^\perp$ block is $P_\perp(\mathcal G_{DD}+\mathsf E)P_\perp$, whose top eigenvalue is at most $\lambda_2(\mathcal G_{DD})+\|P_\perp\mathsf EP_\perp\|=g_D-\gamma_D+\|P_\perp\mathsf EP_\perp\|\triangleq\bar c$ by Weyl's inequality. Under \eqref{eq:diag-eigcond}, $a>\bar c$, and the proof of Lemma~\ref{lem:diag-block} ($\mu-a\le b^2/(\mu-c)\le b^2/(a-\bar c)$ for $c\le\bar c<a\le\mu$) gives $\lambda_{\max}(M)\le a+\|P_\perp\mathsf Eu\|^2/(a-\bar c)$, which is the bound. For the asymptotics: $\mathcal G_{DD}=\frac{1}{2(1-\beta)}\widetilde{\mathcal V}+\frac{1}{2(1+\beta)}\mathcal H_D$ with $0\preceq\mathcal H_D\preceq\lambda_1I$, so by Weyl $\gamma_D\ge\frac{\mathrm{gap}(\widetilde{\mathcal V})}{2(1-\beta)}-\frac{\lambda_1}{1+\beta}$, and by the Davis--Kahan theorem $u=X_\infty+O(1-\beta)$ (up to sign). Since $\theta_\beta\to0$ and $\alpha_\beta\to1$, $\langle u,\mathsf E_\Delta u\rangle=\alpha_\beta\|Au\|^2=\|AX_\infty\|^2+O(1-\beta)=c_\infty+O(1-\beta)$ (Proposition~\ref{prop:diag-expansion}), $\|\mathsf E_\Delta\|=O(1)$, and the second term is $O(1-\beta)$. (5) $0\preceq\mathcal H_D\preceq\lambda_1I$ on $\mathcal D$ and Weyl's inequality give the bracket for $g_D=1/\eta_D$ (with $\frac{\kappa_\beta}{2(1+\beta)}=\frac1{2(1-\beta)}$). $\widetilde{\mathcal V}$ is a positive self-adjoint map on $\mathcal D$; its Rayleigh quotient at $H^{1/2}\in\mathcal D$ is $\tr(\mathcal V(I))/\tr H=\frac{\E\|\Delta H_i\|_F^2}{B\tr H}$. Then $\frac{\eta_D}{\eta^\star}-1=\eta_D\big(\frac1{\eta^\star}-\frac1{\eta_D}\big)\le\frac{2(1-\beta)}{\lambda_{\max}(\widetilde{\mathcal V})}c^{\mathrm{best}}_\beta$, and $\frac1{\eta^\star}=\frac{\lambda_{\max}(\widetilde{\mathcal V})}{2(1-\beta)}+O(1)$. (6) With $\mathcal G_{DO}=0$, $\mathcal G_\beta$ is block diagonal; $\lambda_{\max}(\mathcal G_{OO})\to-\infty$ like $-(1-\beta)^{-2}$ when $\mathcal O\ne\{0\}$, while $g_D=\Theta((1-\beta)^{-1})$.
\end{proof}

\begin{proposition}[First-order correction of the threshold as $\beta\to1$]\label{prop:diag-expansion}
Assume $\lambda_{\max}(\widetilde{\mathcal V})$ is a simple eigenvalue of $\widetilde{\mathcal V}$ on $\mathcal D$, with unit eigenvector $X_\infty\in\mathcal D$, and set $Y_\infty=H^{-1/4}X_\infty H^{-1/4}$. Then, as $\beta\to1$,
\[
\frac1{\eta^\star}=\frac1{\eta_D}+c_\infty+O(1-\beta),\qquad c_\infty\triangleq\sum_{\substack{j\ne k\\\lambda_j\ne\lambda_k}}\frac{(\mathcal V(Y_\infty))_{jk}^2}{(\lambda_j-\lambda_k)^2},
\qquad
\frac{\eta_D}{\eta^\star}-1=\frac{2(1-\beta)c_\infty}{\lambda_{\max}(\widetilde{\mathcal V})}+O\big((1-\beta)^2\big).
\]
In commutator-energy coordinates $c_\infty=\big\|D_\Delta^{-1}\mathcal V_{OD}\Lambda_D^{-1/2}X_\infty\big\|^2\le\kappa_\Delta^2=\lim_{\beta\to1}c'_\beta$: the limit of the scalar bound of Theorem~\ref{thm:diag-threshold}(2) is the same quadratic form evaluated in its worst direction instead of at the dominant diagonal mode $X_\infty$. In the scale of the paper, $\mu=2(1+\beta)/\eta$: $\mu^\star=\mu_D+2(1+\beta)c_\infty+O(1-\beta)$ with $\mu_D=2(1+\beta)/\eta_D=\lambda_{\max}(\mathcal H_D+\kappa_\beta\widetilde{\mathcal V})=\Theta((1-\beta)^{-1})$.
In particular the relative error of the diagonal threshold is of exact order $1-\beta$ iff $\mathcal V(Y_\infty)\notin\mathcal D$, i.e.\ iff the curvature noise maps the dominant diagonal mode off the commutant of $H$.
\end{proposition}
\begin{proof}
Let $\epsilon=1-\beta$, $\mu=\lambda_{\max}(\mathcal G_\beta)$ and $\mathsf G_1=\mathcal L_{0,D}^{-1/2}\mathcal V_{DO}\mathcal L_{0,O}^{-1/2}$, which does not depend on $\beta$; $\mathcal G_{DO}=\mathsf G_1/\epsilon$. By \eqref{eq:diag-G}, $\epsilon\mathcal G_{DD}=\frac12\widetilde{\mathcal V}+O(\epsilon)$ and $\epsilon^2\mathcal G_{OO}=-\frac{\beta}{1+\beta}\mathcal Q_\Delta+O(\epsilon)$ with $\mathcal Q_\Delta$ multiplying entry $(j,k)$ of $\mathcal O$ by $(\lambda_j-\lambda_k)^2/(\lambda_j+\lambda_k)\ge\Delta_{\min}^2/(2\lambda_1)$. By Theorem~\ref{thm:diag-threshold}, $\mu=g_D+O(1)=O(\epsilon^{-1})$, so $\epsilon^2(\mu-\mathcal G_{OO})=\frac{\beta}{1+\beta}\mathcal Q_\Delta+O(\epsilon)$ is invertible for small $\epsilon$ and, as in the proof of Lemma~\ref{lem:diag-block}, $\mu$ is an eigenvalue of $\mathcal G_{DD}+\mathsf S$ with $\mathsf S=\mathcal G_{DO}(\mu-\mathcal G_{OO})^{-1}\mathcal G_{OD}=\mathsf G_1\big(\tfrac{\beta}{1+\beta}\mathcal Q_\Delta+O(\epsilon)\big)^{-1}\mathsf G_1^*=\mathsf S_\infty+O(\epsilon)$, $\mathsf S_\infty=2\mathsf G_1\mathcal Q_\Delta^{-1}\mathsf G_1^*$. The top eigenvalue of $\mathcal G_{DD}$ is separated from the rest of its spectrum by $\frac{\mathrm{gap}(\widetilde{\mathcal V})}{2\epsilon}+O(1)$, with unit eigenvector $x_\beta=X_\infty+O(\epsilon)$. Lemma~\ref{lem:diag-block}, applied to $\mathcal G_{DD}+\mathsf S$ in the basis $(x_\beta,x_\beta^\perp)$, gives $\lambda_{\max}(\mathcal G_{DD}+\mathsf S)=g_D+\langle x_\beta,\mathsf Sx_\beta\rangle+O(\|\mathsf S\|^2\epsilon)$ and shows that all other eigenvalues of $\mathcal G_{DD}+\mathsf S$ are below $g_D-\frac{\mathrm{gap}}{2\epsilon}+O(1)<g_D\le\mu$, so $\mu$ is that top eigenvalue. Hence $\mu=g_D+\langle X_\infty,\mathsf S_\infty X_\infty\rangle+O(\epsilon)$. Finally $\langle X_\infty,\mathsf S_\infty X_\infty\rangle=2\sum_{o\in\mathcal O}\langle F_o,\mathcal V(\mathcal L_{0,D}^{-1/2}X_\infty)\rangle_F^2/(\lambda_j-\lambda_k)^2$, and $\mathcal L_{0,D}^{-1/2}X_\infty=Y_\infty/\sqrt2$, $\langle F_{jk},Z\rangle_F=\sqrt2Z_{jk}$, which gives $c_\infty$ (the sum over $j<k$ counted twice); equivalently $\langle X_\infty,\mathsf S_\infty X_\infty\rangle=2\|D_\Delta^{-1}\mathcal V_{OD}\mathcal L_{0,D}^{-1/2}X_\infty\|^2=\|D_\Delta^{-1}\mathcal V_{OD}\Lambda_D^{-1/2}X_\infty\|^2$. The last display uses $\eta_D=\frac{2(1-\beta)}{\lambda_{\max}(\widetilde{\mathcal V})}(1+O(1-\beta))$.
\end{proof}

\paragraph{The two regimes.} (a) \emph{Fixed $f=\eta/\eta_{\mathrm{det}}$} (Theorem~\ref{thm:diag-K}(4)): $\rho(\mathcal K)$ and $\rho(\mathcal K^D)$ both grow like $\kappa_\beta$, their difference is $O(1-\beta)$, and the relative error is $O((1-\beta)^2)$. This regime lies beyond the ceiling for $\beta$ near $1$. (b) \emph{At the threshold} (Theorem~\ref{thm:diag-threshold}, Proposition~\ref{prop:diag-expansion}): $\eta^\star=\Theta(1-\beta)$, $\mu_D=2(1+\beta)/\eta_D=\Theta((1-\beta)^{-1})$, and $\mu^\star=\mu_D+2(1+\beta)c_\infty+O(1-\beta)$, so the relative shift $\mu^\star/\mu_D-1=\eta_D/\eta^\star-1$ is $\Theta(1-\beta)$ when $c_\infty>0$. The loss of one power is the factor $1/\eta$ in $\varepsilon_{\eta,\beta}$ evaluated at $\eta=\eta^\star$.

\begin{remark}[Scope: the rank-one limit, degenerate, small-gap and flat modes]\label{rem:diag-scope}
(i) The rank-one secular equation of Theorem~\ref{thm:G} is not the $\beta\to1$ limit for arbitrary non-commuting per-sample Hessians, in two ways. The $\beta\to1$ limit is the diagonal M-matrix condition of Proposition~\ref{prop:diag-M} with the full non-negative matrix $\Gamma$ of \eqref{eq:diag-Gamma}; it is a rank-one secular equation only when $\widetilde\Gamma$ is diagonal plus rank one, which is what independent feature coordinates produce (Theorems~\ref{thm:G},~\ref{thm:G5}). And it is not exact at any $\beta<1$ when $c_\infty>0$: by Proposition~\ref{prop:diag-expansion}, $\eta_D/\eta^\star-1=\Theta(1-\beta)$ whenever $c_\infty>0$. What momentum does is suppress the off-commutant entries of $\Sigma_0$ through the weight $\frac{\eta\beta(\lambda_j-\lambda_k)^2}{(1-\beta)^2(1+\beta)}$, which grows faster than the $\frac{\eta}{1-\beta}$ amplification of the curvature noise. As $\beta\to1$, $\frac{2(1-\beta)}{\eta^\star}\to\lambda_{\max}(\widetilde{\mathcal V})$, which for random features is $\lambda_{\max}\big(\frac1B(\sqrt\lambda\sqrt\lambda^\top+(1+\kappa_{\mathrm{feat}})\Lambda)\big)$, the batch trace-sharpness plateau $\tr H/B$ of Corollary~\ref{cor:plateaus} up to a term of size $|1+\kappa_{\mathrm{feat}}|\lambda_1/B$.
(ii) Degenerate eigenvalues. Entries $(j,k)$ inside one eigenspace carry no double-commutator weight, so they belong to $\mathcal D$; the limit is block diagonal, with $\widetilde{\mathcal V}$ acting on block-diagonal matrices through the pinching $\mathcal P$. It is not the diagonal of $\Gamma$ alone: within an eigenspace of multiplicity $m$, $\widetilde{\mathcal V}$ acts on $\mathbb S^m$.
(iii) Small gaps. Every bound scales with $1/\Delta_{\min}^2$, and the explicit form (4) of Theorem~\ref{thm:diag-threshold} requires $\beta\Delta_{\min}^2>(1-\beta^2)\lambda_1\|\mathcal V_{OO}\|/\lambda_P$. For a cluster of eigenvalues closer than $(1-\beta)$ times the scale of the curvature noise, put the cluster into $\mathcal D$: Lemma~\ref{lem:diag-block} and Theorem~\ref{thm:diag-threshold}(3) hold verbatim for any orthogonal splitting of $\mathcal H$ into $\mathcal L$-invariant coordinate subspaces, with $\Delta_{\min}$ the gap between clusters, but inside a cluster the small double-commutator weights remain and $\eta_D$ loses the closed form \eqref{eq:diag-etaD}.
(iv) Flat modes. $\lambda_P$ enters through $\|\widehat D^{1/2}\mathcal V_{DO}\|^2\le\|\mathcal V_{DO}\|^2/(2(1-f)\lambda_P)$, $\|\mathcal G_{DO}\|\le\|\mathcal V_{DO}\|/(2\lambda_P(1-\beta))$ and the condition in (iii); a mode with $\lambda_P\to0$ and curvature noise coupling it to the rest (the collapse of Theorem~\ref{thm:G5}) makes these constants grow like $\lambda_P^{-1}$ to $\lambda_P^{-2}$, so the diagonal reduction needs $\beta$ closer to $1$. A tether $\mu>0$ replaces $\lambda_P$ by $\lambda_P+\mu$. The $\Delta_{\min}$, $\lambda_P$ constants are crude because they charge every pair, including the top couplings $(1,k)$, with the smallest gap and the flattest mode: in the checks $c_\beta$ exceeds the operator form of Theorem~\ref{thm:diag-threshold}(3) by factors $10$ to $5\times10^4$ and is defined in $23$ of $42$ cells. The commutator-energy constant $c'_\beta$ charges each pair with its own gap; it lies between $0.16$ and $11$ times the operator form and is defined in $35$ of $42$ cells, and the weighted form between $0.16$ and $1.1$ times, defined in all $42$. In the near-flat cell $c'_\beta\approx262\,c_\infty$ at $\beta=0.9999$: $\kappa_\Delta^2$ takes the worst direction of $A$, the flat diagonal mode ($\Lambda_D^{-1/2}\approx7$ there), while $c_\infty$ weighs that mode by its small component in $X_\infty$. The eigenvector split of Theorem~\ref{thm:diag-threshold}(1) removes this: its first term is $\langle u,\mathsf Eu\rangle\to c_\infty$, and the flat direction enters only the second-order term, divided by $\gamma_D=\Theta((1-\beta)^{-1})$; in the near-flat cell it is $5.1$, $1.6$, $1.2$, $1.04$, $1.004$, $1.0005$ times the true value at $\beta=0.8,\dots,0.9999$.
\end{remark}

\paragraph{Numerical check (\texttt{paper/checks/diagonalization-burnin.py}).} $\mathcal T$ is built in Kraus form on $\mathbb S^{2P}$ by enumerating all $N^B$ batches of non-commuting positive semidefinite per-sample Hessians ($P\in\{3,4,5\}$, $N\in\{4,5,6,8\}$, $B\in\{1,2\}$; Gram matrices of random Gaussian factors, centered perturbations of a $H$ with a double and a triple eigenvalue, a near-flat mode, and Rademacher features), and $\rho(\mathcal T)$ is computed as $\lim\|\mathcal T^{2^K}\|^{1/2^K}$ by $46$ normalized squarings. In all $14$ cells ($7$ models, $\beta\in\{0.5,0.9\}$), $\rho(\mathcal K)=1$ at $\eta^\star=1/\lambda_{\max}(\mathcal G_\beta)$ to $10^{-12}$, $|\rho(\mathcal T(\eta^\star))-1|\le1.1\times10^{-13}$, and $\rho(\mathcal T)$ crosses $1$ between $0.999\,\eta^\star$ and $1.001\,\eta^\star$. Both inequalities of Theorem~\ref{thm:diag-K}(1) and the bound \eqref{eq:diag-rel} hold in all $30$ cells at $f=\frac12$, $\beta\in\{0.5,0.8,0.9,0.95,0.99\}$; the log-slope of the relative error in $1-\beta$ between $\beta=0.95$ and $0.99$ is $2.015$--$2.024$ in the six models, including the degenerate ones. At the threshold, $(\eta_D/\eta^\star-1)$ has log-slope $0.98$--$1.00$ between $\beta=0.999$ and $0.9999$ and $1/\eta^\star-1/\eta_D$ agrees with $c_\infty$ of Proposition~\ref{prop:diag-expansion} to $4$ digits at $\beta=0.9999$ (e.g.\ $7.051$ against $7.050$, $0.9065$ against $0.9065$). For Rademacher features $\mathcal V_{DO}=0$ and $\eta^\star$, $\eta_D$ and the root of \eqref{eq:secular} agree to $12$ significant digits. Scalar constants against the operator forms (ratio range; cells where the constant is defined): at the threshold, the eigenvector split of Theorem~\ref{thm:diag-threshold}(1) is $0.0038$--$10.3$ times the operator form (3) with $\mathsf E_\Delta$ ($35/42$) and $0.0038$--$1.75$ times with $\bar{\mathsf E}$ ($42/42$), and $1.000$--$17.7$ (resp.\ $1.000$--$17.5$) times the true $1/\eta^\star-1/\eta_D$, the ratio tending to $1$ as $\beta\to1$ in every model ($\le1.001$ at $\beta=0.9999$); condition \eqref{eq:diag-eigcond} held in every cell where the majorant was defined, so its only failures are the $7$ small-$\beta$ cells with $\theta_\beta\ge1$, where $\bar{\mathsf E}$ is used. The other forms: $c_\beta/\text{op}\in[10,5.1\times10^4]$ ($23/42$), $c'_\beta/\text{op}\in[0.16,11]$ ($35/42$), weighted form $\in[0.16,1.1]$ ($42/42$); against the true $1/\eta^\star-1/\eta_D$ the ranges are $[41,1.3\times10^7]$, $[2.0,588]$ and $[1.8,262]$. At fixed $f$, the $\Delta_{\min}$ form of Theorem~\ref{thm:diag-K}(3) is $21$--$1.3\times10^3$ times \eqref{eq:diag-rel} ($21/30$) and the form (3$'$) is $0.016$--$3.9$ times it ($25/30$); against the true relative error, $53$--$5.3\times10^4$ and $1.5$--$40$. Every scalar bound holds in every cell where it is defined; $c_\infty=\|D_\Delta^{-1}\mathcal V_{OD}\Lambda_D^{-1/2}X_\infty\|^2$ matches the sum form to $10^{-10}$ and $\kappa_\Delta^2/c_\infty\in\{2.0,2.2,2.1,3.4,3.8,262\}$ in the six models.

% ---------------------------------------------------------------------------------------------
\subsection{Burn-in in Loewner order and the spectral gap of the second-moment operator}\label{sec:burnin}

Let $z_t=(w_t,w_{t-1})$, $Q_t=\E[z_tz_t^\top]$, $\mathcal T(Q)=\E[F_tQF_t^\top]$ on $\mathbb S^{2P}$ as in Lemma~\ref{lem:russodye}, and $G=\operatorname{blkdiag}(\eta^2\Sigma_\xi,0)$. Assume
\textbf{(H0)} $\E[z_0]=0$, or $\E[(H_{B_t}-H)a\,\xi_t^\top]=0$ for every $a\in\R^P$.
Then $Q_{t+1}=\mathcal T(Q_t)+G$ for every $t$: conditioning on $(F_t,\xi_t)$, independent of $z_t$, the quadratic term is $\mathcal T(Q_t)$, and the cross term $-\eta\E[F_tz_t\xi_t^\top]=-\eta\E[F_t\,\E[z_t]\,\xi_t^\top]$ vanishes under (H0) because $\E[z_t]=\bar F^t\E[z_0]$ with $\bar F=\E[F_t]$ and $\E[F_ta\xi_t^\top]=-\eta\operatorname{blkdiag}(\E[(H_{B_t}-H)a\xi_t^\top],0)$ for the first block row.

\begin{theorem}[Burn-in in Loewner order]\label{thm:burnin}
Assume (H0) and $\rho(\mathcal T)<1$, and let $Q_\infty=(I-\mathcal T)^{-1}(G)=\sum_{k\ge0}\mathcal T^k(G)\succeq0$.
(1) $Q_\infty-Q_t=\mathcal T^t(Q_\infty-Q_0)$ for every $t\ge0$.
(2) For every $\theta\in(\rho(\mathcal T),1)$, $\theta>0$, the series $Q_\theta\triangleq\sum_{k\ge0}\theta^{-k}\mathcal T^k(I)$ converges, $Q_\theta\succeq I$ and $\mathcal T(Q_\theta)=\theta(Q_\theta-I)\preceq\theta Q_\theta$. With $c_\pm\triangleq\min\{s\ge0:\pm(Q_\infty-Q_0)\preceq sQ_\theta\}\le\|(Q_\infty-Q_0)_\pm\|$,
\[
-c_-\theta^tQ_\theta\ \preceq\ Q_\infty-Q_t\ \preceq\ c_+\theta^tQ_\theta\qquad(t\ge0).
\]
If $Q_0\preceq Q_\infty$ (for instance $Q_0=0$, the start at the minimizer), then $c_-=0$ and $0\preceq Q_\infty-Q_t\preceq c_+\theta^tQ_\theta$; in particular $0\le\E\|w_\infty\|^2-\E\|w_t\|^2\le c_+\theta^t\tr\big((Q_\theta)_{11}\big)$.
(3) If $\beta\in(0,1)$ and the transitions are irreducible, i.e.\ no subspace $\{0\}\ne\mathcal W\ne\R^{2P}$ satisfies $F(H_{B})\mathcal W\subseteq\mathcal W$ for almost every batch, then $\mathcal T$ has an eigenmatrix $Q_*\succ0$ with $\mathcal T(Q_*)=\rho(\mathcal T)Q_*$, and (2) holds with $\theta=\rho(\mathcal T)$ and $Q_\theta$ replaced by $Q_*$:
\[
-c_-^*\rho(\mathcal T)^tQ_*\preceq Q_\infty-Q_t\preceq c_+^*\rho(\mathcal T)^tQ_*,\qquad c_\pm^*=\lambda_{\max}\big(\pm Q_*^{-1/2}(Q_\infty-Q_0)Q_*^{-1/2}\big)_+ .
\]
\end{theorem}
\begin{proof}
(1) Subtract $Q_{t+1}=\mathcal T(Q_t)+G$ from $Q_\infty=\mathcal T(Q_\infty)+G$ and iterate. (2) $\rho(\mathcal T/\theta)<1$ gives convergence; $Q_\theta\succeq I$ because every term is positive semidefinite (positivity of $\mathcal T^k$) and the $k=0$ term is $I$; $\mathcal T(Q_\theta)=\sum_{k\ge0}\theta^{-k}\mathcal T^{k+1}(I)=\theta\sum_{k\ge1}\theta^{-k}\mathcal T^k(I)=\theta(Q_\theta-I)$. Since $\mathcal T$ is linear and positive, $X\preceq Y$ implies $\mathcal T(X)\preceq\mathcal T(Y)$, so $\mathcal T^t(Q_\theta)\preceq\theta^tQ_\theta$ by induction. Apply $\mathcal T^t$ to $-c_-Q_\theta\preceq Q_\infty-Q_0\preceq c_+Q_\theta$ and use (1). The bound on $c_\pm$ follows from $Q_\theta\succeq I$; the trace statement takes the trace of the $(1,1)$ block. (3) $\mathcal T$ maps the proper cone $\mathbb S^{2P}_+$ into itself, so by the finite-dimensional Krein--Rutman (Perron--Frobenius) theorem \citep[Thm.~1.3.2]{berman1994} $\rho(\mathcal T)$ is an eigenvalue with an eigenvector $Q_*\in\mathbb S^{2P}_+\setminus\{0\}$. $\rho(\mathcal T)>0$: by Jensen, $\mathcal T(I)=\E[F_tF_t^\top]\succeq\bar F\bar F^\top\succ0$ ($\det\bar F=\beta^P\ne0$), so $\mathcal T(I)\succeq\sigma I$ for some $\sigma>0$ and $\mathcal T^n(I)\succeq\sigma^nI$. If $x\perp\operatorname{range}(Q_*)$, then $0=\rho(\mathcal T)\,x^\top Q_*x=x^\top\mathcal T(Q_*)x=\E\|Q_*^{1/2}F_t^\top x\|^2$, so $F_t^\top x\in\ker Q_*=\operatorname{range}(Q_*)^\perp$ almost surely: $F_t^\top$ preserves $\operatorname{range}(Q_*)^\perp$, hence $F_t$ preserves $\operatorname{range}(Q_*)$, almost surely. Irreducibility and $Q_*\ne0$ give $\operatorname{range}(Q_*)=\R^{2P}$. The rest is (2) with $\mathcal T^t(Q_*)=\rho(\mathcal T)^tQ_*$.
\end{proof}
Irreducibility can fail: if every $H_{B}$ commutes with $H$, each $\operatorname{span}\{(u_k,0),(0,u_k)\}$ is invariant and $Q_*$ may be singular; (2) needs no such condition. The matrix $Q_\varepsilon$ of the proof of Theorem~\ref{thm:G}, which satisfies $\mathcal T(Q_\varepsilon)\preceq\theta Q_\varepsilon$ for some $\theta<1$, can be used in (2) in place of $Q_\theta$.

\paragraph{The regular splitting and the spectral gap.} Write $\mathcal T=\mathcal T_0+\Delta$ with $\mathcal T_0(Q)=\bar FQ\bar F^\top$ (the noise-free part) and $\Delta(Q)=\E[(F_t-\bar F)Q(F_t-\bar F)^\top]=\operatorname{blkdiag}(\eta^2\mathcal V(Q_{11}),0)$; both are positive, and $\rho(\mathcal T_0)=\rho(\bar F)^2<1$ for $\eta\in(0,\eta_{\mathrm{det}})$. Let $\mathcal R=(I-\mathcal T_0)^{-1}\Delta$, $\Pi(Q)=Q_{11}$ and $\mathcal J(X)=(I-\mathcal T_0)^{-1}\big(\operatorname{blkdiag}(\eta^2\mathcal V(X),0)\big)$ for $X\in\mathbb S^P$.

\begin{theorem}[Regular splitting: $\rho(\mathcal R)=\rho(\mathcal K)$ and a two-sided spectral gap]\label{thm:splitting}
Let $\beta\in[0,1)$ and $\eta\in(0,\eta_{\mathrm{det}})$, with $\beta>0$ in parts (3)--(6).
(1) $\mathcal R=\mathcal J\circ\Pi$ and $\Pi\circ\mathcal J=\mathcal K_{\eta,\beta}$. Hence $\mathcal R$ and $\mathcal K_{\eta,\beta}$ have the same non-zero eigenvalues with the same algebraic multiplicities, and $r\triangleq\rho(\mathcal R)=\rho(\mathcal K_{\eta,\beta})$ exactly.
(2) Let $X_*\in\mathbb S^P_+\setminus\{0\}$ be the Krein--Rutman eigenmatrix of $\mathcal K_{\eta,\beta}$ (Theorem~\ref{thm:G6}(1)), $Y_*=\operatorname{blkdiag}(\eta^2\mathcal V(X_*),0)\succeq0$ and $W=\mathcal J(X_*)=\sum_{k\ge0}\mathcal T_0^k(Y_*)\succeq0$. Then $\mathcal R(W)=rW$, $\Pi(W)=rX_*$, $(I-\mathcal T_0)(W)=Y_*$, and
\begin{equation}\label{eq:split-identity}
\mathcal T(W)=rW+(1-r)\mathcal T_0(W)=W-(1-r)Y_* .
\end{equation}
(3) $W\succ0$ iff no eigenvector of $H$ is orthogonal to $\operatorname{range}\mathcal V(X_*)$; in particular if $\mathcal V(X_*)\succ0$.
(4) Lower bound. With $s_{\mathrm{low}}=\max\{s\ge0:\mathcal T_0(W)\succeq sW\}$, $\rho(\mathcal T)\ge r+(1-r)s_{\mathrm{low}}$ if $r\le1$, and $\rho(\mathcal T)\ge1$ if $r\ge1$. If $W\succ0$, $s_{\mathrm{low}}=\lambda_{\min}(W^{-1/2}\mathcal T_0(W)W^{-1/2})>0$.
(5) The one-step upper bound is empty. If $r\le1$, $W\succ0$ and $\mathcal T_0(W)\preceq sW$, then $\rho(\mathcal T)\le r+(1-r)s$; but for this $W$ the least such $s$ is exactly $1$, so the bound reads $\rho(\mathcal T)\le1$.
(6) Two-step upper bound. If $r\le1$ and $\mathcal V(X_*)\succ0$, then $Z\triangleq Y_*+\mathcal T(Y_*)\succ0$, $W\succ0$, $\mathcal T^2(W)=W-(1-r)Z$, and with $m_2\triangleq\lambda_{\min}(W^{-1/2}ZW^{-1/2})>0$, $\rho(\mathcal T)^2\le1-(1-r)m_2$. Therefore, for $r<1$,
\[
\frac{(1-r)\,m_2}{2}\ \le\ 1-\rho(\mathcal T)\ \le\ (1-r)(1-s_{\mathrm{low}}),\qquad r=\rho(\mathcal K_{\eta,\beta}).
\]
More generally, if $Z_m\triangleq\sum_{k<m}\mathcal T^k(Y_*)\succ0$ and $W\succ0$, then $\rho(\mathcal T)^m\le1-(1-r)\lambda_{\min}(W^{-1/2}Z_mW^{-1/2})$.
\end{theorem}
\begin{proof}
(1) $\Delta=\operatorname{blkdiag}(\eta^2\mathcal V(\Pi\,\cdot),0)$ gives $\mathcal R=\mathcal J\circ\Pi$. For $Y\in\mathbb S^P_+$, $Q=(I-\mathcal T_0)^{-1}\operatorname{blkdiag}(\eta^2Y,0)$ is the stationary second moment of the noise-free recursion \eqref{eq:hbP} with $\mathcal V=0$ and $\Sigma_\xi=Y$ (driven by centered Gaussian noise, so $\E[w_t]=0$), and Theorem~\ref{thm:F}(2) gives $\mathcal L_{\eta,\beta}(Q_{11})=\frac{\eta}{1-\beta}Y$, i.e.\ $\Pi(I-\mathcal T_0)^{-1}\operatorname{blkdiag}(\eta^2Y,0)=\frac{\eta}{1-\beta}\mathcal L_{\eta,\beta}^{-1}(Y)$ ($\mathcal L_{\eta,\beta}$ is invertible by Theorem~\ref{thm:G6}(1)); both sides are linear in $Y$ and $\mathbb S^P_+$ spans $\mathbb S^P$, so it holds on $\mathbb S^P$. With $Y=\mathcal V(X)$, $\Pi\mathcal J=\frac{\eta}{1-\beta}\mathcal L^{-1}\mathcal V=\mathcal K$. For linear maps $\mathcal J:\mathbb S^P\to\mathbb S^{2P}$ and $\Pi:\mathbb S^{2P}\to\mathbb S^P$, $\mathcal J\Pi$ and $\Pi\mathcal J$ have the same non-zero eigenvalues with multiplicities ($\det(zI-\mathcal J\Pi)=z^{\dim\mathbb S^{2P}-\dim\mathbb S^P}\det(zI-\Pi\mathcal J)$). The phrase ``compression to the $(1,1)$ block'' in the proof of Theorem~\ref{thm:G6} is made precise by this factorization; a compression alone would not fix the spectral radius.
(2) $\mathcal R(W)=\mathcal J(\Pi\mathcal J X_*)=\mathcal J(\mathcal KX_*)=r\mathcal J(X_*)=rW$ and $\Pi W=\mathcal KX_*=rX_*$. $(I-\mathcal T_0)W=Y_*$ by definition of $\mathcal J$, and $\Delta(W)=\operatorname{blkdiag}(\eta^2\mathcal V(rX_*),0)=rY_*$. Hence $\mathcal T(W)=\mathcal T_0(W)+rY_*=W-Y_*+rY_*=W-(1-r)Y_*$, and $rW+(1-r)\mathcal T_0(W)=rW+(1-r)(W-Y_*)$ is the same matrix. $W\succeq0$ because $Y_*\succeq0$ and $\mathcal T_0$ is positive.
(3) $x\in\ker W$ iff $Y_*^{1/2}(\bar F^\top)^kx=0$ for all $k$, iff $x$ is orthogonal to the reachable subspace of the pair $(\bar F,\operatorname{range}Y_*)$. By the Popov--Belevitch--Hautus test, that subspace is $\R^{2P}$ iff no left eigenvector $z=(p,q)$ of $\bar F$ is orthogonal to $\operatorname{range}Y_*=\operatorname{range}\mathcal V(X_*)\oplus0$. From $z^\top\bar F=\nu z^\top$ with $\bar F=\begin{pmatrix}A&-\beta I\\I&0\end{pmatrix}$, $A=(1+\beta)I-\eta H$: $\nu\ne0$ ($\det\bar F=\beta^P$), $q=-\beta p/\nu$ and $Ap=(\nu+\beta/\nu)p$, so $p$ lies in an eigenspace of $H$; conversely every real eigenvector $p$ of $H$, $Hp=\lambda p$, gives left eigenvectors with $\nu^2-((1+\beta)-\eta\lambda)\nu+\beta=0$. Orthogonality of $z$ to $\operatorname{range}Y_*$ is $p\perp\operatorname{range}\mathcal V(X_*)$ (for complex $p$ in an eigenspace, its real or imaginary part is a non-zero real eigenvector with the same property).
(4) By \eqref{eq:split-identity}, $\mathcal T(W)=rW+(1-r)\mathcal T_0(W)\succeq(r+(1-r)s_{\mathrm{low}})W$ when $r\le1$, and $\mathcal T(W)=W+(r-1)Y_*\succeq W$ when $r\ge1$. If $\mathcal T(W)\succeq cW$ with $c\ge0$ and $W\in\mathbb S^{2P}_+\setminus\{0\}$, then $\mathcal T^n(W)\succeq c^nW$, so $\|\mathcal T^n\|\ge c^n\|W\|/\|W\|$ and $\rho(\mathcal T)\ge c$ by Gelfand's formula. $\mathcal T_0(W)=\bar FW\bar F^\top\succ0$ when $W\succ0$.
(5) If $W\succ0$ and $\mathcal T(W)\preceq\theta W$ then every $X\in\mathbb S^{2P}_+$ satisfies $X\preceq\|X\|\|W^{-1}\|W$, so $\mathcal T^n(X)\preceq\theta^n\|X\|\|W^{-1}\|W$ and $\rho(\mathcal T)\le\theta$ by Lemma~\ref{lem:russodye}; with $\theta=r+(1-r)s$ this is the stated bound. Now $W-\mathcal T_0(W)=Y_*\succeq0$ gives $s\le1$, and for $x=(0,y)$, $y\ne0$, $x^\top\mathcal T_0(W)x=x^\top Wx-x^\top Y_*x=x^\top Wx>0$ because the $(2,2)$ block of $Y_*$ is zero, so $s\ge1$. This is why a one-step contraction of $W$ under $\mathcal T_0$ cannot be used here.
(6) Applying $\mathcal T$ to \eqref{eq:split-identity}: $\mathcal T^2(W)=\mathcal T(W)-(1-r)\mathcal T(Y_*)=W-(1-r)(Y_*+\mathcal T(Y_*))$. With $a=\eta^2\mathcal V(X_*)\succ0$, the quadratic form of $Z$ at $(x,y)$ is $x^\top ax+\E\big[(A_tx+y)^\top a(A_tx+y)\big]$, $A_t=(1+\beta)I-\eta H_{B_t}$ (because $F_t^\top(x,y)=(A_tx+y,-\beta x)$), which vanishes only if $x=0$ and then $y=0$; so $Z\succ0$. $W\succ0$ by (3). Then $\mathcal T^2(W)\preceq(1-(1-r)m_2)W$ and, as in (5) applied to $\mathcal T^2$, $\rho(\mathcal T)^2\le1-(1-r)m_2$. For $r<1$, $\rho(\mathcal T)\le1$ and $1-\rho(\mathcal T)\ge\frac{1-\rho(\mathcal T)^2}{2}\ge\frac{(1-r)m_2}2$; the upper bound on $1-\rho(\mathcal T)$ is (4). The $m$-step statement is the same computation with $\mathcal T^m(W)=W-(1-r)Z_m$.
\end{proof}

So the spectral gap of the full second-moment operator is the gap $1-\rho(\mathcal K_{\eta,\beta})$ of the reduced operator times a factor between $m_2/2$ and $1-s_{\mathrm{low}}$, both read off the single matrix $W=\mathcal J(X_*)$; with Theorem~\ref{thm:burnin}, $t\ge\log(c_+/\delta)/\log(1/\rho(\mathcal T))$ steps bring $Q_\infty-Q_t$ below $\delta Q_*$, and $\log(1/\rho(\mathcal T))\ge1-\rho(\mathcal T)\ge(1-\rho(\mathcal K))m_2/2$. The lower bound is loose in the checks: $m_2$ is small because $W$ is ill-conditioned.

\paragraph{Numerical check (\texttt{paper/checks/diagonalization-burnin.py}).} In five cells ($P\in\{3,4,5\}$, a double eigenvalue, $\beta\in\{0.5,0.8,0.9,0.95\}$, $\eta\in\{0.7,0.8,0.9\}\eta^\star$): $|\rho(\mathcal R)-\rho(\mathcal K)|\le2\times10^{-14}$; $\|\mathcal R W-rW\|/\|W\|\le4\times10^{-14}$ and the residual of \eqref{eq:split-identity} is $\le6\times10^{-16}$; $X_*\succ0$, $W\succ0$, $Q_*\succ0$ in every cell; the one-step constant is $s=1.000000000000$ in every cell, as Theorem~\ref{thm:splitting}(5) predicts; the two-sided gap bound holds, e.g.\ $1.6\times10^{-4}\le1-\rho(\mathcal T)=1.1\times10^{-2}\le2.5\times10^{-2}$ at $P=3$, $\beta=0.9$, $\eta=0.9\eta^\star$. The Loewner envelopes of Theorem~\ref{thm:burnin} hold for $t\le300$ from $Q_0=0$ and from a $Q_0$ with $Q_\infty-Q_0$ indefinite, with both $Q_*$ and $Q_\theta$, $\theta=(1+\rho(\mathcal T))/2$ (smallest normalized eigenvalue of the envelope gap $\ge-4\times10^{-14}$, round-off).

\end{document}
